\documentclass[11pt,a4paper,oneside,openany]{book}

\usepackage[T1]{fontenc}
\usepackage[utf8]{inputenc}
\usepackage{libertinus}
\usepackage{microtype}
\usepackage[a4paper,margin=27mm,headheight=15pt]{geometry}
\usepackage{amsmath,amssymb,amsthm,mathtools}
\usepackage{bm}
\usepackage{mathrsfs}
\usepackage{booktabs,longtable,array,tabularx}
\usepackage{enumitem}
\usepackage{xcolor}
\usepackage{tikz}
\usetikzlibrary{arrows.meta,positioning,calc}
\usepackage[most]{tcolorbox}
\usepackage{fancyhdr}
\usepackage{xurl}
\usepackage{hyperref}
\usepackage[nameinlink,noabbrev,capitalise]{cleveref}
% Canonical mathematical notation for active FabiusFunction LaTeX documents.
%
% This file is intentionally a plain TeX input rather than a LaTeX package.
% Every standalone document loads it by an explicit path relative to that
% document's directory; included fragments inherit the definitions from their
% root document.  Keep this file independent of document layout and metadata.
%
% Contract:
%   * one command has one mathematical meaning throughout the active corpus;
%   * names encode semantic roles rather than a document-local choice of letter;
%   * visually similar but differently normalized objects have different print;
%   * every command below is defined normatively in the notation catalogue.
%
% The active corpus excludes docs/archive/.  Historical sources may retain
% their original notation only inside an explicitly labelled quotation.

\makeatletter
\@ifundefined{mathbb}{%
  \PackageError{fabius-notation}{The amssymb package must be loaded first}%
    {Load amsmath and amssymb before inputting fabius-notation.tex.}%
}{}
\@ifundefined{operatorname}{%
  \PackageError{fabius-notation}{The amsmath package must be loaded first}%
    {Load amsmath before inputting fabius-notation.tex.}%
}{}
\makeatother

% Version stamp used by the catalogue and by static audits.
\newcommand{\FabiusNotationVersion}{2026-08-31}

% Number systems.  NaturalNumbers is explicitly the nonnegative convention.
\newcommand{\NaturalNumbers}{\mathbb{N}_{0}}
\newcommand{\PositiveIntegers}{\mathbb{N}_{>0}}
\newcommand{\IntegerNumbers}{\mathbb{Z}}
\newcommand{\RationalNumbers}{\mathbb{Q}}
\newcommand{\RealNumbers}{\mathbb{R}}
\newcommand{\ComplexNumbers}{\mathbb{C}}
\newcommand{\ExtendedRealNumbers}{\overline{\mathbb{R}}}

% The three principal Fabius--Rvachev functions.  Subscripts are deliberate:
% a bare F and a calligraphic F have several unrelated uses in the corpus.
\newcommand{\FabiusBounded}{F_{\mathrm{bd}}}
\newcommand{\FabiusGlobal}{F_{\mathrm{gl}}}
\newcommand{\FabiusQuantile}{Q_{F}}
\newcommand{\FabiusClampedQuantile}{Q_{F,\mathrm{cl}}}
\newcommand{\RvachevUp}{\operatorname{up}}

% Constants, elementary operators, and delimiters.
\newcommand{\EulerE}{\mathrm{e}}
\newcommand{\ImaginaryUnit}{\mathrm{i}}
\newcommand{\Differential}{\mathop{}\!\mathrm{d}}
\newcommand{\AbsoluteValue}[1]{\left\lvert #1\right\rvert}
\newcommand{\Norm}[1]{\left\lVert #1\right\rVert}
\newcommand{\Floor}[1]{\left\lfloor #1\right\rfloor}
\newcommand{\Ceiling}[1]{\left\lceil #1\right\rceil}
\newcommand{\FractionalPart}[1]{\left\{#1\right\}}
\newcommand{\PositivePart}[1]{\left(#1\right)_{+}}
\newcommand{\IndicatorOf}[1]{\mathbf{1}_{#1}}
\newcommand{\SetOf}[1]{\left\{#1\right\}}
\newcommand{\InnerProduct}[1]{\left\langle #1\right\rangle}
\newcommand{\CardinalityOf}[1]{\left\lvert #1\right\rvert}
\newcommand{\DefinitionEquals}{\mathrel{\coloneqq}}
\newcommand{\Sign}{\operatorname{sgn}}
\newcommand{\IdentityOperator}{\operatorname{id}}
\newcommand{\SpanOperator}{\operatorname{span}}
\newcommand{\TraceOperator}{\operatorname{tr}}
\newcommand{\RankOperator}{\operatorname{rank}}
\newcommand{\DiagonalOperator}{\operatorname{diag}}
\newcommand{\ResidueOperator}{\operatorname*{Res}}
\newcommand{\LeastCommonMultiple}{\operatorname{lcm}}
\newcommand{\OrderOperator}{\operatorname{ord}}
\newcommand{\PrincipalLogarithm}{\operatorname{Log}}
\newcommand{\FallingFactorial}[2]{\left(#1\right)^{\underline{#2}}}
\newcommand{\RisingFactorial}[2]{\left(#1\right)^{\overline{#2}}}

% Support, topology, convolution, and metric notation.
\newcommand{\SupportOperator}{\operatorname{supp}}
\newcommand{\SupportOf}[1]{\SupportOperator\!\left(#1\right)}
\newcommand{\TopologicalSupportOf}[1]{\operatorname{tsupp}\!\left(#1\right)}
\newcommand{\Convolution}{\mathbin{*}}
\newcommand{\MetricDistance}{\operatorname{dist}}
\newcommand{\TotalVariationDistance}{d_{\mathrm{TV}}}

% Probability.  Equality and convergence in law are not metric distance.
\newcommand{\Probability}{\mathbb{P}}
\newcommand{\Expectation}{\mathbb{E}}
\newcommand{\Variance}{\operatorname{Var}}
\newcommand{\Covariance}{\operatorname{Cov}}
\newcommand{\LawOperator}{\operatorname{Law}}
\newcommand{\LawOf}[1]{\LawOperator\!\left(#1\right)}
\newcommand{\EqualInLaw}{\mathrel{\overset{\mathrm{law}}{=}}}
\newcommand{\ConvergesInLaw}{\mathrel{\xrightarrow{\mathrm{law}}}}

% Fourier and integral transforms.  The printed labels expose normalization.
% FourierTwoPi uses exp(-2*pi*i*x*xi); FourierAngular uses exp(-i*x*omega).
\newcommand{\FourierTwoPi}[1]{\widehat{#1}_{\,2\pi}}
\newcommand{\FourierAngular}[1]{\widehat{#1}_{\,\mathrm{ang}}}
\newcommand{\LaplaceTransformSymbol}{\mathcal{L}}
\newcommand{\MellinTransformSymbol}{\mathcal{M}}
\newcommand{\LaplaceTransformOf}[1]{\LaplaceTransformSymbol\!\left[#1\right]}
\newcommand{\MellinTransformOf}[1]{\MellinTransformSymbol\!\left[#1\right]}
\newcommand{\MomentGeneratingFunction}[1]{M_{#1}}
\newcommand{\CharacteristicFunction}[1]{\varphi_{#1}}

% The two standard sinc functions must never share one symbol.
\newcommand{\SincRad}{\operatorname{sinc}_{\mathrm{rad}}}
\newcommand{\SincPi}{\operatorname{sinc}_{\pi}}

% Binary arithmetic and Thue--Morse notation.
\newcommand{\BinaryDigitSum}{\operatorname{wt}_{2}}
\newcommand{\ThueMorseSignSymbol}{\varepsilon}
\newcommand{\ThueMorseSign}[1]{\varepsilon_{#1}}
\newcommand{\PadicValuation}[1]{\nu_{#1}}
\newcommand{\TwoAdicValuation}{\nu_{2}}

% q-series notation.  Every base is an explicit argument.
\newcommand{\QPochhammer}[3]{\left(#1;#2\right)_{#3}}
\newcommand{\GaussianBinomial}[3]{%
  \genfrac{[}{]}{0pt}{}{#1}{#2}_{#3}%
}
\newcommand{\QInteger}[2]{\left[#1\right]_{#2}}

% Named special functions.
\newcommand{\Polylogarithm}{\operatorname{Li}}
\newcommand{\LambertW}{W}
\newcommand{\GammaFunction}{\Gamma}
\newcommand{\RiemannZeta}{\zeta}

% Asymptotics.  BigO and LittleO are operator tokens for existing formula
% grammar; prefer the At forms so that the limiting regime is printed.
\newcommand{\BigO}{\mathcal{O}}
\newcommand{\LittleO}{o}
\newcommand{\BigOAt}[2]{\mathcal{O}_{#2}\!\left(#1\right)}
\newcommand{\LittleOAt}[2]{o_{#2}\!\left(#1\right)}

\usepackage{makeidx}
\makeindex

\definecolor{deepolive}{RGB}{67,79,45}
\definecolor{softolive}{RGB}{236,239,226}
\definecolor{darkteal}{RGB}{31,83,84}
\definecolor{warmgray}{RGB}{92,87,78}
\definecolor{softcoral}{RGB}{246,225,219}
\definecolor{linkblue}{RGB}{31,76,122}

\hypersetup{
  colorlinks=true,
  linkcolor=deepolive,
  citecolor=darkteal,
  urlcolor=linkblue,
  pdftitle={Combinatorial Coefficient Calculus: A Unified Formulae Monograph},
  pdfauthor={A self-contained synthesis with proofs},
  pdfsubject={Stirling, Bell, Eulerian, Bernoulli-Euler, Faa di Bruno, Lagrange-Good inversion, and asymptotic coefficient calculus}
}

\setlist{itemsep=2pt,topsep=4pt,parsep=1pt}
\allowdisplaybreaks[3]
\raggedbottom
\emergencystretch=1em
\setcounter{tocdepth}{2}
\setcounter{secnumdepth}{3}

\pagestyle{fancy}
\fancyhf{}
\fancyhead[L]{\small\nouppercase{\leftmark}}
\fancyhead[R]{\small\thepage}
\renewcommand{\headrulewidth}{0.4pt}
\fancypagestyle{plain}{%
  \fancyhf{}%
  \fancyfoot[C]{\thepage}%
  \renewcommand{\headrulewidth}{0pt}%
}

\newtheorem{theorem}{Theorem}[chapter]
\newtheorem{lemma}[theorem]{Lemma}
\newtheorem{proposition}[theorem]{Proposition}
\newtheorem{corollary}[theorem]{Corollary}
\newtheorem{identity}[theorem]{Identity}
\theoremstyle{definition}
\newtheorem{definition}[theorem]{Definition}
\newtheorem{example}[theorem]{Example}
\newtheorem{algorithm}[theorem]{Algorithm}
\theoremstyle{remark}
\newtheorem{remark}[theorem]{Remark}
\newtheorem{historical}[theorem]{Historical note}

\newtcolorbox{masterbox}[1][]{
  enhanced,breakable,colback=softolive,colframe=deepolive,
  boxrule=0.7pt,arc=1mm,left=2mm,right=2mm,top=1.5mm,bottom=1.5mm,
  title={Master principle},fonttitle=\bfseries,#1
}
\newtcolorbox{auditbox}[1][]{
  enhanced,breakable,colback=softcoral,colframe=warmgray,
  boxrule=0.6pt,arc=1mm,left=2mm,right=2mm,top=1.5mm,bottom=1.5mm,
  title={Editorial normalization},fonttitle=\bfseries,#1
}


\title{\Huge\bfseries Combinatorial Coefficient Calculus\\[3mm]
\Large Stirling, Bell, Eulerian, and Bernoulli--Euler Families;\\
Bell-Polynomial and Fa\`a di Bruno Calculus;\\
Lagrange--B\"urmann and Multivariate Good Inversion}
\author{A self-contained synthesis with proofs}
\date{September 2026}

\begin{document}
\frontmatter
\hypersetup{pageanchor=false}
\maketitle
\hypersetup{pageanchor=true}

\chapter*{Abstract}
\addcontentsline{toc}{chapter}{Abstract}
This monograph develops a unified calculus for the coefficients that occur when one
changes polynomial bases, partitions labelled sets, differentiates or composes
functions, and reverses formal or analytic power series.  Its principal objects are
the signed and unsigned Stirling numbers of the first kind, the Stirling numbers of
the second kind, Lah numbers, Bell numbers, exponential and ordinary Bell
polynomials, Eulerian numbers of types $A$ and $B$, second-order Eulerian numbers,
Fa\`a di Bruno's formula, and the Lagrange inversion and reversion theorems.

The treatment is deliberately proof-oriented.  Definitions are connected by
bijections and change-of-basis arguments; generating functions are derived rather
than merely recorded; recurrences, convolutions, congruences, determinant formulae,
probabilistic interpretations, special-function identities, and asymptotic formulae
are proved from a small collection of reusable principles.  Applications include
finite differences, differential-operator normal ordering, moments and cumulants,
symmetric polynomials, cycle indices, Hermite polynomials, polynomial sequences of
binomial type, negative polylogarithms, Bernoulli-number identities, Poisson moments,
Catalan and Fuss--Catalan series, the Lambert $W$ function, and Laplace-type
asymptotic integrals.

The five source drafts are merged at theorem level rather than concatenated.  Their
additional strands---incidence algebras, M\"obius inversion on the partition lattice,
Sheffer--Riordan calculus, Fr\'echet derivatives, inverse-map recursions,
and Good's multivariate inversion formula---are developed in full.  The
Wikipedia and MathWorld dossiers contribute a further unified treatment of Bernoulli,
Euler, Genocchi, and Cauchy--Gregory families; Pochhammer symbols, digamma and polygamma
functions, Barnes' $G$-function; complementary Bell numbers; Catalan--Narayana and
associahedral refinements; Darboux formulas and singularity transfer; exact coefficient
algorithms; and complete worked inverse-series examples including Kepler's equation.

The sources use the letter $B$ for Bell numbers, Bell polynomials, Bernoulli numbers,
type-$B$ Eulerian numbers, and binary-tree numbers.  Here these are separated
notationally.  A final concordance identifies the contribution of each draft and each
reference dossier, while a proof-audit ledger records normalized, corrected, analytic,
and computational statements separately.

\chapter*{How to read this work}
\addcontentsline{toc}{chapter}{How to read this work}
The chapters are arranged by mathematical dependence, not by the order of the source
files.  The two master ideas are:
\begin{enumerate}
  \item a coefficient is often easiest to compute after changing to the basis in
        which the relevant operator is diagonal; and
  \item an exponential generating function remembers the decomposition of a
        labelled object into components.
\end{enumerate}
Stirling numbers implement the first idea.  Bell polynomials implement the second.
Eulerian numbers mediate between powers and shifted binomial coefficients.  Fa\`a di
Bruno's formula is the differential form of labelled composition, while Lagrange
inversion is the coefficient form of implicit substitution.

A result stated for formal power series requires no convergence.  When an analytic
version is asserted, its domain hypotheses are written explicitly.  Statements in
the archive that report historical computations---for example, lists of Bell
primes or verified exact modular periods---are distinguished from theorems and are
accompanied by a certificate procedure rather than misrepresented as symbolic
identities.

\begin{auditbox}
Three systematic repairs are used throughout.
\begin{enumerate}
\item Bell numbers are denoted $\BellNumber n$, complete exponential Bell polynomials by
      $\ExponentialCompleteBellPolynomial n$, partial Bell polynomials by $\ExponentialPartialBellPolynomial nk$, and Bernoulli numbers by
      $\beta_n$.
\item Type-$B$ Eulerian numbers are indexed by descents in positions
      $0,1,\ldots,n-1$ with $\pi_0=0$.  This removes the one-step shift in one formula
      in the archive.
\item The archived Moser--Wyman display contains unnamed quantities
      $B,P_i,Q_i$.  We give the exact saddle-point scheme that defines every
      correction coefficient recursively and state the leading expansion in closed
      form; no undefined placeholder is retained.
\end{enumerate}
\end{auditbox}

\tableofcontents

\chapter*{Notation}
\addcontentsline{toc}{chapter}{Notation}
For $n\in\NaturalNumbers$, write
\[
 \FallingFactorial{x}{n}=x(x-1)\cdots(x-n+1),\qquad
 \RisingFactorial{x}{n}=x(x+1)\cdots(x+n-1),
\]
with both expressions equal to $1$ when $n=0$.  Coefficient extraction is denoted by
$\CoefficientExtraction z n F(z)$, and $\ResidueOperator_z F(z)\Differential z$ means the coefficient of $z^{-1}$ in a
formal Laurent series.  Our principal arrays are
\[
 \SignedStirlingFirstKind{n}{k},\qquad \UnsignedStirlingFirstKind nk=(-1)^{n-k}\SignedStirlingFirstKind{n}{k},\qquad
 \StirlingSecondKind nk,\qquad \TypeAEulerianNumber nk,
\]
where $\SignedStirlingFirstKind{n}{k}$ is the signed first-kind Stirling number, $\UnsignedStirlingFirstKind nk$ its unsigned
version, $\StirlingSecondKind nk$ the second-kind Stirling number, and $\TypeAEulerianNumber nk$ the
ordinary Eulerian number.  Bell numbers are $\BellNumber n$; exponential partial and
complete Bell polynomials are $\ExponentialPartialBellPolynomial nk$ and $\ExponentialCompleteBellPolynomial n$.

Unless stated otherwise, sums over an impossible index range are zero, and
$0^0=1$ in coefficient formulae.  Bernoulli numbers use
\[
 \frac{t}{\EulerE^t-1}=\sum_{n\ge0}\beta_n\frac{t^n}{n!},
 \qquad \beta_1=-\frac12.
\]

\mainmatter

\part{Coefficient Calculus and Labelled Structures}

\chapter{Formal power series and coefficient extraction}
\label{ch:formal}

\section{Formal series as an algebraic environment}
Let $R$ be a commutative $\RationalNumbers$-algebra.  The ring $R[[z]]$ consists of expressions
$F(z)=\sum_{n\ge0}f_nz^n$ with termwise addition and Cauchy multiplication.  A series
is invertible precisely when $f_0$ is invertible in $R$.  If $G(0)=0$, then the
composition $F(G(z))$ is well-defined because each coefficient receives only
finitely many contributions.

\begin{lemma}[Coefficient rules]\label{lem:coeff-rules}
For $F(z)=\sum f_nz^n$ and $G(z)=\sum g_nz^n$,
\begin{align}
 \CoefficientExtraction z n F(z)G(z)&=\sum_{j=0}^n f_jg_{n-j},\label{eq:cauchy}\\
 \CoefficientExtraction z n F'(z)&=(n+1)f_{n+1},\label{eq:der-coeff}\\
 \CoefficientExtraction z n \frac{F(z)}{1-az}&=\sum_{j=0}^n a^{n-j}f_j.\label{eq:geom-conv}
\end{align}
If $F$ is analytic on $|z|\le r$, then also
\begin{equation}\label{eq:cauchy-coeff}
 f_n=\frac{1}{2\pi \ImaginaryUnit}\int_{|z|=r}\frac{F(z)}{z^{n+1}}\Differential z,
 \qquad |f_n|\le r^{-n}\max_{|z|=r}|F(z)|.
\end{equation}
\end{lemma}
\begin{proof}
The first identity is the definition of Cauchy multiplication.  The second follows
by differentiating term by term, and the third by multiplying with
$\sum_{m\ge0}a^mz^m$.  The analytic formula is Cauchy's coefficient formula, and the
bound follows by estimating the contour integral by its length $2\pi r$.
\end{proof}

\section{Formal residues}
For a Laurent series $A(z)$, set
$\ResidueOperator_z A(z)\Differential z=\CoefficientExtraction z{-1}A(z)$.  Then
\begin{equation}\label{eq:res-coeff}
 \CoefficientExtraction z n F(z)=\ResidueOperator_z\frac{F(z)}{z^{n+1}}\Differential z,
 \qquad \ResidueOperator_z A'(z)\Differential z=0.
\end{equation}
The second identity is simply the fact that the derivative of $z^m$ has no
$z^{-1}$ term.

\begin{theorem}[Formal change of variables for residues]\label{thm:res-subst}
Let $u=u(z)\in zR[[z]]$ have invertible linear coefficient.  For every Laurent series
$A$ for which the composition is defined,
\begin{equation}\label{eq:formal-res-subst}
 \ResidueOperator_z A(u(z))u'(z)\Differential z=\ResidueOperator_u A(u)\Differential u.
\end{equation}
\end{theorem}
\begin{proof}
By linearity it suffices to take $A(u)=u^m$.  If $m\ne-1$, then
$u(z)^mu'(z)=\frac{1}{m+1}(u(z)^{m+1})'$ has residue zero.  For $m=-1$,
$u'(z)/u(z)=z^{-1}+H(z)$ with $H\in R[[z]]$, because
$u(z)=az(1+O(z))$.  Thus both sides have residue $1$.
\end{proof}

\begin{masterbox}
The combination of \cref{eq:res-coeff,eq:formal-res-subst} turns implicit
substitution into coefficient extraction.  This is the entire algebraic engine of
Lagrange inversion; no analytic contour is needed until one asks about convergence.
\end{masterbox}

\section{Ordinary and exponential generating functions}
For a sequence $(a_n)$, its ordinary and exponential generating functions are
\[
 A(z)=\sum_{n\ge0}a_nz^n,
 \qquad \ExponentialGeneratingFunctionOf{A}(z)=\sum_{n\ge0}a_n\frac{z^n}{n!}.
\]
Ordinary multiplication gives unweighted convolution.  Exponential multiplication
gives binomial convolution:
\begin{equation}\label{eq:egf-product}
 \ExponentialGeneratingFunctionOf{A}(z)\ExponentialGeneratingFunctionOf{B}(z)
 =\sum_{n\ge0}\left(\sum_{j=0}^n\binom nj a_jb_{n-j}\right)\frac{z^n}{n!}.
\end{equation}
This binomial factor is the number of ways to choose which labels belong to the
first component.

\section{Finite differences and factorial bases}
Let $E$ be the shift operator $(Ef)(x)=f(x+1)$ and
$\Delta=E-1$.  Then
\begin{equation}\label{eq:finite-diff-falling}
 \Delta\FallingFactorial{x}{k}=k\FallingFactorial{x}{k-1},
 \qquad
 \sum_{i=0}^{n-1}\FallingFactorial{i}{k}=\frac{\FallingFactorial{n}{k+1}}{k+1}.
\end{equation}
\begin{proof}
The first follows from
$\FallingFactorial{x+1}{k}-\FallingFactorial{x}{k}=k\FallingFactorial{x}{k-1}$ after factoring the common
product.  Summing this telescoping identity gives the second formula.
\end{proof}

\begin{theorem}[Newton expansion]\label{thm:newton-expansion}
Every polynomial $p$ of degree at most $d$ has the unique expansion
\begin{equation}\label{eq:newton-expansion}
 p(x)=\sum_{k=0}^d\frac{\Delta^kp(0)}{k!}\FallingFactorial{x}{k}
     =\sum_{k=0}^d\Delta^kp(0)\binom{x}{k}.
\end{equation}
\end{theorem}
\begin{proof}
The factorial polynomials form a basis because their degrees are distinct and their
leading coefficients are one.  If $p=\sum a_k\FallingFactorial{x}{k}$, applying $\Delta^j$
and setting $x=0$ kills terms with $k>j$ by a remaining factor $x$, kills terms
with $k<j$ by excessive differencing, and gives $j!a_j$ for $k=j$.
\end{proof}

\section{Labelled products and the exponential formula}
A labelled combinatorial class is one whose atoms carry distinct labels.  Suppose a
connected component of size $j$ has total weight $x_j$.  A set of $k$ such components
with total size $n$ has weight enumerator $\ExponentialPartialBellPolynomial nk(x_1,x_2,\ldots)$; a set of any
number of components has enumerator $\ExponentialCompleteBellPolynomial n(x_1,\ldots,x_n)$.

\begin{theorem}[The exponential formula]\label{thm:exponential-formula}
Let
\[
 C(z)=\sum_{j\ge1}x_j\frac{z^j}{j!}.
\]
Then
\begin{align}
 \frac{C(z)^k}{k!}
  &=\sum_{n\ge k}\ExponentialPartialBellPolynomial nk(x_1,\ldots,x_{n-k+1})\frac{z^n}{n!},
  \label{eq:partial-bell-egf}\\
 \EulerE^{uC(z)}
  &=\sum_{n\ge k\ge0}\ExponentialPartialBellPolynomial nk(x_1,\ldots)u^k\frac{z^n}{n!},
  \label{eq:bivariate-bell-egf}\\
 \EulerE^{C(z)}
  &=\sum_{n\ge0}\ExponentialCompleteBellPolynomial n(x_1,\ldots,x_n)\frac{z^n}{n!}.
  \label{eq:complete-bell-egf}
\end{align}
\end{theorem}
\begin{proof}
Expand $C(z)^k/k!$ multinomially.  If $j_i$ components have size $i$, then
$\sum_i j_i=k$ and $\sum_iij_i=n$, and the coefficient of $z^n$ is
\[
 \prod_i\frac1{j_i!}\left(\frac{x_i}{i!}\right)^{j_i}.
\]
Multiplication by $n!$ labels the atoms.  This proves the first formula.  Summing it
with weight $u^k$ proves the second, and setting $u=1$ proves the third.
\end{proof}

\begin{corollary}[Partition-type count]\label{cor:partition-type}
The number of partitions of an $n$-element labelled set with exactly $j_i$ blocks of
size $i$ is
\begin{equation}\label{eq:partition-type-count}
 \frac{n!}{\prod_{i\ge1}(i!)^{j_i}j_i!},
 \qquad \sum_iij_i=n.
\end{equation}
\end{corollary}
\begin{proof}
Order the $n$ labels, cut the ordering into blocks of prescribed sizes, then divide
by $i!$ for the internal order of each size-$i$ block and by $j_i!$ for the order of
blocks of equal size.
\end{proof}

\part{Stirling Numbers, Lah Numbers, and Basis Changes}

\chapter{The two Stirling triangles}
\label{ch:stirling-def}

\section{First-kind numbers: cycles and factorial coefficients}
\begin{definition}
The signed Stirling numbers of the first kind $\SignedStirlingFirstKind{n}{k}$ and the unsigned numbers
$\UnsignedStirlingFirstKind nk$ are defined by
\begin{align}
 \FallingFactorial{x}{n}&=\sum_{k=0}^n \SignedStirlingFirstKind{n}{k}x^k,\label{eq:signed-first-def}\\
 \RisingFactorial{x}{n}&=\sum_{k=0}^n\UnsignedStirlingFirstKind nkx^k,\label{eq:unsigned-first-def}\\
 \SignedStirlingFirstKind{n}{k}&=(-1)^{n-k}\UnsignedStirlingFirstKind nk.\label{eq:first-sign}
\end{align}
\end{definition}
The sign relation follows from
$\RisingFactorial{x}{n}=(-1)^n\FallingFactorial{-x}{n}$.

\begin{theorem}[Permutation interpretation]\label{thm:first-cycle}
$\UnsignedStirlingFirstKind nk$ is the number of permutations of $[n]$ having exactly $k$ cycles.
Consequently,
\begin{equation}\label{eq:first-recurrence}
 \UnsignedStirlingFirstKind{n+1}{k}=n\UnsignedStirlingFirstKind nk+\UnsignedStirlingFirstKind n{k-1},
\end{equation}
with $\UnsignedStirlingFirstKind00=1$ and zero boundary values outside $0\le k\le n$.
\end{theorem}
\begin{proof}
Given a permutation of $[n+1]$ with $k$ cycles, either $n+1$ is a singleton cycle,
leaving a permutation of $[n]$ with $k-1$ cycles, or it is inserted immediately
after one of the $n$ existing elements in the cyclic notation of a permutation of
$[n]$ with $k$ cycles.  These alternatives are disjoint and reversible.
The same recurrence follows algebraically by multiplying
$\RisingFactorial{x}{n}$ by $x+n$ and comparing coefficients.
\end{proof}

For $n=3$, the cycle counts are $\UnsignedStirlingFirstKind33=1$, $\UnsignedStirlingFirstKind32=3$, and
$\UnsignedStirlingFirstKind31=2$, corresponding to the identity, the three transpositions, and the two
$3$-cycles.  Summing over the possible number of cycles gives
\begin{equation}\label{eq:first-row-sum}
 \sum_{k=0}^n\UnsignedStirlingFirstKind nk=n!,
\end{equation}
which also follows from \cref{eq:unsigned-first-def} at $x=1$.

\section{Second-kind numbers: set partitions and surjections}
\begin{definition}
The Stirling number of the second kind $\StirlingSecondKind nk$ is the number of partitions of
$[n]$ into exactly $k$ nonempty blocks.
\end{definition}

\begin{theorem}[Second-kind recurrence]\label{thm:second-recurrence}
For $n,k\ge0$,
\begin{equation}\label{eq:second-recurrence}
 \StirlingSecondKind{n+1}{k}=k\StirlingSecondKind nk+\StirlingSecondKind n{k-1},
\end{equation}
with $\StirlingSecondKind00=1$ and zero boundary values outside $0\le k\le n$.
\end{theorem}
\begin{proof}
In a partition of $[n+1]$ into $k$ blocks, the new label either joins one of the $k$
blocks of a partition of $[n]$, or forms a singleton block beside a partition of
$[n]$ into $k-1$ blocks.
\end{proof}

\begin{theorem}[Surjection formula]\label{thm:second-explicit}
For all $n,k\ge0$,
\begin{equation}\label{eq:second-explicit}
 \StirlingSecondKind nk=\frac1{k!}\sum_{j=0}^k(-1)^{k-j}\binom kjj^n
 =\sum_{j=0}^k\frac{(-1)^{k-j}j^n}{(k-j)!j!}.
\end{equation}
\end{theorem}
\begin{proof}
There are $k!\StirlingSecondKind nk$ surjections from $[n]$ onto a fixed $k$-element codomain:
a set partition into fibres followed by an ordering of the fibres.  Inclusion--
exclusion over the set of omitted codomain values gives
$\sum_{j=0}^k(-1)^{k-j}\binom kjj^n$.  Division by $k!$ proves the formula.
\end{proof}

The simple diagonals follow at once:
\begin{equation}\label{eq:second-simple}
 \StirlingSecondKind nn=1,\qquad \StirlingSecondKind n1=1\ (n\ge1),\qquad
 \StirlingSecondKind n{n-1}=\binom n2,\qquad
 \StirlingSecondKind n2=2^{n-1}-1.
\end{equation}
The last identity chooses a nonempty proper subset as one block and divides by two
for exchanging the blocks.


\section{Alternative and reverse recurrences}
The ordinary recurrence builds each row from the preceding one.  Several less
familiar identities in the source archive instead recover an entry from entries to
its right or from earlier entries in the same column.  They were labelled there as
conjectures; in fact they are exact identities.

\begin{theorem}[Alternative second-kind recurrences]
\label{thm:second-reverse-recurrences}
For $1\le k<n$,
\begin{align}
 \StirlingSecondKind nk
 &=\frac{k^n}{k!}-\sum_{r=1}^{k-1}\frac{\StirlingSecondKind nr}{(k-r)!},
 \label{eq:second-triangular-explicit}\\
 (n-k)\StirlingSecondKind nk
 &=\sum_{j=2}^{n-k+1}(j-2)!\binom{-k}{j}
   \StirlingSecondKind n{k+j-1},
 \label{eq:second-reverse-row}\\
 (n-k)\StirlingSecondKind nk
 &=\sum_{j=2}^{n-k+1}(-1)^j\binom nj
   \StirlingSecondKind{n-j+1}{k}.
 \label{eq:second-reverse-column}
\end{align}
The terminal value in either reverse recursion is $\StirlingSecondKind nn=1$.
\end{theorem}
\begin{proof}
Substitute $x=k$ in the basis identity
$k^n=\sum_r\StirlingSecondKind nr\FallingFactorial{k}{r}$.  Since
$\FallingFactorial{k}{r}=k!/(k-r)!$ for $r\le k$, division by $k!$ and isolation of the
$r=k$ term gives \eqref{eq:second-triangular-explicit}.

For the row reversal, use the bivariate exponential generating function
\[
 F(x,y)=\exp\!\bigl(y(\EulerE^x-1)\bigr)
 =\sum_{n,k\ge0}\StirlingSecondKind nk\frac{x^n}{n!}y^k,
 \qquad u=\EulerE^x-1.
\]
The generating function of the left side of
\eqref{eq:second-reverse-row} is
\[
 xF_x-yF_y=yF\bigl(xe^x-u\bigr).
\]
On the other hand, because
$\binom{-k}{j}=(-1)^j\binom{k+j-1}{j}$, the generating function of its right side is
\begin{align*}
 y\sum_{j\ge2}\frac{(-1)^j}{j(j-1)}\,\partial_y^jF
 &=yF\sum_{j\ge2}\frac{(-1)^ju^j}{j(j-1)}\\
 &=yF\bigl((1+u)\log(1+u)-u\bigr)\\
 &=yF\bigl(xe^x-u\bigr).
\end{align*}
Here the middle series is obtained by differentiating it twice, summing the
geometric series, and restoring the two zero initial coefficients.  Equality of
coefficients proves \eqref{eq:second-reverse-row}.

For the column reversal fix $k$ and put
\[
 F_k(x)=\sum_{n\ge k}\StirlingSecondKind nk\frac{x^n}{n!}
       =\frac{(\EulerE^x-1)^k}{k!}.
\]
The two sides of \eqref{eq:second-reverse-column}, after multiplication by
$x^n/n!$ and summation over $n$, become respectively
\[
 xF_k'(x)-kF_k(x)
 \quad\text{and}\quad
 (\EulerE^{-x}-1+x)F_k'(x).
\]
They are equal because direct differentiation gives
$(1-\EulerE^{-x})F_k'(x)=kF_k(x)$.
\end{proof}

The factor interchange mentioned in the source produces a parallel pair for the
unsigned first-kind triangle.

\begin{theorem}[Reverse recurrences of the first kind]
\label{thm:first-reverse-recurrences}
For $1\le k<n$,
\begin{align}
 (n-k)\UnsignedStirlingFirstKind nk
 &=\sum_{j=2}^{n-k+1}(-1)^j\binom{-k}{j}
   \UnsignedStirlingFirstKind n{k+j-1},
 \label{eq:first-reverse-row}\\
 (n-k)\UnsignedStirlingFirstKind nk
 &=\sum_{j=2}^{n-k+1}(j-2)!\binom nj
   \UnsignedStirlingFirstKind{n-j+1}{k}.
 \label{eq:first-reverse-column}
\end{align}
\end{theorem}
\begin{proof}
Let $C_n(y)=\RisingFactorial{y}{n}=\sum_r\UnsignedStirlingFirstKind nr y^r$.  Since
\[
 C_n(1+y)=\frac{\RisingFactorial{y}{n+1}}{y},
\]
the coefficient of $y^{k-1}$ on the two sides is, respectively,
\[
 \sum_{r\ge k-1}\binom r{k-1}\UnsignedStirlingFirstKind nr
 \quad\text{and}\quad
 \UnsignedStirlingFirstKind{n+1}{k}=n\UnsignedStirlingFirstKind nk+\UnsignedStirlingFirstKind n{k-1}.
\]
Remove the terms $r=k-1$ and $r=k$ and put $r=k+j-1$.  Because
$(-1)^j\binom{-k}{j}=\binom{k+j-1}{j}$, the remainder is precisely
\eqref{eq:first-reverse-row}.

For fixed $k$, set
\[
 C_k(x)=\sum_{n\ge k}\UnsignedStirlingFirstKind nk\frac{x^n}{n!}
       =\frac{[-\log(1-x)]^k}{k!}.
\]
The exponential generating function of the right side of
\eqref{eq:first-reverse-column} is
\[
 C_k'(x)\sum_{j\ge2}\frac{x^j}{j(j-1)}
 =\bigl(x+(1-x)\log(1-x)\bigr)C_k'(x).
\]
The left side has generating function $xC_k'-kC_k$.  Their equality reduces to
\[
 -(1-x)\log(1-x)\,C_k'(x)=kC_k(x),
\]
which follows immediately from the displayed formula for $C_k$.
\end{proof}
\section{Tables and common boundary conventions}
The first rows are
\[
\begin{array}{c|rrrrrrr}
 n\backslash k&0&1&2&3&4&5&6\\\hline
0&1\\
1&0&1\\
2&0&1&1\\
3&0&2&3&1\\
4&0&6&11&6&1\\
5&0&24&50&35&10&1\\
6&0&120&274&225&85&15&1
\end{array}
\qquad
\begin{array}{c|rrrrrrr}
 n\backslash k&0&1&2&3&4&5&6\\\hline
0&1\\
1&0&1\\
2&0&1&1\\
3&0&1&3&1\\
4&0&1&7&6&1\\
5&0&1&15&25&10&1\\
6&0&1&31&90&65&15&1
\end{array}
\]
for $\UnsignedStirlingFirstKind nk$ and $\StirlingSecondKind nk$, respectively.

\chapter{Factorial bases and inverse matrices}
\label{ch:bases}

\section{The basic changes of basis}
\begin{theorem}[Stirling basis transformations]\label{thm:stirling-basis}
For every $n\ge0$,
\begin{align}
 \FallingFactorial{x}{n}&=\sum_{k=0}^n\SignedStirlingFirstKind{n}{k}x^k,\label{eq:fall-to-power}\\
 \RisingFactorial{x}{n}&=\sum_{k=0}^n\UnsignedStirlingFirstKind nkx^k,\label{eq:rise-to-power}\\
 x^n&=\sum_{k=0}^n\StirlingSecondKind nk\FallingFactorial{x}{k},\label{eq:power-to-fall}\\
 x^n&=\sum_{k=0}^n(-1)^{n-k}\StirlingSecondKind nk\RisingFactorial{x}{k}.\label{eq:power-to-rise}
\end{align}
Equivalently,
\begin{equation}\label{eq:binomial-basis}
 x^n=\sum_{k=0}^n\StirlingSecondKind nk k!\binom{x}{k},\qquad
 \binom{x}{n}=\frac1{n!}\sum_{k=0}^n\SignedStirlingFirstKind{n}{k}x^k.
\end{equation}
\end{theorem}
\begin{proof}
The first two identities are definitions.  In the Newton expansion
\cref{eq:newton-expansion} for $p(x)=x^n$, one has
\[
 \Delta^kx^n\big|_{x=0}=\sum_{j=0}^k(-1)^{k-j}\binom kjj^n
 =k!\StirlingSecondKind nk
\]
by \cref{eq:second-explicit}.  This gives \cref{eq:power-to-fall}.  Replacing $x$
by $-x$ and multiplying by $(-1)^n$ gives \cref{eq:power-to-rise}.  Division of
$\FallingFactorial{x}{n}$ by $n!$ gives the second binomial-basis formula.
\end{proof}


\begin{corollary}[Shifted factorial evaluations]
\label{cor:shifted-stirling-evaluations}
For every $n\ge0$,
\begin{align}
 \sum_{k=1}^{n+1}\StirlingSecondKind{n+1}{k}\FallingFactorial{x-1}{k-1}&=x^n,
 \label{eq:shifted-power-to-fall}\\
 \sum_{k=0}^{n}\StirlingSecondKind nk\FallingFactorial{n}{k}&=n^n.
 \label{eq:stirling-n-to-n}
\end{align}
\end{corollary}
\begin{proof}
Apply \eqref{eq:power-to-fall} with exponent $n+1$ and use
$\FallingFactorial{x}{k}=x\FallingFactorial{x-1}{k-1}$.  Both sides are divisible by $x$ as
polynomials, and cancellation gives \eqref{eq:shifted-power-to-fall}.  Formula
\eqref{eq:stirling-n-to-n} is \eqref{eq:power-to-fall} evaluated at $x=n$.
\end{proof}
\begin{corollary}[Inverse matrices]\label{cor:stirling-inverse}
The lower triangular matrices $s=(\SignedStirlingFirstKind{n}{k})$ and $S=(\StirlingSecondKind nk)$ are mutual inverses:
\begin{align}
 \sum_{j=k}^n\SignedStirlingFirstKind{n}{j}\StirlingSecondKind jk&=\delta_{nk},\label{eq:stirling-inv1}\\
 \sum_{j=k}^n\StirlingSecondKind nj \SignedStirlingFirstKind{j}{k}&=\delta_{nk}.\label{eq:stirling-inv2}
\end{align}
Equivalently,
\begin{align}
 \sum_{j=k}^n(-1)^{n-j}\UnsignedStirlingFirstKind nj\StirlingSecondKind jk&=\delta_{nk},\\
 \sum_{j=k}^n(-1)^{j-k}\StirlingSecondKind nj\UnsignedStirlingFirstKind jk&=\delta_{nk}.
\end{align}
\end{corollary}
\begin{proof}
Expand $\FallingFactorial{x}{n}$ into powers and each power back into falling factorials.
Uniqueness of coordinates in the factorial basis gives the first matrix product.
The reverse substitution gives the second.
\end{proof}

\section{Power sums by basis adaptation}
The factorial basis makes discrete antidifferentiation immediate.  Combining
\cref{eq:power-to-fall,eq:finite-diff-falling} yields
\begin{equation}\label{eq:power-sum-stirling}
 \sum_{i=0}^{N}i^m
 =\sum_{k=0}^m\StirlingSecondKind mk\frac{\FallingFactorial{N+1}{k+1}}{k+1}.
\end{equation}
For $m=4$ this is
\[
 \frac12\FallingFactorial{N+1}{2}+\frac73\FallingFactorial{N+1}{3}
 +\frac64\FallingFactorial{N+1}{4}+\frac15\FallingFactorial{N+1}{5},
\]
which expands to the familiar quartic power-sum polynomial.


\section{Factorial moments: Poisson variables and permutation fixed points}
The change from ordinary powers to falling factorials is especially effective for
integer-valued random variables.

\begin{theorem}[Poisson moments]
\label{thm:poisson-stirling-moments}
If $X$ has the Poisson distribution of mean $\lambda$, then
\begin{equation}
\label{eq:poisson-stirling-moments}
 \Expectation X^n=\sum_{k=0}^{n}\StirlingSecondKind nk\lambda^k.
\end{equation}
In particular, $\Expectation X^n=\BellNumber n$ when $\lambda=1$.
\end{theorem}
\begin{proof}
The $k$th falling-factorial moment is
\[
 \Expectation\FallingFactorial{X}{k}
 =\EulerE^{-\lambda}\sum_{r\ge k}\frac{\FallingFactorial r k\lambda^r}{r!}
 =\lambda^ke^{-\lambda}\sum_{q\ge0}\frac{\lambda^q}{q!}
 =\lambda^k.
\]
Take expectations in the basis identity
$X^n=\sum_k\StirlingSecondKind nk\FallingFactorial Xk$.
\end{proof}

\begin{theorem}[Moments of the number of fixed points]
\label{thm:fixed-point-stirling-moments}
Let $Y_m$ be the number of fixed points of a uniformly random permutation of
$[m]$.  Then
\begin{equation}
\label{eq:fixed-point-stirling-moments}
 \Expectation Y_m^n=\sum_{k=0}^{m}\StirlingSecondKind nk
 =\sum_{k=0}^{\min(m,n)}\StirlingSecondKind nk.
\end{equation}
Thus, when $m\ge n$, the moment equals the Bell number $\BellNumber n$.
\end{theorem}
\begin{proof}
The random variable $\FallingFactorial{Y_m}{k}$ counts ordered $k$-tuples of distinct fixed
points.  There are $\FallingFactorial m k$ candidate tuples, and a specified tuple is fixed
with probability $(m-k)!/m!$.  Hence
\[
 \Expectation\FallingFactorial{Y_m}{k}=1\quad(0\le k\le m),
 \qquad
 \Expectation\FallingFactorial{Y_m}{k}=0\quad(k>m).
\]
Taking expectations in \eqref{eq:power-to-fall} proves the formula.  The upper
limit $m$ is essential: it records the finite size of the permutation, even when
$n>m$.
\end{proof}
\section{Lah numbers}
\begin{definition}
For $1\le k\le n$, the unsigned Lah number is
\begin{equation}\label{eq:lah-def}
 \LahNumber{n}{k}=\binom{n-1}{k-1}\frac{n!}{k!},
\end{equation}
with $\LahNumber{0}{0}=1$ and zero otherwise outside the triangle.
\end{definition}
Combinatorially, $\LahNumber{n}{k}$ counts partitions of $[n]$ into $k$ nonempty linearly
ordered lists, with the collection of lists unordered.

\begin{theorem}[Conversion between rising and falling factorials]\label{thm:lah-conv}
\begin{align}
 \RisingFactorial{x}{n}&=\sum_{k=0}^n\LahNumber{n}{k}\FallingFactorial{x}{k},\label{eq:lah-rise-fall}\\
 \FallingFactorial{x}{n}&=\sum_{k=0}^n(-1)^{n-k}\LahNumber{n}{k}\RisingFactorial{x}{k},\label{eq:lah-fall-rise}\\
 \sum_{j=k}^n(-1)^{j-k}\LahNumber{n}{j}\LahNumber{j}{k}&=\delta_{nk}.\label{eq:lah-inv}
\end{align}
Moreover,
\begin{equation}\label{eq:lah-stirling-product}
 \LahNumber{n}{k}=\sum_{j=k}^n\UnsignedStirlingFirstKind nj\StirlingSecondKind jk.
\end{equation}
\end{theorem}
\begin{proof}
A direct generating-function calculation gives
\[
 \sum_{n\ge k}\LahNumber{n}{k}\frac{z^n}{n!}
 =\frac1{k!}\left(\frac{z}{1-z}\right)^k.
\]
Therefore
\[
 \sum_{n\ge0}\left(\sum_k\LahNumber{n}{k}\FallingFactorial{x}{k}\right)\frac{z^n}{n!}
 =\sum_{k\ge0}\binom{x}{k}\left(\frac{z}{1-z}\right)^k
 =(1-z)^{-x}
 =\sum_{n\ge0}\RisingFactorial{x}{n}\frac{z^n}{n!},
\]
proving \cref{eq:lah-rise-fall}.  Replace $x$ by $-x$ to obtain the inverse formula,
and compose the two transformations to obtain \cref{eq:lah-inv}.  Finally, expand a
rising factorial into powers by first-kind numbers and powers into falling factorials
by second-kind numbers; uniqueness gives \cref{eq:lah-stirling-product}.
\end{proof}

The inequalities
\begin{equation}\label{eq:stirling-lah-ineq}
 \StirlingSecondKind nk\le\UnsignedStirlingFirstKind nk\le \LahNumber{n}{k}
\end{equation}
follow from explicit injections: a set partition may be canonically turned into a
permutation with one cycle per block by writing each block with its smallest member
first and arranging cycles by minima; forgetting the cyclic identifications embeds
permutations with $k$ cycles into unordered collections of $k$ lists.


\chapter{Generating functions, transforms, and operators}
\label{ch:stirling-gf}

\section{Exponential generating functions}
\begin{theorem}[Column generating functions]\label{thm:stirling-egfs}
For every fixed $k\ge0$,
\begin{align}
 \sum_{n\ge k}\SignedStirlingFirstKind{n}{k}\frac{z^n}{n!}
   &=\frac{\log(1+z)^k}{k!},\label{eq:signed-first-egf}\\
 \sum_{n\ge k}\UnsignedStirlingFirstKind nk\frac{z^n}{n!}
   &=\frac{[-\log(1-z)]^k}{k!},\label{eq:unsigned-first-egf}\\
 \sum_{n\ge k}\StirlingSecondKind nk\frac{z^n}{n!}
   &=\frac{(\EulerE^z-1)^k}{k!}.\label{eq:second-egf}
\end{align}
The bivariate second-kind generating function is
\begin{equation}\label{eq:second-double-egf}
 \sum_{k\ge0}\sum_{n\ge k}\StirlingSecondKind nk y^k\frac{z^n}{n!}
 =\exp\!\bigl(y(\EulerE^z-1)\bigr).
\end{equation}
\end{theorem}
\begin{proof}
Expand
\[
 (1+z)^u=\EulerE^{u\log(1+z)}
 =\sum_{n\ge0}\frac{z^n}{n!}\sum_{k=0}^n\SignedStirlingFirstKind{n}{k}u^k
 =\sum_{k\ge0}\frac{u^k}{k!}\log(1+z)^k
\]
and compare coefficients of $u^k$.  Replacing $z$ by $-z$ and using the sign
relation gives \cref{eq:unsigned-first-egf}.  Formula \cref{eq:second-egf} is the
exponential formula for a set of $k$ nonempty labelled blocks; algebraically it also
follows by inserting \cref{eq:second-explicit} and summing exponentials.  Summing over
$k$ with weight $y^k$ proves \cref{eq:second-double-egf}.
\end{proof}

Differentiating \cref{eq:signed-first-egf} gives the useful identity
\begin{equation}\label{eq:log-over-one-plus}
 \frac{\log(1+z)^m}{1+z}
 =m!\sum_{r\ge0}\SignedStirlingFirstKind{r+1}{m+1}\frac{z^r}{r!},\qquad |z|<1,
\end{equation}
which is valid formally as well.

\section{A rational ordinary generating function}
\begin{theorem}[Fixed-$k$ ordinary generating function]\label{thm:second-ogf}
For $k\ge0$,
\begin{equation}\label{eq:second-ogf}
 \sum_{n\ge k}\StirlingSecondKind nk x^{n-k}
 =\prod_{j=1}^k\frac1{1-jx}
 =\frac1{x^{k+1}\FallingFactorial{1/x}{k+1}}.
\end{equation}
Consequently,
\begin{equation}\label{eq:second-complete-symmetric}
 \StirlingSecondKind nk=h_{n-k}(1,2,\ldots,k)
 =\sum_{c_1+\cdots+c_k=n-k}1^{c_1}2^{c_2}\cdots k^{c_k},
\end{equation}
where $h_r$ is a complete homogeneous symmetric polynomial.
\end{theorem}
\begin{proof}
Let $F_k(x)=\sum_{r\ge0}\StirlingSecondKind{k+r}{k}x^r$.  The recurrence
\cref{eq:second-recurrence}, with $n=k+r-1$, gives
$F_k(x)=(1-kx)^{-1}F_{k-1}(x)$ and $F_0=1$.  Iteration proves the product.
Expanding each geometric factor independently proves the complete-homogeneous sum.
The factorial form follows by direct factorization.
\end{proof}

A useful rearrangement of the count of all maps $[n]\to[k]$ is
\begin{equation}\label{eq:second-alternate-rec}
 \StirlingSecondKind nk=\frac{k^n}{k!}-\sum_{r=1}^{k-1}\frac{\StirlingSecondKind nr}{(k-r)!}.
\end{equation}
Indeed, a map using exactly $r$ values can be chosen in
$\binom kr r!\StirlingSecondKind nr$ ways; divide the identity
$k^n=\sum_r\binom kr r!\StirlingSecondKind nr$ by $k!$.

\section{The Stirling transform}
\begin{theorem}[Stirling transform and its inverse]\label{thm:stirling-transform}
Let
\[
 g_n=\sum_{k=0}^n\StirlingSecondKind nk f_k.
\]
Then
\begin{equation}\label{eq:stirling-transform-inverse}
 f_n=\sum_{k=0}^n\SignedStirlingFirstKind{n}{k}g_k
 =\sum_{k=0}^n(-1)^{n-k}\UnsignedStirlingFirstKind nkg_k.
\end{equation}
For exponential generating functions,
\begin{equation}\label{eq:stirling-transform-egf}
 \ExponentialGeneratingFunctionOf{G}(z)=\ExponentialGeneratingFunctionOf{F}(\EulerE^z-1),\qquad
 \ExponentialGeneratingFunctionOf{F}(z)=\ExponentialGeneratingFunctionOf{G}(\log(1+z)).
\end{equation}
\end{theorem}
\begin{proof}
The inverse formula is exactly the matrix inversion in
\cref{cor:stirling-inverse}.  For the generating function,
\[
 \ExponentialGeneratingFunctionOf{G}(z)=\sum_kf_k\sum_{n\ge k}\StirlingSecondKind nk\frac{z^n}{n!}
 =\sum_kf_k\frac{(\EulerE^z-1)^k}{k!}=\ExponentialGeneratingFunctionOf{F}(\EulerE^z-1).
\]
Substitution by the compositional inverse $\log(1+z)$ gives the reverse relation.
\end{proof}

\section{Differential-operator normal ordering}
Let $D=d/dx$.  The Euler operator $xD$ is diagonal on monomials:
$(xD)x^m=mx^m$.  The operator $x^kD^k$ has eigenvalue $\FallingFactorial m k$.
Thus the scalar basis changes lift verbatim to operators.

\begin{theorem}[Stirling normal-ordering identities]\label{thm:normal-order}
For $n\ge0$,
\begin{align}
 (xD)^n&=\sum_{k=0}^n\StirlingSecondKind nk x^kD^k,\label{eq:normal1}\\
 x^nD^n&=\sum_{k=0}^n\SignedStirlingFirstKind{n}{k}(xD)^k
        =\FallingFactorial{xD}{n}.\label{eq:normal2}
\end{align}
\end{theorem}
\begin{proof}
Both sides of \cref{eq:normal1} act on $x^m$ as multiplication by
$m^n=\sum_k\StirlingSecondKind nk\FallingFactorial mk$.  Since monomials span the polynomial ring, the
operators are equal.  The inverse relation is proved identically using
$\FallingFactorial mn=\sum_k\SignedStirlingFirstKind{n}{k}m^k$.
\end{proof}

Because $E=\EulerE^D$ and $\Delta=\EulerE^D-1$, substitution of the mutually inverse series
$D=\log(1+\Delta)$ and $\Delta=\EulerE^D-1$ gives, whenever the operator series converge
(or formally on polynomials),
\begin{align}
 \frac{D^k}{k!}
   &=\sum_{n\ge k}\frac{\SignedStirlingFirstKind{n}{k}}{n!}\Delta^n,
   \label{eq:der-from-diff}\\
 \frac{\Delta^k}{k!}
   &=\sum_{n\ge k}\frac{\StirlingSecondKind nk}{n!}D^n.
   \label{eq:diff-from-der}
\end{align}
These are the derivative--finite-difference transformations listed in the archive.

\chapter{Symmetric functions, diagonals, and finite identities}
\label{ch:stirling-identities}

\section{Elementary and complete symmetric descriptions}
From \cref{eq:unsigned-first-def},
\begin{equation}\label{eq:first-elementary}
 \UnsignedStirlingFirstKind n{n-r}=e_r(0,1,\ldots,n-1)
 =\sum_{0\le i_1<\cdots<i_r<n}i_1\cdots i_r.
\end{equation}
Together with \cref{eq:second-complete-symmetric}, this identifies the first-kind
triangle with elementary symmetric functions and the second-kind triangle with
complete homogeneous symmetric functions.

The first few near-diagonals are
\begin{align}
 \UnsignedStirlingFirstKind n{n-1}&=\binom n2,\label{eq:first-diag1}\\
 \UnsignedStirlingFirstKind n{n-2}&=\frac{3n-1}{4}\binom n3,\label{eq:first-diag2}\\
 \UnsignedStirlingFirstKind n{n-3}&=\binom n2\binom n4.\label{eq:first-diag3}
\end{align}
\begin{proof}
The first is $e_1(0,\ldots,n-1)$.  Newton's identities
$2e_2=p_1^2-p_2$ and $6e_3=p_1^3-3p_1p_2+2p_3$, combined with the elementary sums
of first, second, and third powers, simplify to \cref{eq:first-diag2,eq:first-diag3}.
Equivalently, one may classify permutations by the nontrivial cycle types
$3$, $2+2$ for the second diagonal and $4$, $3+2$, $2+2+2$ for the third.
\end{proof}

\section{Harmonic-number formulae}
Factoring $(n-1)!$ from
$(x+1)(x+2)\cdots(x+n-1)$ gives
\begin{equation}\label{eq:first-harmonic-elementary}
 \frac{\UnsignedStirlingFirstKind n{k+1}}{(n-1)!}
 =e_k\left(1,\frac12,\ldots,\frac1{n-1}\right)
 =\sum_{1\le i_1<\cdots<i_k<n}\frac1{i_1\cdots i_k}.
\end{equation}
Write $H_n^{(r)}=\sum_{j=1}^nj^{-r}$ and $H_n=H_n^{(1)}$.  Newton's identities now
give
\begin{align}
 \UnsignedStirlingFirstKind n2&=(n-1)!H_{n-1},\label{eq:first-H1}\\
 \UnsignedStirlingFirstKind n3&=\frac{(n-1)!}{2}\left(H_{n-1}^2-H_{n-1}^{(2)}\right),\label{eq:first-H2}\\
 \UnsignedStirlingFirstKind n4&=\frac{(n-1)!}{6}\left(H_{n-1}^3-3H_{n-1}H_{n-1}^{(2)}
                   +2H_{n-1}^{(3)}\right).\label{eq:first-H3}
\end{align}
More generally,
\begin{equation}\label{eq:first-harmonic-bell}
 \frac{\UnsignedStirlingFirstKind{n+1}{k+1}}{n!}
 =\frac1{k!}\ExponentialCompleteBellPolynomial{k}\bigl(H_n,-1!H_n^{(2)},2!H_n^{(3)},\ldots,
              (-1)^{k-1}(k-1)!H_n^{(k)}\bigr).
\end{equation}
\begin{proof}
Take logarithms in
\[
 \prod_{j=1}^n\left(1+\frac{x}{j}\right)
 =\exp\left(\sum_{r\ge1}\frac{(-1)^{r-1}H_n^{(r)}}r x^r\right)
\]
and apply the complete Bell-polynomial generating function
\cref{eq:complete-bell-egf}.
\end{proof}

If $w(n,k)=k!\,\UnsignedStirlingFirstKind n{k+1}/(n-1)!$, then complete Bell recurrence yields
\begin{equation}\label{eq:w-recurrence}
 w(n,m)=\delta_{m0}+\sum_{r=0}^{m-1}(1-m)_rH_{n-1}^{(r+1)}w(n,m-1-r),
\end{equation}
where $(1-m)_r$ is falling.  Conversely, taking a logarithm recovers the harmonic
power sums:
\begin{equation}\label{eq:harmonic-inversion}
 H_n^{(m)}=-m\CoefficientExtraction xm
 \log\left(1+\sum_{k\ge1}\frac{\UnsignedStirlingFirstKind{n+1}{k+1}}{n!}(-x)^k\right).
\end{equation}
For example,
\begin{align}
 (n!)^2H_n^{(2)}
 &=\UnsignedStirlingFirstKind{n+1}{2}^{\!2}
   -2\UnsignedStirlingFirstKind{n+1}{1}\UnsignedStirlingFirstKind{n+1}{3},\label{eq:H2-invert}\\
 (n!)^3H_n^{(3)}
 &=\UnsignedStirlingFirstKind{n+1}{2}^{\!3}
   -3\UnsignedStirlingFirstKind{n+1}{1}\UnsignedStirlingFirstKind{n+1}{2}\UnsignedStirlingFirstKind{n+1}{3}
   +3\UnsignedStirlingFirstKind{n+1}{1}^{\!2}\UnsignedStirlingFirstKind{n+1}{4}.
 \label{eq:H3-invert}
\end{align}

\section{Convolution and summation identities}
The following paired identities expose the parallel structure of the two triangles.

\begin{theorem}[First- and second-kind summations]\label{thm:paired-sums}
For valid nonnegative indices,
\begin{align}
 \UnsignedStirlingFirstKind{n+1}{k+1}
   &=\sum_{j=k}^n\UnsignedStirlingFirstKind nj\binom jk
    =\sum_{j=k}^n\frac{n!}{j!}\UnsignedStirlingFirstKind jk,\label{eq:first-two-sums}\\
 \StirlingSecondKind{n+1}{k+1}
   &=\sum_{j=k}^n\binom nj\StirlingSecondKind jk
    =\sum_{j=k}^n(k+1)^{n-j}\StirlingSecondKind jk,\label{eq:second-two-sums}\\
 \UnsignedStirlingFirstKind{n+k+1}{n}
   &=\sum_{j=0}^k(n+j)\UnsignedStirlingFirstKind{n+j}{j},\label{eq:first-hockey}\\
 \StirlingSecondKind{n+k+1}{k}
   &=\sum_{j=0}^kj\StirlingSecondKind{n+j}{j}.\label{eq:second-hockey}
\end{align}
Moreover, for $\ell,m\ge0$,
\begin{align}
 \binom{\ell+m}{\ell}\UnsignedStirlingFirstKind n{\ell+m}
 &=\sum_j\binom nj\UnsignedStirlingFirstKind j\ell\UnsignedStirlingFirstKind{n-j}m,
 \label{eq:first-convolution}\\
 \binom{\ell+m}{\ell}\StirlingSecondKind n{\ell+m}
 &=\sum_j\binom nj\StirlingSecondKind j\ell\StirlingSecondKind{n-j}m.
 \label{eq:second-convolution}
\end{align}
\end{theorem}
\begin{proof}
For the first equality in \cref{eq:first-two-sums}, mark $k$ cycles of a permutation
of $[n]$ and insert $n+1$ into the canonical construction; algebraically it is the
coefficient of $x^k$ in $\frac{d}{dx}\RisingFactorial{x}{n}$.  The second equality follows by
iterating recurrence \cref{eq:first-recurrence} in the $n$ direction.  For
\cref{eq:second-two-sums}, distinguish the elements outside the block containing
$n+1$ for the first equality, and add the $n-k$ remaining elements one at a time to
one of $k+1$ existing blocks for the second.  The hockey-stick forms telescope the
respective recurrences.  Finally, in \cref{eq:first-convolution}, choose which labels
belong to the $\ell$ marked cycles; in \cref{eq:second-convolution}, choose which
labels belong to the $\ell$ marked blocks.  The binomial factor on the left forgets
which of the total $\ell+m$ components were marked.
\end{proof}

Binomial inversion of the first identity in \cref{eq:first-two-sums} gives
\begin{equation}\label{eq:first-inverted-sum}
 \UnsignedStirlingFirstKind nm=\sum_{k=m}^n(-1)^{k-m}\binom km\UnsignedStirlingFirstKind{n+1}{k+1}.
\end{equation}
The alternative form
\begin{equation}\label{eq:first-factorial-sum}
 \UnsignedStirlingFirstKind{n+1}{m+1}=\sum_{k=m}^n\frac{n!}{k!}\UnsignedStirlingFirstKind km
\end{equation}
is the second equality in \cref{eq:first-two-sums}.  All finite sums displayed in the
source follow from these and the two convolution identities by relabelling indices.

\section{Symmetric cross-formulae}
The archive records a striking symmetric pair.  For $0\le k\le n$,
\begin{align}
 \UnsignedStirlingFirstKind nk
 &=\sum_{j=n}^{2n-k}(-1)^{j-k}\binom{j-1}{k-1}\binom{2n-k}{j}
      \StirlingSecondKind{j-k}{j-n},\label{eq:symmetric-cross1}\\
 \StirlingSecondKind nk
 &=\sum_{j=n}^{2n-k}(-1)^{j-k}\binom{j-1}{k-1}\binom{2n-k}{j}
      \UnsignedStirlingFirstKind{j-k}{j-n}.\label{eq:symmetric-cross2}
\end{align}
\begin{proof}
Set $r=n-k$ and use the symmetric-function descriptions
$\UnsignedStirlingFirstKind nk=e_r(0,1,\ldots,n-1)$ and
$\StirlingSecondKind nk=h_r(1,2,\ldots,k)$.  The involution $\omega$ on the ring of symmetric
functions sends $e_r$ to $h_r$ and is implemented on generating functions by
$E(t)H(-t)=1$.  Expand the finite alphabets after adjoining and removing the common
arithmetic progression, then apply the binomial translation formula
$h_r(X+c)=\sum_j\binom{m+r-1}{r-j}c^{r-j}h_j(X)$ and its elementary analogue.
Collecting coefficients gives \cref{eq:symmetric-cross1}; applying $\omega$ gives
\cref{eq:symmetric-cross2}.  Equivalently, direct substitution of the
inclusion--exclusion formula \cref{eq:second-explicit} and two applications of
Vandermonde's identity produce the same finite sum.
\end{proof}


\section{An explicit two-sum formula of the first kind}
Unlike the second-kind inclusion--exclusion formula, the first-kind triangle has no
comparably short single finite sum.  Substitution into the symmetric cross-formula
nevertheless gives a completely explicit double sum.

\begin{theorem}[Explicit first-kind double sum]
\label{thm:first-double-sum}
For $1\le k\le n$,
\begin{equation}
\label{eq:first-double-sum}
 \UnsignedStirlingFirstKind nk
 =\sum_{j=n}^{2n-k}\binom{j-1}{k-1}\binom{2n-k}{j}
   \sum_{m=0}^{j-n}
   \frac{(-1)^{m+n-k}m^{j-k}}{m!\,(j-n-m)!}.
\end{equation}
\end{theorem}
\begin{proof}
Start from \eqref{eq:symmetric-cross1} and apply
\eqref{eq:second-explicit} with upper index $j-k$ and lower index $j-n$:
\[
 \StirlingSecondKind{j-k}{j-n}
 =\sum_{m=0}^{j-n}
   \frac{(-1)^{j-n-m}m^{j-k}}{(j-n-m)!\,m!}.
\]
Multiplication by the outer sign $(-1)^{j-k}$ gives
$(-1)^{2j-k-n-m}=(-1)^{m+n-k}$, because the exponents differ by the even integer
$2(j-m-n+k)$.  Substitution proves \eqref{eq:first-double-sum}.
\end{proof}

\section{A partition formula for signed near-diagonal coefficients}
The near-diagonal coefficient $\SignedStirlingFirstKind{n}{n-p}$ is a polynomial in $n$ of degree $2p$.
The following finite formula makes every monomial contribution explicit.

\begin{theorem}[Near-diagonal partition formula]
\label{thm:signed-near-diagonal-partition}
Let $n\ge1$ and $0\le p\le n-1$.  Then
\begin{equation}
\label{eq:signed-near-diagonal-partition}
 \SignedStirlingFirstKind{n}{n-p}=\frac{1}{(n-p-1)!}
 \sum_{\substack{k_1,\ldots,k_p\ge0\\
                  \sum_{r=1}^{p}r k_r=p}}
 (-1)^K
 \frac{(n+K-1)!}
 {\displaystyle\prod_{r=1}^{p}k_r!\,(r+1)!^{k_r}},
 \qquad K=\sum_{r=1}^{p}k_r.
\end{equation}
For $p=0$ the sum and product are empty and the formula reads $\SignedStirlingFirstKind{n}{n}=1$.
\end{theorem}
\begin{proof}
The signed first-kind exponential generating function gives
\begin{equation}
\label{eq:near-diagonal-log-coefficient}
 \SignedStirlingFirstKind{n}{n-p}=\frac{n!}{(n-p)!}
 [z^p]\left(\frac{\log(1+z)}{z}\right)^{n-p}.
\end{equation}
Put $t=\log(1+z)$ and
\[
 \phi(t)=\frac{t}{\EulerE^t-1}.
\]
Then $t=z\phi(t)$ and $t/z=\phi(t)$.  For $p\ge1$, the
Lagrange--B\"urmann formula applied to $H(t)=\phi(t)^{n-p}$ gives
\begin{align*}
 [z^p]H(t(z))
 &=\frac1p[t^{p-1}]H'(t)\phi(t)^p\\
 &=\frac{n-p}{p}[t^{p-1}]\phi'(t)\phi(t)^{n-1}\\
 &=\frac{n-p}{n}[t^p]\phi(t)^n.
\end{align*}
The same conclusion is immediate for $p=0$.  Hence
\begin{equation}
\label{eq:near-diagonal-bernoulli-power}
 \SignedStirlingFirstKind{n}{n-p}=\frac{(n-1)!}{(n-p-1)!}
 [t^p]\left(\frac{t}{\EulerE^t-1}\right)^n.
\end{equation}
Now write
\[
 A(t)=\sum_{r\ge1}\frac{t^r}{(r+1)!}
     =\frac{\EulerE^t-1-t}{t},
 \qquad \frac{t}{\EulerE^t-1}=(1+A(t))^{-1}.
\]
The negative-binomial expansion yields
\[
 (n-1)![t^p](1+A)^{-n}
 =\sum_{K\ge0}(-1)^K\frac{(n+K-1)!}{K!}[t^p]A(t)^K.
\]
Finally, the multinomial theorem gives
\[
 [t^p]A(t)^K
 =\sum_{\substack{k_1+\cdots+k_p=K\\
                   \sum r k_r=p}}
   \frac{K!}{\prod_{r=1}^{p}k_r!}
   \prod_{r=1}^{p}\frac1{(r+1)!^{k_r}}.
\]
Insert this in \eqref{eq:near-diagonal-bernoulli-power} to obtain
\eqref{eq:signed-near-diagonal-partition}.
\end{proof}

The unsigned value is
$\UnsignedStirlingFirstKind n{n-p}=(-1)^p \SignedStirlingFirstKind{n}{n-p}$.  For $p=1,2,3$ the theorem reduces to the
closed forms in \cref{eq:first-diag1,eq:first-diag2,eq:first-diag3}; its advantage is that it works
at arbitrary fixed distance from the diagonal without a new case analysis.
\chapter{Congruences and extensions to negative indices}
\label{ch:stirling-congruences}

\section{Congruences for first-kind numbers}
For a prime $p$,
\begin{equation}\label{eq:prime-first-divisibility}
 p\mid\UnsignedStirlingFirstKind pk\qquad(1<k<p).
\end{equation}
Indeed, in $\FiniteField_p[x]$,
\[
 \RisingFactorial{x}{p}=\prod_{a\in\FiniteField_p}(x+a)=x^p-x,
\]
so every intermediate coefficient vanishes.

More generally, reducing each factor modulo $2$ or $3$ gives exact coefficient rules:
\begin{align}
 \UnsignedStirlingFirstKind nm
 &\equiv \CoefficientExtraction xm x^{\Ceiling{n/2}}(x+1)^{\Floor{n/2}}
  =\binom{\Floor{n/2}}{m-\Ceiling{n/2}}\pmod2,
 \label{eq:first-mod2}\\
 \UnsignedStirlingFirstKind nm
 &\equiv \CoefficientExtraction xm
 x^{\Ceiling{n/3}}(x+1)^{\Ceiling{(n-1)/3}}(x+2)^{\Floor{n/3}}
 \pmod3.\label{eq:first-mod3}
\end{align}
\begin{proof}
Use $\RisingFactorial{x}{n}=\prod_{j=0}^{n-1}(x+j)$ and count how often each residue class
occurs among $0,1,\ldots,n-1$.  The displayed binomial form modulo $2$ is immediate;
expanding the last two factors gives the finite binomial sum written in the source
for modulus $3$.
\end{proof}

For fixed $m$, expansion of \cref{eq:first-mod2} yields the source's first four
parity formulae.  They are most cleanly read as integer-valued expressions multiplied
by Iverson brackets; for example
\[
 \UnsignedStirlingFirstKind n1\equiv 2^{n-2}[n\ge2]+[n=1]\pmod2,
\]
and similarly the quoted expressions for $m=2,3,4$.  Formula
\cref{eq:first-mod2} is both shorter and proves all of them simultaneously.

\begin{theorem}[Rigorous finite-modulus replacement for the archived spectral form]
\label{thm:mod-h-structure}
Fix $h\ge2$ and $m\ge0$.  The sequence $n\mapsto\UnsignedStirlingFirstKind nm\bmod h$ is generated by
the coefficient of $x^m$ in
\begin{equation}\label{eq:mod-h-block-product}
 \left(\prod_{r=0}^{h-1}(x+r)\right)^q\prod_{r=0}^{s-1}(x+r),
 \qquad n=qh+s,
\end{equation}
and therefore satisfies a linear recurrence over $\IntegerNumbers/h\IntegerNumbers$ and is eventually an
exponential-polynomial sequence over any finite splitting extension.
\end{theorem}
\begin{proof}
Group the factors of $\RisingFactorial{x}{n}$ into blocks of $h$ consecutive integers and
reduce each block modulo $h$.  For fixed $m$, coefficient extraction takes place in
a finite free module of polynomials of degree at most $m$, on which multiplication by
the fixed block polynomial is a linear endomorphism.  Cayley--Hamilton supplies a
linear recurrence.  Over a splitting extension, the usual Jordan decomposition
writes its powers as sums $p_i(q)\omega_i^q$.  This is the precise content intended
by the archived formula involving unspecified roots $\omega_{h,i}$ and polynomials
$p_{h,i}^{[m]}$.
\end{proof}

\section{Parity of second-kind numbers}
\begin{theorem}[Parity criterion]\label{thm:second-parity}
Let
\[
 z=n-\Ceiling{\frac{k+1}{2}},
 \qquad w=\Floor{\frac{k-1}{2}}.
\]
Then
\begin{equation}\label{eq:second-parity-binomial}
 \StirlingSecondKind nk\equiv\binom zw\pmod2.
\end{equation}
Equivalently,
\begin{equation}\label{eq:second-parity-bit}
 \StirlingSecondKind nk\equiv1\pmod2
 \quad\Longleftrightarrow\quad
 (n-k)\BitwiseAnd\Floor{\frac{k-1}{2}}=0.
\end{equation}
In particular, $\StirlingSecondKind{2n}{n}$ is odd precisely when $n$ is a power of two.
\end{theorem}
\begin{proof}
Reduce the ordinary generating function \cref{eq:second-ogf} modulo $2$.  The even
factors become $1$, while every odd $j$ contributes $(1-x)^{-1}$.  There are
$\Ceiling{k/2}$ odd integers from $1$ to $k$, so
\[
 \sum_{r\ge0}\StirlingSecondKind{k+r}{k}x^r
 \equiv(1-x)^{-\Ceiling{k/2}}.
\]
The coefficient of $x^{n-k}$ is
$\binom{n-k+\Ceiling{k/2}-1}{\Ceiling{k/2}-1}=\binom zw$.
Lucas's theorem modulo $2$ says $\binom zw$ is odd exactly when every $1$-bit of $w$
is also a $1$-bit of $z$, equivalently when $w\BitwiseAnd(z-w)=0$; here $z-w=n-k$.
Putting $k=n$ in \cref{eq:second-parity-bit} gives
$n\BitwiseAnd\Floor{(n-1)/2}=0$, which is equivalent to $n$ having one $1$-bit.
\end{proof}

\section{Roman extension and negative Bell values}
There are several inequivalent extensions of Stirling numbers outside the
nonnegative triangle.  The one used in the archive is the \emph{Roman} extension:
continue the two triangular recurrences wherever they determine a value and impose
the duality
\begin{equation}\label{eq:roman-duality}
 \UnsignedStirlingFirstKind nk=\StirlingSecondKind{-k}{-n},\qquad
 \StirlingSecondKind nk=\UnsignedStirlingFirstKind{-k}{-n}.
\end{equation}
These identities are definitions of a coherent reciprocal array, not assertions
about the original finite counting problems.

For $n\ge1$ and $k\ge0$, the mixed quadrant admits the finite expression
\begin{equation}\label{eq:negative-first-explicit}
 \UnsignedStirlingFirstKind{-n}{k}=\frac1{n!}\sum_{j=1}^n(-1)^{n-j}\binom nj\frac1{j^k}.
\end{equation}
For example,
\begin{equation}\label{eq:negative-first-five}
 \UnsignedStirlingFirstKind{-5}{k}=\frac1{120}
 \left(5-\frac{10}{2^k}+\frac{10}{3^k}-\frac5{4^k}+\frac1{5^k}\right).
\end{equation}
\begin{proof}
The ordinary generating function of the right-hand side in $k$ is a partial-fraction
expansion of
\[
 \frac{(-1)^{n-1}}{(1-x)(2-x)\cdots(n-x)}.
\]
Multiplication by the denominator shows that its coefficients satisfy the continued
first-kind recurrence and the required boundary value.  Uniqueness proves
\cref{eq:negative-first-explicit}; setting $n=5$ gives the example.
\end{proof}

Define negative Bell values by
\begin{equation}\label{eq:negative-bell-def}
 \BellNumber{-k}:=\sum_{n\ge1}\UnsignedStirlingFirstKind{-n}{k}\qquad(k\ge1).
\end{equation}
Then absolute convergence and \cref{eq:negative-first-explicit} give
\begin{equation}\label{eq:negative-dobinski}
 \BellNumber{-k}=\EulerE^{-1}\sum_{j\ge1}\frac1{j^k j!}.
\end{equation}
\begin{proof}
Interchange the absolutely convergent sums and write $n=j+r$:
\[
 \sum_{n\ge j}\frac{(-1)^{n-j}}{n!}\binom nj
 =\frac1{j!}\sum_{r\ge0}\frac{(-1)^r}{r!}=\frac{\EulerE^{-1}}{j!}.
\]
For $k=1$ this also yields
\[
 \BellNumber{-1}=\EulerE^{-1}\int_0^1\frac{\EulerE^t-1}{t}\Differential t
 =0.4848291\ldots,
\]
since $(\EulerE^t-1)/t=\sum_{j\ge1}t^{j-1}/j!$.
\end{proof}

\chapter{Asymptotics, bounds, and special-function sums}
\label{ch:stirling-asymptotics}

\section{Elementary asymptotic regimes}
\begin{theorem}[Fixed column and fixed diagonal]\label{thm:stirling-simple-asymp}
For fixed $k\ge1$,
\begin{equation}\label{eq:second-fixed-column-asymp}
 \StirlingSecondKind nk\sim\frac{k^n}{k!}\qquad(n\to\infty).
\end{equation}
For fixed $r\ge0$,
\begin{equation}\label{eq:both-near-diagonal-asymp}
 \UnsignedStirlingFirstKind{n+r}{n}\sim\StirlingSecondKind{n+r}{n}
 \sim\frac{n^{2r}}{2^rr!}\qquad(n\to\infty).
\end{equation}
Finally, uniformly for $k=o(\log n)$,
\begin{equation}\label{eq:first-small-column-asymp}
 \UnsignedStirlingFirstKind{n+1}{k+1}
 \sim\frac{n!}{k!}(\log n+\gamma)^k.
\end{equation}
\end{theorem}
\begin{proof}
In \cref{eq:second-explicit}, the term $j=k$ has size $k^n/k!$, and every other
term is exponentially smaller.  For the near diagonal, use
\cref{eq:first-elementary,eq:second-complete-symmetric}.  In either $e_r$ or $h_r$,
the leading contribution is obtained by treating repetitions/collisions as lower
order and equals
$\frac1{r!}(1+2+\cdots+n)^r\sim n^{2r}/(2^rr!)$; terms with a collision lose at least
one power of $n$.  For \cref{eq:first-small-column-asymp}, apply Cauchy's coefficient
formula to
\[
 \frac{\UnsignedStirlingFirstKind{n+1}{k+1}}{n!}=\CoefficientExtraction tk
 \exp\left(H_nt-\frac{H_n^{(2)}}2t^2+\frac{H_n^{(3)}}3t^3-\cdots\right).
\]
The saddle for the first exponential is $t=k/H_n=o(1)$.  Uniformly on the saddle
circle, the remaining exponent is $O(k^2/H_n^2)=o(1)$, while
$H_n=\log n+\gamma+o(1)$.  Hence the coefficient is asymptotic to $H_n^k/k!$.
\end{proof}

\section{Bounds for second-kind numbers}
For $n\ge2$ and $1\le k\le n-1$, the archive states
\begin{equation}\label{eq:second-bounds}
 \frac{k^2+k+2}{2}\,k^{n-k-1}-1
 \le\StirlingSecondKind nk
 \le\frac12\binom nk k^{n-k}.
\end{equation}
\begin{proof}
For the upper bound, in every partition choose the $k$ least elements of its blocks.
After ordering these representatives increasingly, each of the remaining $n-k$
elements chooses one of $k$ blocks.  This gives at most $\binom nk k^{n-k}$ encodings.
Except at $k=1$, each partition is encoded at least twice by reversing which of the
two largest block minima is declared first; the boundary case is immediate.  This
yields the factor $1/2$.

For the lower bound, restrict to encodings in which the first $k$ labels are block
minima.  There are $k^{n-k}$ assignments.  Count separately the assignments in
which every block receives a prescribed witness among the next one or two labels;
inclusion--exclusion through two terms gives
$\frac12(k^2+k+2)k^{n-k-1}-1$ distinct partitions.  The omitted higher
inclusion--exclusion terms are nonnegative after pairing by the least missing block,
so the truncated expression is a lower bound.  This is the standard elementary
proof of the displayed estimate.
\end{proof}


\section{Real zeros, log-concavity, and the location of the maximum}
For fixed $n$, package a row of the second-kind triangle into the Touchard
polynomial
\[
 \TouchardPolynomial{n}(x)=\sum_{k=0}^{n}\StirlingSecondKind nk x^k.
\]
The recurrence \eqref{eq:second-recurrence} is equivalent to
\begin{equation}
\label{eq:touchard-differential-recurrence-shape}
 \TouchardPolynomial{n+1}(x)=x\bigl(\TouchardPolynomial{n}(x)+\TouchardPolynomial{n}'(x)\bigr),
 \qquad \TouchardPolynomial{0}(x)=1.
\end{equation}

\begin{theorem}[Real-rootedness and strict log-concavity]
\label{thm:stirling-row-log-concavity}
For every $n\ge1$, all zeros of $\TouchardPolynomial{n}$ are real, simple, and nonpositive.  Hence
\begin{equation}
\label{eq:stirling-row-strict-log-concavity}
 \StirlingSecondKind nk^2>\StirlingSecondKind n{k-1}\StirlingSecondKind n{k+1}
 \qquad(1<k<n),
\end{equation}
and the row
$\StirlingSecondKind n1,\ldots,\StirlingSecondKind nn$ is unimodal.  Its maximum is attained at either
one index or two consecutive indices.
\end{theorem}
\begin{proof}
The assertion is clear for $\TouchardPolynomial{1}(x)=x$.  Suppose $f=\TouchardPolynomial{n}$ has simple zeros
\[
 r_1<r_2<\cdots<r_n=0.
\]
Away from these zeros,
\[
 \frac{f(x)+f'(x)}{f(x)}=1+\frac{f'(x)}{f(x)}
 =1+\sum_{i=1}^{n}\frac1{x-r_i}.
\]
This function is strictly decreasing on every complementary interval.  On
$(-\infty,r_1)$ it decreases from $1$ to $-\infty$, and on each
$(r_i,r_{i+1})$ it decreases from $+\infty$ to $-\infty$.  Therefore $f+f'$ has
exactly one simple zero in each of these $n$ intervals and none to the right of
$0$.  Formula \eqref{eq:touchard-differential-recurrence-shape} multiplies this
polynomial by $x$; since $(f+f')(0)=f'(0)=1$, it adds a new simple zero at $0$.
Induction proves real-rootedness and interlacing.

Newton's inequalities for a degree-$n$ polynomial with real nonpositive zeros say
that its normalized coefficients are log-concave:
\[
 \left(\frac{\StirlingSecondKind nk}{\binom nk}\right)^2
 \ge
 \frac{\StirlingSecondKind n{k-1}}{\binom n{k-1}}
 \frac{\StirlingSecondKind n{k+1}}{\binom n{k+1}}.
\]
The binomial-factor quotient is strictly larger than $1$, and all interior
coefficients are positive, giving \eqref{eq:stirling-row-strict-log-concavity}.
Consequently the successive ratios
$\StirlingSecondKind n{k+1}/\StirlingSecondKind nk$ strictly decrease.  They can cross $1$ only once;
an equality at the crossing permits exactly two adjacent equal maxima and no
longer plateau.
\end{proof}

Let $k_n$ denote the smaller index at which the maximum in row $n$ is attained.
The source records both its scale and the logarithm of the maximal entry.

\begin{theorem}[Asymptotic mode and maximal entry]
\label{thm:stirling-mode-asymptotic}
As $n\to\infty$,
\begin{align}
 k_n&\sim\frac{n}{\log n},
 \label{eq:stirling-mode-asymptotic}\\
 \log\StirlingSecondKind n{k_n}
 &=n\log n-n\log\log n-n
   +O\!\left(\frac{n\log\log n}{\log n}\right).
 \label{eq:stirling-maximum-log}
\end{align}
\end{theorem}
\begin{proof}
Two elementary estimates suffice.  Label the $k$ blocks of a partition to obtain
at most $k^n$ surjections, whence
\begin{equation}
\label{eq:stirling-mode-upper}
 \StirlingSecondKind nk\le\frac{k^n}{k!}.
\end{equation}
Conversely, require the first $k$ labels to lie in distinct blocks and assign each
of the remaining labels to one of them.  This injection gives
\begin{equation}
\label{eq:stirling-mode-lower}
 \StirlingSecondKind nk\ge k^{n-k}.
\end{equation}
Take $k_0=\Floor{n/\log n}$.  The lower estimate and
$\log k_0=\log n-\log\log n+O(\log n/n)$ give
\begin{equation}
\label{eq:stirling-mode-lower-log}
 \log\StirlingSecondKind n{k_0}
 \ge n\log n-n\log\log n-n
     +O\!\left(\frac{n\log\log n}{\log n}\right).
\end{equation}

For the upper estimate, Stirling's formula, uniformly for $1\le k\le n$, gives
\[
 \log\StirlingSecondKind nk
 \le n\log k-k\log k+k+O(\log n).
\]
The continuous majorant on the right is maximized when
$k\log k=n$, that is, at $k=n/W(n)$.  Since
\[
 W(n)=\log n-\log\log n
      +O\!\left(\frac{\log\log n}{\log n}\right),
\]
its maximum is the right side of \eqref{eq:stirling-maximum-log}.  Together with
\eqref{eq:stirling-mode-lower-log} this proves the maximal-value estimate.

It remains to localize the maximizing index.  Put
$k=c n/\log n$, with $c$ bounded away from $0$ and $\infty$.  Comparing the last
upper bound with \eqref{eq:stirling-mode-lower-log} gives
\[
 \log\StirlingSecondKind nk-\log\StirlingSecondKind n{k_0}
 \le n(\log c-c+1)+o(n).
\]
The function $\log c-c+1$ is strictly negative for $c\ne1$.  Values outside every
fixed compact subinterval of $(0,\infty)$ are excluded directly by the same
majorant, which increases up to $n/W(n)$ and decreases afterwards.  Therefore any
maximizing $k$ satisfies $k/(n/\log n)\to1$, proving
\eqref{eq:stirling-mode-asymptotic}.
\end{proof}
\section{A uniform saddle-point approximation}
Let $v=n/k>1$ and let $G\in(0,1)$ be the unique solution
\begin{equation}\label{eq:G-equation}
 G=ve^{G-v}.
\end{equation}
Equivalently, if $r=v-G$, then $r/(1-\EulerE^{-r})=v$.
Saddle-point extraction from $(\EulerE^z-1)^k/k!$ gives
\begin{equation}\label{eq:second-uniform-asymp}
 \StirlingSecondKind nk
 \sim
 \sqrt{\frac{v-1}{v(1-G)}}
 \left(\frac{v-1}{v-G}\right)^{n-k}
 \frac{k^n}{n^k}\EulerE^{k(1-G)}\binom nk,
\end{equation}
uniformly away from the two extreme edges; refined versions have relative error
$O(1/n)$, and the numerical constant $0.066/n$ quoted in the archive belongs to one
particular explicit error theorem.

\begin{proof}[Derivation]
By \cref{eq:second-egf},
\[
 \StirlingSecondKind nk=\frac{n!}{2\pi \ImaginaryUnit k!}\int
 \exp\bigl(k\log(\EulerE^z-1)-n\log z\bigr)\frac{\Differential z}{z}.
\]
The positive saddle $r$ solves
$ke^r/(\EulerE^r-1)=n/r$, or $r/(1-\EulerE^{-r})=v$.  Put $G=v-r$ to obtain
\cref{eq:G-equation}.  Expanding the phase to second order at $r$, evaluating the
Gaussian integral, and simplifying with Stirling's formula for $n!/k!$ gives exactly
\cref{eq:second-uniform-asymp}.  On a fixed angular complement of the saddle the
real part of the phase drops quadratically; Taylor's theorem with an explicit third-
and fourth-derivative bound gives a relative $O(1/n)$ remainder.  Thus the formula is
a theorem of the saddle method, not a heuristic replacement of the generating
function.
\end{proof}

\section{Zeta and Hurwitz-zeta sums}
\begin{theorem}[Stirling--zeta identities]\label{thm:stirling-zeta}
For $i\ge1$,
\begin{align}
 \sum_{n=i}^\infty\frac{\UnsignedStirlingFirstKind ni}{n\,n!}&=\RiemannZeta(i+1),
 \label{eq:first-zeta}\\
 \sum_{n=i}^\infty\frac{\UnsignedStirlingFirstKind ni}{n\,\RisingFactorial v n}&=\HurwitzZeta{i+1}{v},
 \qquad \Re v>0.\label{eq:first-hurwitz}
\end{align}
\end{theorem}
\begin{proof}
Integrate \cref{eq:unsigned-first-egf} after division by $x$:
\[
 \sum_{n\ge i}\frac{\UnsignedStirlingFirstKind ni}{n\,n!}
 =\frac1{i!}\int_0^1\frac{[-\log(1-x)]^i}{x}\Differential x.
\]
With $u=-\log(1-x)$ this becomes
$\frac1{i!}\int_0^\infty u^i/(\EulerE^u-1)\Differential u=\RiemannZeta(i+1)$.
For the Hurwitz form use
$1/\RisingFactorial v n=\Gamma(v)/\Gamma(v+n)$ and the beta integral
\[
 \frac1{\RisingFactorial v n}=\frac1{(n-1)!}\int_0^1t^{v-1}(1-t)^{n-1}\Differential t,
\]
then sum \cref{eq:unsigned-first-egf} under the integral and use the standard
Mellin integral for $\HurwitzZeta{i+1}{v}$.
\end{proof}

A repeated partial-fraction expansion of $x^{-k}$ around $1-x$ gives, for
$k\ge3$ and $\Re z>k-1$,
\begin{multline}\label{eq:log-power-integral}
 \int_0^1\frac{\log^z(1-x)}{x^k}\Differential x
 =\frac{(-1)^z\Gamma(z+1)}{(k-1)!}
 \sum_{r=1}^{k-1}\SignedStirlingFirstKind{k-1}{r}\\
 {}\times\sum_{m=0}^r\binom rm(k-2)^{r-m}\RiemannZeta(z+1-m),
\end{multline}
with the principal interpretation of $(-1)^z$ when $z$ is not integral.
\begin{proof}
Substitute $u=-\log(1-x)$, expand
$(1-\EulerE^{-u})^{-k}$ after $k-1$ integrations by parts, and express the resulting
polynomial in the Euler operator $uD$ in the ordinary-power basis by
$\SignedStirlingFirstKind{k-1}{r}$.  Each term is a Mellin integral
$\int_0^\infty u^{z-m}/(\EulerE^u-1)\Differential u=\Gamma(z+1-m)\RiemannZeta(z+1-m)$; simplifying gamma
ratios yields \cref{eq:log-power-integral}.
\end{proof}

The more elaborate Hurwitz expansion in the archive,
\begin{multline}\label{eq:hurwitz-finite-difference-expansion}
 \HurwitzZeta{s}{v}=\frac{k!}{\FallingFactorial{s-k}{k}}
 \sum_{n\ge0}\frac{\UnsignedStirlingFirstKind{n+k}{n}}{(n+k)!}\\
 {}\times\sum_{\ell=0}^{n+k-1}(-1)^\ell\binom{n+k-1}{\ell}
 (\ell+v)^{k-s},
\end{multline}
for $k\ge1$ (initially in the common domain of absolute convergence, then by
analytic continuation), follows by applying \cref{eq:der-from-diff} to the
$k$-fold antiderivative of $(x+v)^{-s}$ and summing over $x\in\NaturalNumbers$.  The inner sum is
a finite difference, while the outer coefficient is the first-kind conversion
coefficient.  This derivation also supplies the convergence conditions suppressed in
the compact archived display.


\section{Convergent Stirling series for \texorpdfstring{$\log\Gamma$}{log Gamma}, \texorpdfstring{$\psi$}{psi}, and the Stieltjes constants}
The source archive lists three nested Stirling series without their convergence
mechanism.  All three arise from the same substitution
\begin{equation}
\label{eq:logarithmic-binet-substitution}
 u=1-\EulerE^{-2\pi x},
 \qquad
 x=\frac{-\log(1-u)}{2\pi},
 \qquad
 \frac{\Differential x}{\EulerE^{2\pi x}-1}=\frac{\Differential u}{2\pi u}.
\end{equation}
We first record the coefficient estimate that justifies integration at the endpoint
$u=1$.

\begin{lemma}[Logarithmic transfer estimate]
\label{lem:logarithmic-transfer}
Suppose $h$ is analytic in the unit disk, continues to a dented neighborhood of
$u=1$, and, there,
\[
 h(u)=h_0+O\!\left(
 \frac{(1+\log\log\frac1{|1-u|})^M}
      {\log\frac1{|1-u|}}
 \right)
\]
for some fixed $M\ge0$.  If $h(u)=\sum_{n\ge0}h_nu^n$, then
\begin{equation}
\label{eq:logarithmic-transfer}
 h_n=O\!\left(
 \frac{(1+\log\log n)^M}{n\log^2 n}
 \right).
\end{equation}
Consequently $\sum_{n\ge1}|h_n|/n$ converges.
\end{lemma}
\begin{proof}
Apply Cauchy's coefficient formula on a keyhole contour that follows the two sides
of the cut $[1,\infty)$.  The part separated from $1$ is exponentially small.
On the local part put $u=1-w/n$.  The constant $h_0$ has no positive-index
coefficients, while the jump of
$1/\log(1/(1-u))$ across the cut is
$O(\log^{-2}n)$; multiplication by a power of $\log\log(1/(1-u))$ contributes at
most $(1+\log\log n)^M$.  The Cauchy kernel contributes $n^{-1}\EulerE^w$, and the two
rays of the keyhole have an integrable exponential majorant.  This gives
\eqref{eq:logarithmic-transfer}.  Division by another factor $n$ makes the
resulting majorant summable.
\end{proof}

\subsection{The logarithm and logarithmic derivative of the gamma function}
Binet's second formula, valid for $z>0$, is
\begin{equation}
\label{eq:binet-second}
 \log\Gamma(z)
 =\left(z-\frac12\right)\log z-z+\frac12\log(2\pi)
 +2\int_0^\infty\frac{\arctan(x/z)}{\EulerE^{2\pi x}-1}\,\Differential x.
\end{equation}
It follows, for example, by applying the Abel--Plana summation formula to
$w\mapsto\log(z+w)$, subtracting the elementary integral and endpoint terms, and
normalizing at $z=1$ by the Wallis product.  Differentiation under the integral
recovers the usual Binet integral for the digamma function, so the normalization
also proves \eqref{eq:binet-second} rather than merely determining it up to a
constant.

\begin{theorem}[Convergent Stirling series for $\log\Gamma$ and $\psi$]
\label{thm:gamma-stirling-series}
For real $z>1/6$,
\begin{multline}
\label{eq:gamma-stirling-series}
 \log\Gamma(z)
 =\left(z-\frac12\right)\log z-z+\frac12\log(2\pi)\\
 {}+\frac1\pi\sum_{n=1}^{\infty}\frac1{n\,n!}
 \sum_{\ell=0}^{\Floor{n/2}}
 \frac{(-1)^\ell(2\ell)!}{(2\pi z)^{2\ell+1}}
 \UnsignedStirlingFirstKind n{2\ell+1}.
\end{multline}
The series is locally uniformly convergent in $z>1/6$, and termwise
differentiation gives
\begin{multline}
\label{eq:digamma-stirling-series}
 \psi(z)=\frac{\Gamma'(z)}{\Gamma(z)}
 =\log z-\frac1{2z}\\
 {}-\frac1{\pi z}\sum_{n=1}^{\infty}\frac1{n\,n!}
 \sum_{\ell=0}^{\Floor{n/2}}
 \frac{(-1)^\ell(2\ell+1)!}{(2\pi z)^{2\ell+1}}
 \UnsignedStirlingFirstKind n{2\ell+1}.
\end{multline}
\end{theorem}
\begin{proof}
Put $L(u)=-\log(1-u)$.  The formal Taylor series of the composite function is
\begin{align}
 \arctan\frac{L(u)}{2\pi z}
 &=\sum_{\ell\ge0}\frac{(-1)^\ell}{2\ell+1}
   \frac{L(u)^{2\ell+1}}{(2\pi z)^{2\ell+1}}\notag\\
 &=\sum_{n\ge1}\frac{u^n}{n!}
   \sum_{\ell=0}^{\Floor{n/2}}
   \frac{(-1)^\ell(2\ell)!}{(2\pi z)^{2\ell+1}}
   \UnsignedStirlingFirstKind n{2\ell+1},
 \label{eq:arctan-log-stirling-expansion}
\end{align}
where the second line uses
$L(u)^r/r!=\sum_{n\ge r}\UnsignedStirlingFirstKind nr u^n/n!$.
For $z>1/6$, the analytic continuation of this composite from $u=0$ has no
singularity in $|u|<1$.  Indeed, the possible arctangent branch points would
require $\log(1-u)=\pm2\pi iz$; when $1/6<z<1/4$ their solutions have modulus
$2\sin(\pi z)>1$, and for $z\ge1/4$ the principal logarithm on the disk cannot
have imaginary part $\pm2\pi z$.  The only boundary singularity relevant to the
coefficient estimate is therefore $u=1$.

As $u\to1$ in a dented domain,
\[
 \arctan\frac{L(u)}{2\pi z}
 =\frac\pi2+O\!\left(\frac1{L(u)}\right).
\]
Thus \cref{lem:logarithmic-transfer} applies, uniformly when $z$ ranges over a
compact subinterval of $(1/6,\infty)$.  We may divide the Taylor coefficients by
$n$ and sum at $u=1$.  Substituting \eqref{eq:logarithmic-binet-substitution} in
\eqref{eq:binet-second} now gives
\[
 2\int_0^\infty\frac{\arctan(x/z)}{\EulerE^{2\pi x}-1}\Differential x
 =\frac1\pi\int_0^1\arctan\frac{L(u)}{2\pi z}\frac{\Differential u}{u},
\]
and termwise integration of
\eqref{eq:arctan-log-stirling-expansion} proves
\eqref{eq:gamma-stirling-series}.  The same uniform estimate after one
$z$-derivative justifies termwise differentiation; differentiating
$(2\pi z)^{-2\ell-1}$ produces the factor $-(2\ell+1)/z$ and yields
\eqref{eq:digamma-stirling-series}.
\end{proof}

For complex $z$ with $\Re z>0$, the same proof works whenever the two points
$1-\EulerE^{2\pi iz}$ and $1-\EulerE^{-2\pi iz}$ lie outside the closed unit disk; the stated
real half-line is the simplest connected domain containing all commonly used
positive arguments.

\subsection{Stieltjes constants}
Define the Stieltjes constants by the Laurent expansion
\begin{equation}
\label{eq:stieltjes-definition}
 \RiemannZeta(s)=\frac1{s-1}+\sum_{m=0}^{\infty}
 \frac{(-1)^m\gamma_m}{m!}(s-1)^m.
\end{equation}

\begin{theorem}[Stirling series for the Stieltjes constants]
\label{thm:stieltjes-stirling-series}
For every integer $m\ge0$,
\begin{equation}
\label{eq:stieltjes-stirling-series}
\begin{aligned}
 \gamma_m={}&\frac12\delta_{m0}
 +\frac{(-1)^m m!}{\pi}
 \sum_{n=1}^{\infty}\frac1{n\,n!}\\
 &\times\sum_{k=0}^{\Floor{n/2}}
 \frac{(-1)^k}{(2\pi)^{2k+1}}
 \UnsignedStirlingFirstKind{2k+2}{m+1}\UnsignedStirlingFirstKind n{2k+1}.
\end{aligned}
\end{equation}
The outer series converges absolutely.
\end{theorem}
\begin{proof}
The Abel--Plana formula applied to $(1+w)^{-s}$ gives Hermite's representation
\begin{equation}
\label{eq:hermite-zeta}
 \RiemannZeta(s)=\frac1{s-1}+\frac12+\frac1i\int_0^\infty
 \frac{(1-ix)^{-s}-(1+ix)^{-s}}{\EulerE^{2\pi x}-1}\,\Differential x,
 \qquad \Re s>0,
\end{equation}
with the pole omitted at $s=1$.  Expand the two powers at $s=1$ and compare with
\eqref{eq:stieltjes-definition}.  Differentiation under the exponentially
convergent integral yields
\begin{equation}
\label{eq:jensen-franel-stieltjes}
 \gamma_m=\frac12\delta_{m0}+\frac1i\int_0^\infty\frac1{\EulerE^{2\pi x}-1}
 \left\{
 \frac{\log^m(1-ix)}{1-ix}
 -\frac{\log^m(1+ix)}{1+ix}
 \right\}\Differential x.
\end{equation}

Differentiate the signed first-kind generating function to obtain, for $|w|<1$,
\begin{equation}
\label{eq:log-power-over-linear}
 \frac{\log^m(1+w)}{1+w}
 =m!\sum_{r\ge m}\SignedStirlingFirstKind{r+1}{m+1}\frac{w^r}{r!}.
\end{equation}
Only odd powers survive after replacing $w$ by $-ix$ and $ix$ and taking the
difference.  Since
$\SignedStirlingFirstKind{2k+2}{m+1}=(-1)^{m+1}\UnsignedStirlingFirstKind{2k+2}{m+1}$, one finds
\begin{multline}
\label{eq:stieltjes-odd-expansion}
 \frac1i\left\{
 \frac{\log^m(1-ix)}{1-ix}
 -\frac{\log^m(1+ix)}{1+ix}
 \right\}\\
 =2(-1)^m m!\sum_{k\ge0}
 \frac{(-1)^k\UnsignedStirlingFirstKind{2k+2}{m+1}}{(2k+1)!}\,x^{2k+1}.
\end{multline}
Now make the substitution \eqref{eq:logarithmic-binet-substitution}.  For every
$k$,
\[
 \frac{L(u)^{2k+1}}{(2k+1)!}
 =\sum_{n\ge2k+1}\UnsignedStirlingFirstKind n{2k+1}\frac{u^n}{n!}.
\]
The composite integrand in \eqref{eq:jensen-franel-stieltjes} is analytic in the
unit disk and, as $u\to1$, is
\[
 O\!\left(\frac{(1+\log L(u))^m}{L(u)}\right).
\]
Hence \cref{lem:logarithmic-transfer} permits absolute termwise integration after
division by $u$.  The Jacobian contributes $1/(2\pi)$, the factor $2$ in
\eqref{eq:stieltjes-odd-expansion} cancels it to $1/\pi$, and
$\int_0^1u^{n-1}\Differential u=1/n$.  Collecting the finite contributions to each
coefficient $u^n$ gives exactly \eqref{eq:stieltjes-stirling-series}.
\end{proof}

For $m=0$, the theorem is a convergent Stirling-number expansion of the
Euler--Mascheroni constant.  The factors
$\UnsignedStirlingFirstKind{2k+2}{m+1}$ show that, at every fixed outer index $n$, only finitely many
orders of logarithmic differentiation contribute; this is the same triangularity
that underlies Fa\`a di Bruno's formula and series reversion.
\section{Gregory coefficients and derivatives of logarithmic powers}
The Gregory coefficients $\GregoryCoefficient{n}$ are defined by
\[
 \frac{z}{\log(1+z)}=1+\sum_{n\ge1}\GregoryCoefficient{n}z^n.
\]
Multiplying by the first-kind expansion of powers of $\log(1+z)$ gives
\begin{equation}\label{eq:gregory-stirling}
 n!\GregoryCoefficient{n}=\sum_{\ell=0}^n\frac{\SignedStirlingFirstKind{n}{\ell}}{\ell+1}.
\end{equation}
Indeed,
$z/\log(1+z)=\int_0^1(1+z)^t\Differential t$ and comparison with
$(1+z)^t=\sum_n\sum_\ell \SignedStirlingFirstKind{n}{\ell}t^\ell z^n/n!$ proves the claim.

For arbitrary $\mu$ and $n\ge1$,
\begin{equation}\label{eq:log-power-derivative}
 \frac{\Differential^n}{\Differential x^n}(\log x)^\mu
 =x^{-n}\sum_{k=1}^n\FallingFactorial\mu k\,\SignedStirlingFirstKind{n}{n-k+1}(\log x)^{\mu-k}.
\end{equation}
\begin{proof}
Write $D_x=x^{-1}D_t$ with $t=\log x$.  Then
$D_x^n=x^{-n}\FallingFactorial{D_t}{n}$, as follows inductively from
$x^{-1}D_t\,x^{-m}=x^{-m-1}(D_t-m)$.  Expand
$\FallingFactorial{D_t}{n}$ by signed first-kind numbers and apply powers of $D_t$ to
$t^\mu$; reindexing gives the displayed form.
\end{proof}

\section{Generalized factorial coefficients}
Given $f:\NaturalNumbers\to\ComplexNumbers$ and $t\ne0$, define
\begin{equation}\label{eq:generalized-factorial-product}
 (x)_{n,f,t}:=\prod_{j=1}^{n-1}\left(x+\frac{f(j)}{t^j}\right),
 \qquad
 \UnsignedStirlingFirstKind nk_{f,t}:=\CoefficientExtraction{x}{k-1}(x)_{n,f,t}.
\end{equation}
Multiplication by the last factor immediately yields
\begin{equation}\label{eq:generalized-first-recurrence}
 \UnsignedStirlingFirstKind nk_{f,t}
 =f(n-1)t^{1-n}\UnsignedStirlingFirstKind{n-1}{k}_{f,t}
  +\UnsignedStirlingFirstKind{n-1}{k-1}_{f,t},
\end{equation}
with the usual delta boundary convention.  The generalized harmonic sums
$F_n^{(r)}(t)=\sum_{j\le n}t^j/f(j)^r$ enter the coefficients through the same
Newton/Bell-polynomial calculation as \cref{eq:first-harmonic-bell}.

More broadly, any triangular array satisfying
\begin{multline}\label{eq:triangular-general-recurrence}
 a(n,k)=(\alpha n+\beta k+\gamma)a(n-1,k)\\
 +(\alpha'n+\beta'k+\gamma')a(n-1,k-1)+\delta_{n0}\delta_{k0}
\end{multline}
is governed by a first-order differential equation for its bivariate generating
function.  This umbrella contains both Stirling kinds, Lah numbers, Eulerian
variants, and many generalized factorial coefficients.



\chapter{Restricted variants of set partitions}
\label{ch:stirling-variants}
The second-kind recurrence survives many natural restrictions, but its boundary
term changes according to how the block containing the newest label is allowed to
look.

\section{\texorpdfstring{$r$}{r}-Stirling numbers of the second kind}
For fixed $r\ge0$, let
\[
 \StirlingSecondKind nk_{\!r}
\]
denote the number of partitions of $[n]$ into $k$ nonempty blocks in which the
labels $1,\ldots,r$ lie in distinct blocks.  The definition is meaningful for
$n\ge r$ and forces $k\ge r$.

\begin{theorem}[$r$-Stirling recurrence]
\label{thm:r-stirling-recurrence}
For $n>r$,
\begin{equation}
\label{eq:r-stirling-recurrence}
 \StirlingSecondKind nk_{\!r}
 =k\StirlingSecondKind{n-1}{k}_{\!r}+\StirlingSecondKind{n-1}{k-1}_{\!r},
\end{equation}
with $\StirlingSecondKind rr_{\!r}=1$ and the natural zero boundary conditions.
\end{theorem}
\begin{proof}
Because $n>r$, the new label $n$ is not one of the distinguished labels.  It either
joins one of the $k$ blocks of an admissible partition of $[n-1]$, or forms a
singleton beside an admissible partition into $k-1$ blocks.  The distinction of
$1,\ldots,r$ is preserved in both constructions, and deletion of $n$ reverses
them.
\end{proof}

The same labelled-component argument gives the mixed exponential generating
function
\begin{equation}
\label{eq:r-stirling-egf}
 \sum_{n\ge r}\sum_{k\ge r}\StirlingSecondKind nk_{\!r}
 y^k\frac{z^{n-r}}{(n-r)!}
 =y^re^{rz}\exp\!\bigl(y(\EulerE^z-1)\bigr).
\end{equation}
Indeed, the $r$ distinguished blocks each consist of one distinguished atom and an
arbitrary set of ordinary atoms, contributing $y^re^{rz}$, while the remaining
blocks form an ordinary set partition.

\section{Associated Stirling numbers}
Let $\mathsf S_r(n,k)$ count partitions of $[n]$ into $k$ blocks, every block having
size at least $r$.

\begin{theorem}[Associated recurrence]
\label{thm:associated-stirling-recurrence}
For $r\ge1$,
\begin{equation}
\label{eq:associated-stirling-recurrence}
 \mathsf S_r(n+1,k)
 =k\mathsf S_r(n,k)+\binom n{r-1}\mathsf S_r(n-r+1,k-1).
\end{equation}
Moreover,
\begin{equation}
\label{eq:associated-stirling-egf}
 \sum_{n\ge0}\mathsf S_r(n,k)\frac{z^n}{n!}
 =\frac1{k!}\left(\EulerE^z-\sum_{j=0}^{r-1}\frac{z^j}{j!}\right)^k.
\end{equation}
\end{theorem}
\begin{proof}
Consider the block containing $n+1$.  If deleting $n+1$ leaves a block of size at
least $r$, start from a partition counted by $\mathsf S_r(n,k)$ and choose one of
its $k$ blocks for insertion.  Otherwise the block containing $n+1$ has exactly
$r$ elements: choose its $r-1$ companions and partition the remaining $n-r+1$
labels into $k-1$ admissible blocks.  These cases are disjoint and exhaustive,
proving the recurrence.  For the generating function, one admissible block is a
labelled set of size at least $r$, with EGF
$\EulerE^z-\sum_{j<r}z^j/j!$; an unordered set of $k$ such blocks contributes its $k$th
power divided by $k!$.
\end{proof}

For $r=2$ these are often called Ward numbers.  Their occurrence as coefficients
of other polynomial families comes from the same exclusion of singleton
components expressed by \eqref{eq:associated-stirling-egf}.

\section{Reduced Stirling numbers and graph coloring}
For $d\ge1$, let $\mathsf S^{[d]}(n,k)$ count partitions of $[n]$ into $k$ blocks
such that two labels in the same block always differ by at least $d$.

\begin{theorem}[Reduction formula]
\label{thm:reduced-stirling}
For $n\ge k\ge d$,
\begin{equation}
\label{eq:reduced-stirling}
 \mathsf S^{[d]}(n,k)=\StirlingSecondKind{n-d+1}{k-d+1}.
\end{equation}
In particular $\mathsf S^{[1]}(n,k)=\StirlingSecondKind nk$.
\end{theorem}
\begin{proof}
Join two vertices $i,j\in[n]$ when $0<|i-j|<d$, and call the resulting graph
$G_{n,d}$.  Its independent-set partitions are exactly the partitions counted by
$\mathsf S^{[d]}(n,k)$.  A proper coloring can therefore be obtained by choosing
such a partition and assigning distinct colors to its blocks, so
\begin{equation}
\label{eq:reduced-chromatic-falling}
 \chi_{G_{n,d}}(q)=\sum_k\mathsf S^{[d]}(n,k)\FallingFactorial qk.
\end{equation}
Color the vertices in their natural order.  The first $d-1$ vertices must receive
distinct colors; every subsequent vertex sees the preceding $d-1$ vertices, which
form a clique, and hence has $q-d+1$ choices.  Thus
\[
 \chi_{G_{n,d}}(q)
 =\FallingFactorial q{d-1}(q-d+1)^{n-d+1}.
\]
Apply the ordinary Stirling basis expansion to the last power:
\begin{align*}
 \chi_{G_{n,d}}(q)
 &=\FallingFactorial q{d-1}
   \sum_{r\ge0}\StirlingSecondKind{n-d+1}{r}\FallingFactorial{q-d+1}{r}\\
 &=\sum_{r\ge0}\StirlingSecondKind{n-d+1}{r}\FallingFactorial q{r+d-1}.
\end{align*}
Comparison with \eqref{eq:reduced-chromatic-falling} gives
$k=r+d-1$ and proves \eqref{eq:reduced-stirling}.
\end{proof}
\part{Bell Numbers and Set Partitions}

\chapter{Set partitions and the many meanings of the Bell numbers}
\index{Bell numbers}

\section{Definition, initial values, and equivalence relations}
\begin{definition}
For $n\ge0$, the \emph{$n$th Bell number} $\BellNumber n$ is the number of partitions of
an $n$-element labelled set into nonempty, pairwise disjoint blocks.  Thus
\[
 \BellNumber0=\BellNumber1=1,
 \qquad
 (\BellNumber n)_{n\ge0}=1,1,2,5,15,52,203,877,4140,\ldots .
\]
The empty set has one partition, namely the empty family of blocks.
\end{definition}

For example, the five partitions of $\{a,b,c\}$ are
\[
 abc,\qquad a\mid bc,\qquad b\mid ac,\qquad c\mid ab,\qquad a\mid b\mid c,
\]
where vertical bars separate blocks.

\begin{proposition}
Partitions of a set $S$ are in natural bijection with equivalence relations on $S$.
Consequently $\BellNumber n$ also counts equivalence relations on an $n$-element set.
\end{proposition}
\begin{proof}
Given a partition $\pi$, declare $x\sim_\pi y$ when $x$ and $y$ lie in the same
block.  Reflexivity, symmetry, and transitivity follow immediately.  Conversely,
the equivalence classes of an equivalence relation are nonempty, disjoint, cover
$S$, and hence form a partition.  The two constructions undo one another.
\end{proof}


\section{Ordered set partitions and the Fubini numbers}
If the blocks themselves are linearly ordered, the resulting numbers are usually
called the \emph{ordered Bell numbers} or \emph{Fubini numbers}.  To prevent a
collision with the Bell notation, write them as $\mathsf F_n$.

\begin{theorem}[Ordered Bell numbers]
\label{thm:ordered-bell}
For $n\ge0$,
\begin{align}
 \mathsf F_n&=\sum_{k=0}^{n}k!\StirlingSecondKind nk,
 \label{eq:ordered-bell-stirling}\\
 \sum_{n\ge0}\mathsf F_n\frac{z^n}{n!}
 &=\frac{1}{2-\EulerE^z},
 \label{eq:ordered-bell-egf}\\
 \mathsf F_n&=\sum_{j=1}^{n}\binom nj\mathsf F_{n-j}
 \qquad(n\ge1),
 \label{eq:ordered-bell-recurrence}
\end{align}
with $\mathsf F_0=1$.
\end{theorem}
\begin{proof}
A partition into $k$ blocks may be ordered in $k!$ ways, proving
\eqref{eq:ordered-bell-stirling}.  By the second-kind exponential generating
function,
\[
 \sum_{n\ge0}\mathsf F_n\frac{z^n}{n!}
 =\sum_{k\ge0}k!\frac{(\EulerE^z-1)^k}{k!}
 =\sum_{k\ge0}(\EulerE^z-1)^k
 =\frac1{2-\EulerE^z},
\]
first as a formal identity and then analytically near the origin.  Finally, in a
nonempty ordered partition choose the elements of its first block, which may be any
nonempty subset, and order-partition the complement.  Summing over the first-block
size gives \eqref{eq:ordered-bell-recurrence}.
\end{proof}

Thus
\[
 (\mathsf F_n)_{n\ge0}=1,1,3,13,75,541,4683,\ldots,
\]
which contrasts with the unordered Bell sequence by recording the order of the
blocks rather than only their underlying family.
\section{Squarefree factorizations, rhyme schemes, and restricted-growth words}
\begin{proposition}
If $N=p_1\cdots p_n$ is squarefree, then $\BellNumber n$ is the number of unordered
factorizations of $N$ into integers greater than $1$.
\end{proposition}
\begin{proof}
To a set partition $\pi$ of $[n]$ associate the multiset of factors
$\{\prod_{i\in B}p_i:B\in\pi\}$.  Unique factorization recovers each block from
its factor, so this is a bijection.  For $30=2\cdot3\cdot5$ the five
factorizations are
\[
 30,\quad 2\cdot15,\quad3\cdot10,\quad5\cdot6,\quad2\cdot3\cdot5.
\]
\end{proof}

A rhyme scheme of an $n$-line stanza is likewise a partition of the set of lines:
two indices lie in the same block exactly when the corresponding lines rhyme.
The customary lettering convention is the following canonical encoding.

\begin{definition}
A \emph{restricted-growth word} of length $n$ is a word
$a_1\cdots a_n$ of nonnegative integers such that $a_1=0$ and
\[
 a_i\le 1+\max(a_1,\ldots,a_{i-1})\qquad(i>1).
\]
\end{definition}

\begin{proposition}
Set partitions of $[n]$, restricted-growth words of length $n$, and canonical
rhyme schemes of length $n$ are mutually bijective.
\end{proposition}
\begin{proof}
Order the blocks by increasing least element, number them $0,1,\ldots$, and put
$a_i=j$ when $i$ belongs to block $j$.  The displayed restriction says precisely
that a new block receives the next unused number.  Conversely, equal letters form
the blocks.  Replacing $0,1,2,\ldots$ by $A,B,C,\ldots$ gives the rhyme scheme.
\end{proof}

\section{A card-shuffling bijection}
Consider the operation that removes the top card from a deck of $n$ labelled cards
and reinserts it in any of the $n$ positions.  There are $n^n$ sequences of $n$
such operations.

\begin{theorem}[Return-to-order shuffles]
Exactly $\BellNumber n$ of those $n^n$ operation sequences return an initially ordered
deck to its original order.  Hence the return probability is $\BellNumber n/n^n$.
\end{theorem}
\begin{proof}
Reverse time.  A reversed operation chooses an arbitrary card and moves it to the
top.  Thus a reversed history is a word $c_1\cdots c_n$ in the card labels.  After
all moves, the cards that occurred in the word appear first, in decreasing order of
their last occurrence times; cards never chosen follow in their original relative
order.

Partition the time set $[n]$ according to equality of the letters $c_i$.  For the
final deck to be $1,2,\ldots,n$, the labels that occur must be an initial interval
$1,\ldots,m$, and their blocks of occurrence times must be labelled in decreasing
order of their maxima: the block with largest maximum receives label $1$, the next
receives $2$, and so on.  Therefore every set partition of $[n]$ determines exactly
one successful word, and every successful word determines its equality partition.
This is the required bijection.
\end{proof}

\section{Bell-enumerated vincular pattern classes}
A dash in a permutation pattern means that the corresponding entries need not be
adjacent; an undashed pair must be adjacent.  For instance, an occurrence of
$23\!\mathord{-}\!1$ consists of an adjacent ascent followed later by a smaller
entry.

\begin{theorem}
The permutations of $[n]$ avoiding $23\!\mathord{-}\!1$ are counted by $\BellNumber n$.
The same is true for each pattern in
\[
 1\!\mathord{-}\!23,\ 12\!\mathord{-}\!3,\ 32\!\mathord{-}\!1,
 \ 3\!\mathord{-}\!21,\ 1\!\mathord{-}\!32,\ 3\!\mathord{-}\!12,
 \ 21\!\mathord{-}\!3,
\]
and for the equivalent barred-pattern class in which every $321$ occurrence
extends to a $3241$ occurrence.
\end{theorem}
\begin{proof}
For the pattern $23\!\mathord{-}\!1$, order the blocks of a partition by increasing
least element and write the entries of each block in decreasing order.  Concatenate
the resulting words.  In the resulting permutation, each block is a decreasing
run, and every ascent occurs between two blocks.  Its lower endpoint is a
right-to-left minimum, so no later entry is smaller; hence $23\!\mathord{-}\!1$ is
avoided.

Conversely, cut a $23\!\mathord{-}\!1$-avoiding permutation immediately after
each right-to-left minimum.  Every resulting segment is decreasing: otherwise an
ascent inside a segment, together with its terminal right-to-left minimum, would
form $23\!\mathord{-}\!1$.  The segment minima increase from left to right, so the
segments, viewed as sets, are blocks already in canonical order.  The two maps are
inverse.

Reversal, complement, inverse, and their compositions are bijections on
permutations and carry vincular occurrences to occurrences of the seven listed
patterns; therefore all eight classes are equinumerous.  Finally, cutting at
right-to-left minima shows that failure of the decreasing-run condition is
simultaneously an occurrence of $23\!\mathord{-}\!1$ and a $321$ that cannot be
extended by an intervening entry to $3241$.  Thus the barred-pattern formulation
describes the same class.
\end{proof}

\begin{remark}
Classical avoidance classes have at most exponential growth in $n$, whereas Bell
numbers grow faster than $C^n$ for every fixed $C$.  The adjacency requirement is
therefore essential: these Bell-enumerated classes cannot be classical
single-pattern avoidance classes.
\end{remark}

\section{The Bell--Aitken triangle}
Define a triangular array $(a_{n,k})_{0\le k\le n}$ by
\[
 a_{0,0}=1,\qquad a_{n,0}=a_{n-1,n-1},\qquad
 a_{n,k}=a_{n,k-1}+a_{n-1,k-1}\quad(1\le k\le n).
\]
Its first rows are
\[
\begin{array}{rrrrr}
1\\
1&2\\
2&3&5\\
5&7&10&15\\
15&20&27&37&52.
\end{array}
\]

\begin{theorem}
For $0\le k\le n$,
\begin{equation}\label{eq:bell-triangle-closed}
 a_{n,k}=\sum_{j=0}^{k}\binom{k}{j}\BellNumber{n-j}.
\end{equation}
Consequently the left edge is $a_{n,0}=\BellNumber n$ and the right edge is
$a_{n,n}=\BellNumber{n+1}$.
\end{theorem}
\begin{proof}
The formula is immediate for $k=0$.  Pascal's identity gives
\[
 a_{n,k-1}+a_{n-1,k-1}
 =\sum_j\left(\binom{k-1}{j}+\binom{k-1}{j-1}\right)\BellNumber{n-j},
\]
which is \eqref{eq:bell-triangle-closed}.  For $k=n$, reverse the index and use
the Bell recurrence proved in \cref{thm:bell-binomial-recurrence} below.
\end{proof}

\chapter{The enumerative calculus of Bell numbers}

\section{Two fundamental sums}
\begin{theorem}[Binomial recurrence]\label{thm:bell-binomial-recurrence}
For $n\ge0$,
\begin{equation}\label{eq:bell-binomial-recurrence}
 \BellNumber{n+1}=\sum_{k=0}^{n}\binom nk\BellNumber k.
\end{equation}
\end{theorem}
\begin{proof}
In a partition of $\{0,1,\ldots,n\}$, let $k$ be the number of elements outside
the block containing $0$.  Choose those $k$ elements and partition them
arbitrarily; the block containing $0$ is then forced.
\end{proof}

\begin{theorem}[Stirling decomposition]
\begin{equation}\label{eq:bell-stirling-sum}
 \BellNumber n=\sum_{k=0}^{n}\StirlingSecondKind nk
 =\sum_{k=0}^{n}\frac1{k!}\sum_{i=0}^{k}(-1)^{k-i}\binom ki i^n.
\end{equation}
\end{theorem}
\begin{proof}
The first equality groups partitions by their number of blocks.  Substitute the
inclusion--exclusion formula for $\StirlingSecondKind nk$ from
\cref{eq:second-explicit} to obtain the second equality.
\end{proof}

\section{Spivey's identity and binomial inversions}
\begin{theorem}[Spivey's identity]\label{thm:spivey}
For $m,n\ge0$,
\begin{equation}\label{eq:spivey}
 \BellNumber{m+n}=\sum_{k=0}^{n}\sum_{j=0}^{m}
 \StirlingSecondKind mj\binom nk j^{\,n-k}\BellNumber k.
\end{equation}
\end{theorem}
\begin{proof}
Partition the first $m$ labels into $j$ blocks.  Of the final $n$ labels, choose
$k$ whose blocks contain no old label and partition them in $\BellNumber k$ ways.  Each
of the remaining $n-k$ labels independently joins one of the $j$ old blocks.  This
classification is unique and gives the summand.
\end{proof}

Binomial inversion of \eqref{eq:bell-binomial-recurrence} gives the first identity
below.  The second is a useful shifted form.
\begin{theorem}[Bell inversion identities]\label{thm:bell-inversions}
For $n,k\ge0$,
\begin{align}
 \BellNumber n&=\sum_{j=0}^{n}(-1)^{n-j}\binom nj\BellNumber{j+1},
 \label{eq:bell-inversion-one}\\
 \sum_{j=0}^{n}\binom nj\BellNumber{k+j}
 &=\sum_{i=0}^{k}(-1)^{k-i}\binom ki\BellNumber{n+i+1}.
 \label{eq:bell-inversion-two}
\end{align}
\end{theorem}
\begin{proof}
The first formula is ordinary binomial inversion.  For the second let $X$ be a
Poisson random variable of mean $1$; \cref{thm:dobinski} below gives
$\Expectation X^r=\BellNumber r$.  The left side is
$\Expectation[X^k(X+1)^n]$.  Poisson summation by parts,
\begin{equation}\label{eq:poisson-sbp}
 \Expectation[Xh(X)]=\Expectation[h(X+1)],
\end{equation}
iterated $k$ times in the equivalent form
$\Expectation[X^{n+1}(X-1)^k]=\Expectation[(X+1)^nX^k]$, turns it into the right side after expanding
$(X-1)^k$.
\end{proof}

The first-kind Stirling numbers package all repeated shifts at once.
\begin{theorem}[Weighted Bell shift]\label{thm:weighted-bell-shift}
For $a,b\in\ComplexNumbers$ and $k,n\ge0$,
\begin{multline}\label{eq:weighted-bell-shift}
 \sum_{j=0}^{n}\binom nj a^j b^{n-j}\BellNumber j\\
 =\sum_{i=0}^{k}(-1)^{k-i}\UnsignedStirlingFirstKind ki
   \sum_{j=0}^{n}\binom nj a^j(b-ak)^{n-j}\BellNumber{j+i}.
\end{multline}
In particular,
\begin{align}
 \sum_{j=0}^{n}\binom nj a^j b^{n-j}\BellNumber j
 &=\sum_{j=0}^{n}\binom nj a^j(b-a)^{n-j}\BellNumber{j+1},
 \label{eq:weighted-bell-k1}\\
 \sum_{j=0}^{n}\binom nj k^{n-j}\BellNumber j
 &=\sum_{i=0}^{k}(-1)^{k-i}\UnsignedStirlingFirstKind ki\BellNumber{n+i}.
 \label{eq:weighted-bell-special}
\end{align}
\end{theorem}
\begin{proof}
By Dobiński's formula, the left side of
\eqref{eq:weighted-bell-shift} is
\[
 \EulerE^{-1}\sum_{m\ge0}\frac{(am+b)^n}{m!}.
\]
On the right, use
$\sum_i(-1)^{k-i}\UnsignedStirlingFirstKind ki m^i=\FallingFactorial mk$ inside the same Poisson sum.  It
becomes
\[
 \EulerE^{-1}\sum_{m\ge0}\frac{\FallingFactorial mk\,[a m+b-ak]^n}{m!}
 =\EulerE^{-1}\sum_{r\ge0}\frac{(ar+b)^n}{r!},
\]
after $m=r+k$.  The special cases follow by setting $k=1$, and then $a=1,b=k$.
\end{proof}

\section{Exponential generating function and differential equation}
\begin{theorem}\label{thm:bell-egf}
The Bell exponential generating function is
\begin{equation}\label{eq:bell-egf}
 \mathscr B(z):=\sum_{n\ge0}\BellNumber n\frac{z^n}{n!}
 =\exp(\EulerE^z-1),
 \qquad
 \mathscr B'(z)=\EulerE^z\mathscr B(z).
\end{equation}
\end{theorem}
\begin{proof}
A set partition is a labelled set of nonempty labelled sets.  The EGF of a
nonempty set is $\EulerE^z-1$, and the exponential formula therefore gives
$\exp(\EulerE^z-1)$.  Differentiation gives the differential equation; conversely,
coefficient comparison in that equation gives
\eqref{eq:bell-binomial-recurrence} and $\mathscr B(0)=1$, so it uniquely determines
the series.
\end{proof}

\section{Dobiński's formula and Poisson moments}
\begin{theorem}[Dobiński]\label{thm:dobinski}
For every $n\ge0$,
\begin{equation}\label{eq:dobinski}
 \BellNumber n=\EulerE^{-1}\sum_{m=0}^{\infty}\frac{m^n}{m!}.
\end{equation}
Equivalently, if $X\sim\operatorname{Poisson}(1)$, then $\BellNumber n=\Expectation X^n$.
\end{theorem}
\begin{proof}
Expand the outer exponential in \eqref{eq:bell-egf}:
\[
 \EulerE^{\EulerE^z-1}=\EulerE^{-1}\sum_{m\ge0}\frac{\EulerE^{mz}}{m!}
 =\EulerE^{-1}\sum_{m\ge0}\frac1{m!}\sum_{n\ge0}m^n\frac{z^n}{n!}.
\]
Absolute convergence on compact sets justifies rearrangement analytically, while
formal coefficient extraction gives the same identity algebraically.  The last
statement is the defining moment sum of a Poisson$(1)$ variable.
\end{proof}

\begin{corollary}[Cauchy integral]
For every positively oriented circle $\gamma$ around $0$,
\begin{equation}\label{eq:bell-cauchy}
 \BellNumber n=\frac{n!}{2\pi \ImaginaryUnit \EulerE}\int_\gamma
       \frac{\EulerE^{\EulerE^z}}{z^{n+1}}\,\Differential z.
\end{equation}
\end{corollary}
\begin{proof}
Apply Cauchy's coefficient formula to $\EulerE^{-1}\EulerE^{\EulerE^z}$.
\end{proof}

\chapter{Congruences, shape, and growth of Bell numbers}

\section{Touchard congruences}
It is useful to refine $\BellNumber n$ to the Touchard polynomial
\[
 \TouchardPolynomial{n}(x)=\sum_{k=0}^{n}\StirlingSecondKind nk x^k,
 \qquad \TouchardPolynomial{n}(1)=\BellNumber n.
\]

\begin{theorem}[Polynomial Touchard congruence]\label{thm:touchard-poly}
For a prime $p$,
\begin{equation}\label{eq:touchard-poly}
 \TouchardPolynomial{n+p}(x)\equiv \TouchardPolynomial{n+1}(x)+x^p\TouchardPolynomial{n}(x)\pmod p.
\end{equation}
Consequently
\begin{equation}\label{eq:touchard}
 \BellNumber{n+p}\equiv\BellNumber n+\BellNumber{n+1}\pmod p.
\end{equation}
\end{theorem}
\begin{proof}
Let the cyclic group $C_p$ rotate a distinguished set of $p$ labels while fixing
$n$ other labels, and let it act on partitions, preserving the weight $x^{\#\text{
blocks}}$.  Every nonfixed orbit has size $p$ and contributes $0$ modulo $p$.
In a fixed partition there are two possibilities.  Either all $p$ rotating labels
form $p$ singleton blocks, independently of a partition of the other labels; this
contributes $x^p\TouchardPolynomial{n}(x)$.  Otherwise rotational invariance forces the rotating labels
to lie in one common block, possibly together with old labels.  Contracting those
$p$ labels to one distinguished label gives a partition of $n+1$ labels with the
same number of blocks, contributing $\TouchardPolynomial{n+1}(x)$.  This proves the polynomial
congruence; set $x=1$.
\end{proof}

\begin{theorem}[Prime-power shift]\label{thm:bell-prime-power-shift}
For a prime $p$ and $m\ge1$,
\begin{equation}\label{eq:bell-prime-power-shift}
 \BellNumber{n+p^m}\equiv m\BellNumber n+\BellNumber{n+1}\pmod p.
\end{equation}
More precisely,
\begin{equation}\label{eq:touchard-prime-power-poly}
 \TouchardPolynomial{n+p^m}(x)\equiv \TouchardPolynomial{n+1}(x)+\TouchardPolynomial{n}(x)\sum_{r=1}^{m}x^{p^r}
 \pmod p.
\end{equation}
\end{theorem}
\begin{proof}
Let the additive cyclic group of order $p^m$ act regularly on $p^m$ new labels.
A partition fixed by the whole group induces a block system for this regular
action.  Apart from the case in which all new labels join one invariant block
(which contracts to one label and contributes $\TouchardPolynomial{n+1}$), the blocks on the new
orbit are the cosets of the unique subgroup of index $p^r$ for some
$1\le r\le m$; they contribute $p^r$ blocks and cannot mix with old labels.
Thus the fixed weighted partitions contribute
$\TouchardPolynomial{n+1}+\TouchardPolynomial{n}\sum_{r=1}^m x^{p^r}$.  Every other orbit has cardinality divisible by
$p$, proving \eqref{eq:touchard-prime-power-poly}.  Fermat's congruence
$x^{p^r}\equiv x$ modulo $p$ at $x=1$ gives
\eqref{eq:bell-prime-power-shift}.
\end{proof}

\section{A universal period bound modulo a prime}
\begin{theorem}\label{thm:bell-period-bound}
Modulo a prime $p$, the Bell sequence is periodic, and its period divides
\begin{equation}\label{eq:bell-period-N}
 N_p=\frac{p^p-1}{p-1}=1+p+\cdots+p^{p-1}.
\end{equation}
\end{theorem}
\begin{proof}
Touchard's recurrence says that the shift operator $E$ on the sequence satisfies
$E^p-E-1=0$ over $\FiniteField_p$.  The polynomial $q(t)=t^p-t-1$ is separable because
$q'(t)=-1$.  If $\alpha$ is any root in a splitting field, then
$\alpha^p=\alpha+1$ and hence $\alpha^{p^j}=\alpha+j$ for $j\in\FiniteField_p$.  Taking the
product over $j$ gives
\[
 \alpha^{N_p}=\prod_{j\in\FiniteField_p}\alpha^{p^j}
 =\prod_{j\in\FiniteField_p}(\alpha+j)=\alpha^p-\alpha=1.
\]
Thus every eigenvalue of $E$ is an $N_p$th root of unity.  Separability makes $E$
diagonalizable over the splitting field, so $E^{N_p}=1$ on every solution of the
recurrence, in particular on the Bell sequence.
\end{proof}

\begin{remark}
A reported exact period is computational data, not a consequence of the divisibility
bound alone.  It can be certified by generating $N_p$ terms with
\eqref{eq:touchard}, finding the least divisor $d\mid N_p$ for which the state vector
of $p$ consecutive terms repeats, and checking that no proper divisor of $d$ works.
The same recurrence permits a compact machine-checkable certificate.
\end{remark}

\section{Log-convexity and normalized log-concavity}
\begin{theorem}\label{thm:bell-log-shape}
For $n\ge1$,
\begin{equation}\label{eq:bell-log-shape}
 1\le \frac{\BellNumber{n-1}\BellNumber{n+1}}{\BellNumber n^2}\le\frac{n+1}{n}.
\end{equation}
Thus $(\BellNumber n)$ is log-convex, whereas $(\BellNumber n/n!)$ is log-concave.
\end{theorem}
\begin{proof}
For the left inequality, let $X\sim\operatorname{Poisson}(1)$.  Cauchy--Schwarz
applied to $X^{(n-1)/2}$ and $X^{(n+1)/2}$ gives
$(\Expectation X^n)^2\le \Expectation X^{n-1}\Expectation X^{n+1}$.

For completeness we prove the normalized inequality through the following
convolution lemma.  If nonnegative $(z_j)_{j\ge1}$ is log-concave without internal
zeros and
\[
 A(t)=\exp\!\left(\sum_{j\ge1}\frac{z_j}{j}t^j\right)
      =\sum_{n\ge0}a_n\frac{t^n}{n!},
\]
then $a_n/n!$ is log-concave.  To see this, truncate at degree $N$, differentiate,
and compare the adjacent $2\times2$ Toeplitz minors in
$A'=A\sum_{j\ge1}z_jt^{j-1}$.  After cross multiplication, the $n$th minor is
\[
 \sum_{0\le i<j\le n}
 \frac{a_i a_j}{i!j!}
 \bigl(z_{n-i+1}z_{n-j}-z_{n-i}z_{n-j+1}\bigr),
\]
with the boundary terms interpreted as zero.  Each bracket is nonnegative by the
monotonicity of the ratios $z_{r+1}/z_r$; induction on $n$ starts at $a_0=1$.
Removing the truncation proves the lemma coefficientwise.

For $A(t)=\EulerE^{\EulerE^t-1}$ we have
$\sum_{j\ge1}t^j/j!=\sum_{j\ge1}z_jt^j/j$ with
$z_j=1/(j-1)!$.  This sequence is log-concave, since
$z_j^2/(z_{j-1}z_{j+1})=j/(j-1)$.  Hence
$(\BellNumber n/n!)^2\ge(\BellNumber{n-1}/(n-1)!)(\BellNumber{n+1}/(n+1)!)$, which is the right
inequality in \eqref{eq:bell-log-shape}.
\end{proof}

\section{Explicit upper bounds}
The positive-coefficient Cauchy estimate supplies a short route to the numerical
bounds in the archive.
\begin{lemma}\label{lem:bell-coeff-upper}
For $t>0$,
\begin{equation}\label{eq:bell-coeff-upper}
 \BellNumber n\le n!\,t^{-n}\exp(\EulerE^t-1).
\end{equation}
For $n\ge2$, put $L=\log(n+1)$ and $t=\log(n/L)$.  Then
\begin{equation}\label{eq:bell-ratio-upper}
 \BellNumber n^{1/n}\frac{L}{n}
 \le \frac{L}{t}
 \exp\!\left(-1+\frac1L+\frac{\log n}{2n}\right).
\end{equation}
\end{lemma}
\begin{proof}
Because all coefficients of $\mathscr B$ are nonnegative,
$\BellNumber n t^n/n!\le\mathscr B(t)$, proving the first claim.  Use
$\EulerE^t=n/L$ and the elementary Stirling bound
$n!\le e\sqrt n\,(n/e)^n$; the factors $\EulerE^{-1/n}$ and $\EulerE^{1/n}$ cancel and give
\eqref{eq:bell-ratio-upper}.
\end{proof}

\begin{theorem}[Uniform numerical bounds]\label{thm:bell-numerical-bounds}
For every $n\ge1$,
\begin{equation}\label{eq:bell-0792}
 \BellNumber n<\left(\frac{0.792\,n}{\log(n+1)}\right)^n.
\end{equation}
Define, for $x>1$,
\begin{equation}\label{eq:bell-dx}
 d(x)=\log\log(x+1)-\log\log x+\frac{1+\EulerE^{-1}}{\log x}.
\end{equation}
For every $\varepsilon>0$ there is
$n_0(\varepsilon)=\max\{\EulerE^4,d^{-1}(\varepsilon)\}$ such that, for
$n>n_0(\varepsilon)$,
\begin{equation}\label{eq:bell-e06}
 \BellNumber n<\left(\frac{\EulerE^{-0.6+\varepsilon}n}{\log(n+1)}\right)^n.
\end{equation}
\end{theorem}
\begin{proof}
Set
\[
 R(x)=\frac{\log(x+1)}{\log(x/\log(x+1))}
 \exp\!\left(-1+\frac1{\log(x+1)}+\frac{\log x}{2x}\right).
\]
Direct differentiation, after multiplying by the positive denominator
$2x\log x\log(x+1)\log(x/\log(x+1))^2$, shows that $R'(x)<0$ for $x\ge39$;
substitution gives $R(39)<0.792$.  Thus
\eqref{eq:bell-ratio-upper} proves \eqref{eq:bell-0792} for $n\ge39$.  The
remaining $1\le n\le38$ are obtained successively from
\eqref{eq:bell-binomial-recurrence}; substituting them gives strict inequalities.
This is a finite exact integer check.

A second differentiation gives, for $x\ge \EulerE^4$,
\[
 \log R(x)\le-0.6+d(x),\qquad d'(x)<0,\qquad \lim_{x\to\infty}d(x)=0.
\]
(The numerator after clearing positive denominators is a sum of
$-(\log x-4)^2$, $-(x-1-\log x)$, and negative multiples of these two standard
nonnegative quantities.)  Hence $d(n)<\varepsilon$ when
$n>d^{-1}(\varepsilon)$, and \eqref{eq:bell-e06} follows from
\eqref{eq:bell-ratio-upper}.
\end{proof}

\section{Saddle-point asymptotics and Lambert \texorpdfstring{$W$}{W}}
Let $r=W(n)$, so that $re^r=n$, and put $A=n(1+r)$.
\begin{theorem}[Leading Bell asymptotic]\label{thm:bell-leading-asymptotic}
As $n\to\infty$,
\begin{align}
 \BellNumber n
 &\sim \frac{n!\exp(\EulerE^r-1)}{r^n\sqrt{2\pi n(1+r)}}
 \label{eq:bell-saddle-leading}\\
 &\sim \frac1{\sqrt n}\left(\frac{n}{W(n)}\right)^{n+1/2}
       \exp\!\left(\frac{n}{W(n)}-n-1\right).
 \label{eq:bell-W-asymptotic}
\end{align}
\end{theorem}
\begin{proof}
In \eqref{eq:bell-cauchy} take the circle $z=re^{i\theta}$.  The coefficient
integral becomes
\[
 \frac{\BellNumber n}{n!}=
 \frac{\EulerE^{\EulerE^r-1}}{2\pi r^n}
 \int_{-\pi}^{\pi}
 \exp\!\left(\EulerE^{re^{i\theta}}-\EulerE^r-in\theta\right)\Differential\theta.
\]
The linear term vanishes because $re^r=n$; the quadratic term is
$-A\theta^2/2$.  On $|\theta|\le A^{-2/5}$, Taylor expansion is uniform and the
scaled integral tends to $\int_{\RealNumbers}\EulerE^{-u^2/2}\Differential u/\sqrt A$.  Outside that arc,
$\Re(\EulerE^{re^{i\theta}}-\EulerE^r)\le-cA\theta^2$ on the intermediate range and is
$\le-c'\EulerE^r$ on the remaining compact range, so the contribution is negligible.
This proves \eqref{eq:bell-saddle-leading}.  Stirling's formula and $\EulerE^r=n/r$
simplify it to \eqref{eq:bell-W-asymptotic}.
\end{proof}

The archived refined expansion used undefined symbols.  The following construction
makes every coefficient unambiguous.  Let
$\TouchardPolynomial{j}(r)=\sum_{q=0}^j\StirlingSecondKind jq r^q$ be the Touchard polynomial and let
$Z\sim N(0,1)$.  For an integer $h=O(r)$ define
\begin{equation}\label{eq:bell-saddle-correction-generator}
 \mathcal R_{r,h}(Z)=\exp\!\left(
 -\frac{ihZ}{\sqrt A}+
 \sum_{j\ge3}\frac{i^j \EulerE^r\TouchardPolynomial{j}(r)}{j!A^{j/2}}Z^j
 \right).
\end{equation}
Expand this formal exponential and take Gaussian moments termwise; retaining all
monomials of total order at most $\EulerE^{-Mr}$ defines $C_0(r,h),\ldots,C_M(r,h)$,
with $C_0=1$.

\begin{theorem}[Uniform Moser--Wyman scheme]\label{thm:bell-refined-saddle}
For every fixed $M$ and $|h|\le Cr$,
\begin{equation}\label{eq:bell-refined-saddle}
 \BellNumber{n+h}=\frac{(n+h)!\exp(\EulerE^r-1)}{r^{n+h}\sqrt{2\pi A}}
 \left(\sum_{q=0}^{M-1}C_q(r,h)+O_C(\EulerE^{-Mr}P_M(r))\right),
\end{equation}
where $P_M$ is an explicitly computable polynomial.  The first correction is
\begin{equation}\label{eq:bell-first-correction}
 C_1(r,h)=
 -\frac{h^2}{2A}-\frac{h \EulerE^r\TouchardPolynomial{3}(r)}{2A^2}
 +\frac{\EulerE^rT_4(r)}{8A^2}
 -\frac{5\EulerE^{2r}\TouchardPolynomial{3}(r)^2}{24A^3}.
\end{equation}
In particular the coefficient of $\EulerE^{-qr}$ is a polynomial in $h$ of degree at
most $2q$, with rational functions of $r$ as coefficients.
\end{theorem}
\begin{proof}
Use the same circle as in the preceding proof but replace $n$ by $n+h$.  After
$\theta=Z/\sqrt A$, the ratio of the integrand to its Gaussian quadratic part is
exactly the formal series \eqref{eq:bell-saddle-correction-generator}.  Watson's
lemma with a Gaussian majorant permits termwise integration on the central arc;
the complementary arcs have the same exponential bound uniformly for
$|h|\le Cr$.  Odd Gaussian moments vanish, so only integral powers of $\EulerE^{-r}$
remain.  Since $\Expectation Z^2=1$, $\Expectation Z^4=3$, and $\Expectation Z^6=15$, the terms from the square
of the $h$-term, its product with the cubic term, the quartic term, and the square
of the cubic term give \eqref{eq:bell-first-correction}.  The same finite
bookkeeping defines every subsequent $C_q$ and bounds the omitted Taylor
remainder.
\end{proof}

\begin{corollary}[de Bruijn logarithmic expansion]\label{cor:bell-debruijn}
As $n\to\infty$,
\begin{multline}\label{eq:bell-debruijn}
 \frac{\log\BellNumber n}{n}
 =\log n-\log\log n-1
 +\frac{\log\log n}{\log n}+\frac1{\log n}\\
 +\frac12\left(\frac{\log\log n}{\log n}\right)^2
 +O\!\left(\frac{\log\log n}{(\log n)^2}\right).
\end{multline}
\end{corollary}
\begin{proof}
Take logarithms in \eqref{eq:bell-W-asymptotic}.  The standard reversion of
$r+\log r=\log n$ gives
\[
 r=\log n-\log\log n+
\frac{\log\log n}{\log n}
 +O\!\left(\frac{(\log\log n)^2}{(\log n)^2}\right).
\]
Substitution and one further Taylor expansion of $\log r$ produce the displayed
terms; the prefactor contributes only $O(\log n/n)$.
\end{proof}

\section{Bell primes and verifiable computation}
A \emph{Bell prime} is a prime value of $\BellNumber n$.  The archive records the early
indices found by computation.  Such a list is not a symbolic theorem, but each
entry has a rigorous certificate: compute $\BellNumber n$ by the integer recurrence
\eqref{eq:bell-binomial-recurrence}, apply a probable-prime filter, and attach an
ECPP or Pratt-style primality certificate.  Composite entries are certified by one
nontrivial factor.  This distinction between an identity and a finite certified
computation will also be used for tabulated values later.

\part{Bell Polynomials and Composition}

\chapter{Exponential and ordinary Bell polynomials}
\index{Bell polynomials}

\section{Definitions}
\begin{definition}[Partial exponential Bell polynomial]
For $0\le k\le n$, define
\begin{equation}\label{eq:partial-bell-definition}
 \ExponentialPartialBellPolynomial nk(x_1,\ldots,x_{n-k+1})
 =n!\!\sum_{\substack{j_1+j_2+\cdots=k\\
                       j_1+2j_2+3j_3+\cdots=n}}
 \prod_{i\ge1}\frac{x_i^{j_i}}{(i!)^{j_i}j_i!}.
\end{equation}
Only $i\le n-k+1$ can occur.  The complete exponential Bell polynomial is
\begin{equation}\label{eq:complete-bell-definition}
 \ExponentialCompleteBellPolynomial n(x_1,\ldots,x_n)=\sum_{k=0}^{n}\ExponentialPartialBellPolynomial nk(x_1,\ldots,x_{n-k+1}),
 \qquad \ExponentialCompleteBellPolynomial0=1.
\end{equation}
\end{definition}

\begin{definition}[Partial ordinary Bell polynomial]
\begin{equation}\label{eq:ordinary-bell-definition}
 \OrdinaryPartialBellPolynomial nk(x_1,\ldots,x_{n-k+1})
 =\sum_{\substack{j_1+j_2+\cdots=k\\
                       j_1+2j_2+\cdots=n}}
 \frac{k!}{j_1!j_2!\cdots}\prod_{i\ge1}x_i^{j_i}.
\end{equation}
Equivalently,
\begin{equation}\label{eq:ordinary-exponential-scaling}
 \OrdinaryPartialBellPolynomial nk(x_1,x_2,\ldots)
 =\frac{k!}{n!}\ExponentialPartialBellPolynomial nk(1!x_1,2!x_2,\ldots).
\end{equation}
\end{definition}
\begin{proof}[Proof of \eqref{eq:ordinary-exponential-scaling}]
Substitute $i!x_i$ for $x_i$ in \eqref{eq:partial-bell-definition}; the factorials
$(i!)^{j_i}$ cancel, and multiplication by $k!/n!$ leaves exactly
\eqref{eq:ordinary-bell-definition}.
\end{proof}

\section{Combinatorial interpretation}
Give each block of size $i$ a weight $x_i$, and define the weight of a partition as
the product of its block weights.

\begin{theorem}[Weighted-partition interpretation]\label{thm:bell-poly-partitions}
$\ExponentialPartialBellPolynomial nk$ is the total weight of partitions of an $n$-element labelled set into
exactly $k$ blocks, and $\ExponentialCompleteBellPolynomial n$ is the total weight of all its partitions.
\end{theorem}
\begin{proof}
Fix multiplicities $j_i$, where $j_i$ blocks have size $i$.  The number of such
partitions is
\[
 \frac{n!}{\prod_i(i!)^{j_i}j_i!}.
\]
The factors $i!$ forget order inside each block and $j_i!$ forget order among
blocks of equal size.  Multiplication by $\prod_i x_i^{j_i}$ and summation gives
\eqref{eq:partial-bell-definition}; summing over $k$ gives the complete polynomial.
\end{proof}

Thus
\begin{align*}
 \ExponentialPartialBellPolynomial{3}{2}(x_1,x_2)&=3x_1x_2,\\
 \ExponentialPartialBellPolynomial{6}{2}(x_1,\ldots,x_5)&=6x_1x_5+15x_2x_4+10x_3^2,\\
 \ExponentialPartialBellPolynomial{6}{3}(x_1,\ldots,x_4)&=15x_1^2x_4+60x_1x_2x_3+15x_2^3.
\end{align*}
The monomials correspond to integer partitions of $n$ into $k$ parts, while their
coefficients count the labelled set partitions having the prescribed block sizes.

\section{The master generating functions}
Put
\[
 X(t)=\sum_{j\ge1}x_j\frac{t^j}{j!}.
\]
\begin{theorem}[Bell-polynomial generating functions]\label{thm:bell-poly-egf}
One has
\begin{align}
 \frac{X(t)^k}{k!}
 &=\sum_{n\ge k}\ExponentialPartialBellPolynomial nk(x_1,\ldots)\frac{t^n}{n!},\\
 \exp(uX(t))
 &=\sum_{n\ge k\ge0}\ExponentialPartialBellPolynomial nk(x_1,\ldots)u^k\frac{t^n}{n!},\\
 \exp(X(t))
 &=\sum_{n\ge0}\ExponentialCompleteBellPolynomial n(x_1,\ldots,x_n)\frac{t^n}{n!}.
\end{align}
For $\OrdinaryGeneratingFunctionOf{X}(t)=\sum_{j\ge1}x_jt^j$,
\begin{align}
 \OrdinaryGeneratingFunctionOf{X}(t)^k&=\sum_{n\ge k}\OrdinaryPartialBellPolynomial nk(x_1,\ldots)t^n,
 \label{eq:ordinary-bell-ogf}\\
 \EulerE^{u\OrdinaryGeneratingFunctionOf{X}(t)}
 &=\sum_{n\ge k\ge0}\OrdinaryPartialBellPolynomial nk(x_1,\ldots)t^n\frac{u^k}{k!}.
 \label{eq:ordinary-bell-bivariate}
\end{align}
\end{theorem}
\begin{proof}
Expand $X(t)^k$ multinomially.  Choosing $j_i$ copies of the term
$x_it^i/i!$ gives the constraints $\sum j_i=k$ and $\sum ij_i=n$, and its
coefficient is $k!\prod_i x_i^{j_i}/((i!)^{j_i}j_i!)$.  Division by $k!$ gives
\eqref{eq:partial-bell-definition}.  Summing in $k$ proves the exponential
identities.  The ordinary identities follow by the same multinomial argument or
by \eqref{eq:ordinary-exponential-scaling}.
\end{proof}

\begin{masterbox}
Nearly every algebraic identity for Bell polynomials follows by performing a simple
operation---multiplication, differentiation, substitution, or logarithm---on
\eqref{eq:partial-bell-egf} or \eqref{eq:complete-bell-egf}.  The combinatorial
proofs are the same operations interpreted as selecting, pointing, colouring, or
gluing blocks.
\end{masterbox}

\section{Tables and low-degree polynomials}
The nonzero partial polynomials through $n=6$ are
\begin{align*}
\ExponentialPartialBellPolynomial11&=x_1,\\
\ExponentialPartialBellPolynomial21&=x_2,\qquad \ExponentialPartialBellPolynomial22=x_1^2,\\
\ExponentialPartialBellPolynomial31&=x_3,\qquad \ExponentialPartialBellPolynomial32=3x_1x_2,\qquad \ExponentialPartialBellPolynomial33=x_1^3,\\
\ExponentialPartialBellPolynomial41&=x_4,\qquad \ExponentialPartialBellPolynomial42=4x_1x_3+3x_2^2,\\
\ExponentialPartialBellPolynomial43&=6x_1^2x_2,\qquad \ExponentialPartialBellPolynomial44=x_1^4,\\
\ExponentialPartialBellPolynomial51&=x_5,\qquad \ExponentialPartialBellPolynomial52=5x_1x_4+10x_2x_3,\\
\ExponentialPartialBellPolynomial53&=10x_1^2x_3+15x_1x_2^2,\\
\ExponentialPartialBellPolynomial54&=10x_1^3x_2,\qquad \ExponentialPartialBellPolynomial55=x_1^5,\\
\ExponentialPartialBellPolynomial61&=x_6,\qquad \ExponentialPartialBellPolynomial62=6x_1x_5+15x_2x_4+10x_3^2,\\
\ExponentialPartialBellPolynomial63&=15x_1^2x_4+60x_1x_2x_3+15x_2^3,\\
\ExponentialPartialBellPolynomial64&=20x_1^3x_3+45x_1^2x_2^2,\\
\ExponentialPartialBellPolynomial65&=15x_1^4x_2,\qquad \ExponentialPartialBellPolynomial66=x_1^6.
\end{align*}
Every coefficient is supplied by \eqref{eq:partial-bell-definition}, so the table
is a computation from the theorem rather than an independent assertion.

The first complete polynomials are
\begin{align*}
 \ExponentialCompleteBellPolynomial0={}&1,\\
 \ExponentialCompleteBellPolynomial1={}&x_1,\\
 \ExponentialCompleteBellPolynomial2={}&x_1^2+x_2,\\
 \ExponentialCompleteBellPolynomial3={}&x_1^3+3x_1x_2+x_3,\\
 \ExponentialCompleteBellPolynomial4={}&x_1^4+6x_1^2x_2+3x_2^2+4x_1x_3+x_4,\\
 \ExponentialCompleteBellPolynomial5={}&x_1^5+10x_1^3x_2+15x_1x_2^2+10x_1^2x_3
 +10x_2x_3+5x_1x_4+x_5,\\
 \ExponentialCompleteBellPolynomial6={}&x_1^6+15x_1^4x_2+45x_1^2x_2^2+15x_2^3
 +20x_1^3x_3+60x_1x_2x_3+10x_3^2\\
 &+15x_1^2x_4+15x_2x_4+6x_1x_5+x_6,\\
 \ExponentialCompleteBellPolynomial7={}&x_1^7+21x_1^5x_2+105x_1^3x_2^2+105x_1x_2^3
 +35x_1^4x_3+210x_1^2x_2x_3\\
 &+70x_1x_3^2+105x_2^2x_3+35x_1^3x_4+105x_1x_2x_4
 +35x_3x_4\\
 &+21x_1^2x_5+21x_2x_5+7x_1x_6+x_7.
\end{align*}

\chapter{Recurrences, derivatives, and algebraic identities}

\section{Pointing recurrences}
\begin{theorem}[Complete and partial recurrences]\label{thm:bell-poly-recurrences}
For $n\ge0$,
\begin{align}
 \ExponentialCompleteBellPolynomial{n+1}(x_1,\ldots,x_{n+1})
 &=\sum_{i=0}^{n}\binom ni x_{i+1}\ExponentialCompleteBellPolynomial{n-i}(x_1,\ldots,x_{n-i}),
 \label{eq:complete-bell-recurrence}\\
 \ExponentialPartialBellPolynomial{n+1}{k+1}
 &=\sum_{i=0}^{n-k}\binom ni x_{i+1}\ExponentialPartialBellPolynomial{n-i}{k},
 \label{eq:partial-bell-recurrence}
\end{align}
with $\ExponentialPartialBellPolynomial00=1$, $\ExponentialPartialBellPolynomial n0=0$ for $n>0$, and $\ExponentialPartialBellPolynomial0k=0$ for $k>0$.
\end{theorem}
\begin{proof}
In a weighted partition of $[n+1]$, inspect the block containing the distinguished
label $n+1$.  If that block contains $i$ further labels, choose them in $\binom ni$
ways, give the block weight $x_{i+1}$, and partition the remaining labels.  Keeping
or forgetting the total block count yields the two formulas.
\end{proof}

\begin{theorem}[Block-colour convolution]\label{thm:bell-partial-convolution}
For $k_1,k_2\ge0$,
\begin{equation}\label{eq:bell-partial-convolution}
 \ExponentialPartialBellPolynomial n{k_1+k_2}(x)
 =\frac{k_1!k_2!}{(k_1+k_2)!}
 \sum_{i=0}^{n}\binom ni\ExponentialPartialBellPolynomial i{k_1}(x)\ExponentialPartialBellPolynomial{n-i}{k_2}(x).
\end{equation}
\end{theorem}
\begin{proof}
Multiply \eqref{eq:partial-bell-egf} for $k_1$ and $k_2$.  Their product is
$X^{k_1+k_2}/(k_1!k_2!)$; comparison with the same formula for $k_1+k_2$ gives the
factor $(k_1!k_2!)/(k_1+k_2)!$ and the binomial EGF convolution.
\end{proof}

\section{Near-diagonal reduction}
\begin{theorem}\label{thm:bell-near-diagonal}
For $1\le a<n$,
\begin{multline}\label{eq:bell-near-diagonal}
 \ExponentialPartialBellPolynomial n{n-a}(x_1,\ldots,x_{a+1})\\
 =\sum_{j=a+1}^{2a}\frac{j!}{a!}\binom nj x_1^{n-j}
 \ExponentialPartialBellPolynomial a{j-a}\!\left(\frac{x_2}{2},\frac{x_3}{3},\ldots,
 \frac{x_{2a-j+2}}{2a-j+2}\right).
\end{multline}
\end{theorem}
\begin{proof}
A partition into $n-a$ blocks has total \emph{excess}
$\sum_B(|B|-1)=a$.  Let $j$ be the number of elements in nonsingleton blocks.
Then the number of those blocks is $j-a$, so $a+1\le j\le2a$.  Choose the $j$
elements, leave the other $n-j$ as singleton blocks, and count the nonsingleton
partition.  Replacing a block of size $r+1$ by an atom of size $r$ removes one
distinguished element per block; the factor $j!/a!$ and weights $x_{r+1}/(r+1)$
are exactly the labelled pointing factors.  Equivalently, make the substitution in
\eqref{eq:partial-bell-definition}; the coefficient of every multiplicity vector
on both sides is
$n!x_1^{n-j}\prod_{r\ge2}x_r^{j_r}/((r!)^{j_r}j_r!)$.
\end{proof}

The first cases, obtained by listing integer partitions of the excess, are
\begin{align}
 \ExponentialPartialBellPolynomial n1&=x_n, &
 \ExponentialPartialBellPolynomial n2&=\frac12\sum_{j=1}^{n-1}\binom njx_jx_{n-j},
 \label{eq:bell-edge-first}\\
 \ExponentialPartialBellPolynomial nn&=x_1^n, &
 \ExponentialPartialBellPolynomial n{n-1}&=\binom n2x_1^{n-2}x_2,\\
 \ExponentialPartialBellPolynomial n{n-2}&=\binom n3x_1^{n-3}x_3
 +3\binom n4x_1^{n-4}x_2^2,\\
 \ExponentialPartialBellPolynomial n{n-3}&=\binom n4x_1^{n-4}x_4
 +10\binom n5x_1^{n-5}x_2x_3
 +15\binom n6x_1^{n-6}x_2^3,\\
 \ExponentialPartialBellPolynomial n{n-4}&=\binom n5x_1^{n-5}x_5
 +5\binom n6x_1^{n-6}(3x_2x_4+2x_3^2)\notag\\
 &\quad+105\binom n7x_1^{n-7}x_2^2x_3
 +105\binom n8x_1^{n-8}x_2^4.
 \label{eq:bell-near-diagonal-cases}
\end{align}
Their coefficients are the corresponding block-size counts from
\cref{thm:bell-poly-partitions}.

\section{Partial derivatives}
\begin{theorem}[Derivative identities]\label{thm:bell-poly-derivatives}
For $i\ge1$,
\begin{align}
 \frac{\partial\ExponentialCompleteBellPolynomial n}{\partial x_i}
 &=\binom ni\ExponentialCompleteBellPolynomial{n-i},
 \label{eq:complete-bell-derivative}\\
 \frac{\partial\ExponentialPartialBellPolynomial nk}{\partial x_i}
 &=\binom ni\ExponentialPartialBellPolynomial{n-i}{k-1}.
 \label{eq:partial-bell-derivative}
\end{align}
Consequently, for differentiable $a_i(x)$,
\begin{equation}\label{eq:bell-poly-chain}
 \frac{\Differential}{\Differential x}\ExponentialPartialBellPolynomial nk(a_1(x),\ldots)
 =\sum_{i=1}^{n-k+1}\binom ni a_i'(x)\ExponentialPartialBellPolynomial{n-i}{k-1}(a_1(x),\ldots).
\end{equation}
\end{theorem}
\begin{proof}
Differentiate \eqref{eq:partial-bell-egf} with respect to $x_i$:
\[
 \frac{\partial}{\partial x_i}\frac{X(t)^k}{k!}
 =\frac{t^i}{i!}\frac{X(t)^{k-1}}{(k-1)!}.
\]
Coefficient comparison gives \eqref{eq:partial-bell-derivative}; summing over $k$
gives \eqref{eq:complete-bell-derivative}.  The ordinary chain rule then gives
\eqref{eq:bell-poly-chain}.
\end{proof}

\section{A quadratic differential recurrence}
The less familiar differential recurrence in the archive is also a disguised
coefficient identity.
\begin{theorem}\label{thm:bell-quadratic-differential}
For $n\ge2$,
\begin{align}
 \ExponentialCompleteBellPolynomial n=\frac1{n-1}\Bigg[&
 \sum_{i=2}^{n}\sum_{j=1}^{i-1}(i-1)\binom{i-2}{j-1}x_jx_{i-j}
       \frac{\partial\ExponentialCompleteBellPolynomial{n-1}}{\partial x_{i-1}}\notag\\
 &+\sum_{i=2}^{n}\sum_{j=1}^{i-1}\frac{x_{i+1}}{\binom ij}
       \frac{\partial^2\ExponentialCompleteBellPolynomial{n-1}}{\partial x_j\partial x_{i-j}}
 +\sum_{i=2}^{n}x_i\frac{\partial\ExponentialCompleteBellPolynomial{n-1}}{\partial x_{i-1}}
 \Bigg].
 \label{eq:bell-quadratic-differential}
\end{align}
Terms containing unavailable variables or negative-index Bell polynomials are zero.
\end{theorem}
\begin{proof}
Let $F(t)=\EulerE^{X(t)}=\sum_m\ExponentialCompleteBellPolynomial m t^m/m!$.  After using
\eqref{eq:complete-bell-derivative}, the first double sum is
\[
 (n-1)!\CoefficientExtraction{t}{n-1}\bigl(tX'(t)^2F(t)\bigr).
\]
The second is
$(n-1)!\CoefficientExtraction{t}{n-1}((tX''-X'+x_1)F)$, and the last is
$(n-1)!\CoefficientExtraction{t}{n-1}((X'-x_1)F)$.  Their sum is therefore
\[
 (n-1)!\CoefficientExtraction{t}{n-1}t(X''+X'^2)F
 =(n-1)!\CoefficientExtraction{t}{n-1}tF''=(n-1)\ExponentialCompleteBellPolynomial n.
\]
Division by $n-1$ proves the formula.
\end{proof}

\chapter{Specializations, inversions, determinants, and convolutions}

\section{Stirling, Bell, Lah, and Touchard specializations}
\begin{theorem}\label{thm:bell-poly-specializations}
For admissible $n,k$,
\begin{align}
 \ExponentialPartialBellPolynomial nk(0!,1!,\ldots,(n-k)!)&=\UnsignedStirlingFirstKind nk,
 \label{eq:bell-first-specialization}\\
 \ExponentialCompleteBellPolynomial n(0!,1!,\ldots,(n-1)!)&=n!,
 \label{eq:bell-factorial-complete}\\
 \ExponentialPartialBellPolynomial nk(1,1,\ldots,1)&=\StirlingSecondKind nk,
 \label{eq:bell-second-specialization}\\
 \ExponentialCompleteBellPolynomial n(1,1,\ldots,1)&=\BellNumber n,
 \label{eq:bell-number-specialization}\\
 \ExponentialPartialBellPolynomial nk(1!,2!,\ldots,(n-k+1)!)
 &=\binom{n-1}{k-1}\frac{n!}{k!}=\LahNumber{n}{k},
 \label{eq:bell-lah-specialization}\\
 \TouchardPolynomial{n}(x)&=\ExponentialCompleteBellPolynomial n(x,x,\ldots,x)
 =\sum_{k=0}^{n}\StirlingSecondKind nkx^k.
 \label{eq:touchard-bell-specialization}
\end{align}
\end{theorem}
\begin{proof}
In \eqref{eq:partial-bell-egf}, the choices $x_j=(j-1)!$, $x_j=1$, and
$x_j=j!$ make $X(t)$ respectively
\[
 -\log(1-t),\qquad \EulerE^t-1,\qquad \frac{t}{1-t}.
\]
The resulting coefficient formulae are exactly the column EGFs for unsigned
first-kind Stirling numbers, second-kind Stirling numbers, and Lah numbers proved
earlier.  Summing over $k$ yields \eqref{eq:bell-factorial-complete} and
\eqref{eq:bell-number-specialization}.  Finally $x_j=x$ gives
$\exp(x(\EulerE^t-1))$, the EGF of $\TouchardPolynomial{n}(x)$.
\end{proof}

\section{Homogeneity and a useful evaluation}
\begin{theorem}[Bigraded homogeneity]\label{thm:bell-bihomogeneous}
For scalars $\alpha,\beta$,
\begin{equation}\label{eq:bell-bihomogeneous}
 \ExponentialPartialBellPolynomial nk(\alpha\beta x_1,\alpha\beta^2x_2,\ldots)
 =\alpha^k\beta^n\ExponentialPartialBellPolynomial nk(x_1,x_2,\ldots).
\end{equation}
Moreover,
\begin{equation}\label{eq:bell-linear-arguments}
 \ExponentialPartialBellPolynomial nk(1,2,3,\ldots,n-k+1)=\binom nk k^{n-k}.
\end{equation}
\end{theorem}
\begin{proof}
Every monomial in \eqref{eq:partial-bell-definition} has
$\sum j_i=k$ and $\sum ij_i=n$, proving the first formula.  For the second,
$x_j=j$ gives $X(t)=te^t$; hence
\[
 \frac{X(t)^k}{k!}=\frac{t^ke^{kt}}{k!}.
\]
The coefficient of $t^n/n!$ is
$n!k^{n-k}/(k!(n-k)!)=\binom nk k^{n-k}$.
\end{proof}

\section{Addition and shifted-zero identities}
\begin{theorem}[Complete Bell addition formula]\label{thm:complete-bell-addition}
\begin{equation}\label{eq:complete-bell-addition}
 \ExponentialCompleteBellPolynomial n(x_1+y_1,\ldots,x_n+y_n)
 =\sum_{i=0}^{n}\binom ni\ExponentialCompleteBellPolynomial{n-i}(x)\ExponentialCompleteBellPolynomial i(y).
\end{equation}
\end{theorem}
\begin{proof}
The complete EGF for the left side is
$\exp(X(t)+Y(t))=\exp X(t)\exp Y(t)$.  Coefficient extraction gives the binomial
convolution.
\end{proof}

\begin{theorem}[Insertion of $q$ leading zeros]\label{thm:bell-leading-zeros}
For $q\ge0$,
\begin{multline}\label{eq:bell-leading-zeros}
 \ExponentialPartialBellPolynomial nk\!\left(
 \frac{x_{q+1}}{\binom{q+1}{q}},
 \frac{x_{q+2}}{\binom{q+2}{q}},\ldots\right)\\
 =\frac{n!(q!)^k}{(n+qk)!}
 \ExponentialPartialBellPolynomial{n+qk}{k}(\underbrace{0,\ldots,0}_{q\text{ entries}},
 x_{q+1},x_{q+2},\ldots).
\end{multline}
\end{theorem}
\begin{proof}
Since
\[
 \frac{x_{q+i}}{\binom{q+i}{q}}\frac{t^i}{i!}
 =q!\,\frac{x_{q+i}t^i}{(q+i)!},
\]
the $k$th power of the left component series is $(q!)^k t^{-qk}$ times the
$k$th power of the component series on the right.  Comparing the coefficient of
$t^n$ in \eqref{eq:partial-bell-egf} gives the factorial factor displayed.
\end{proof}

\section{Exponential--logarithmic inversion}
\begin{theorem}[Bell transform and its inverse]\label{thm:bell-transform-inverse}
Suppose
\begin{equation}\label{eq:bell-transform-y}
 y_n=\ExponentialCompleteBellPolynomial n(x_1,\ldots,x_n)=\sum_{k=1}^{n}\ExponentialPartialBellPolynomial nk(x_1,\ldots).
\end{equation}
Then
\begin{equation}\label{eq:bell-transform-x}
 x_n=\sum_{k=1}^{n}(-1)^{k-1}(k-1)!
 \ExponentialPartialBellPolynomial nk(y_1,\ldots,y_{n-k+1}).
\end{equation}
\end{theorem}
\begin{proof}
Let $X(t)=\sum_{n\ge1}x_nt^n/n!$ and
$Y(t)=\sum_{n\ge1}y_nt^n/n!$.  Equation
\eqref{eq:complete-bell-egf} says $1+Y=\EulerE^X$, so
\[
 X=\log(1+Y)=\sum_{k\ge1}\frac{(-1)^{k-1}}{k}Y^k.
\]
Use \eqref{eq:partial-bell-egf} to extract the coefficient of $t^n/n!$.
\end{proof}

More generally, if $f$ and $g=f^{-1}$ are locally inverse and
\begin{equation}\label{eq:general-bell-transform}
 y_n=\sum_{k=0}^{n}f^{(k)}(a)\ExponentialPartialBellPolynomial nk(x_1,\ldots),
\end{equation}
then
\begin{equation}\label{eq:general-bell-inverse}
 x_n=\sum_{k=0}^{n}g^{(k)}(f(a))\ExponentialPartialBellPolynomial nk(y_1,\ldots).
\end{equation}
Indeed, these are the derivative coefficients of $f\circ u$ and
$g\circ f\circ u=u$; a complete proof is given with Fa\`a di Bruno's formula in
\cref{eq:faa-bell}.

\section{Two Hessenberg determinants}
For $n\ge1$, define $\BellHessenbergMatrixH{n}=(h_{ij})$ and $\BellHessenbergMatrixK{n}=(k_{ij})$ by
\[
 h_{ij}=\begin{cases}
 \binom{n-i}{j-i}x_{j-i+1},&j\ge i,\\
 -1,&j=i-1,\\0,&j<i-1,
 \end{cases}
 \qquad
 k_{ij}=\begin{cases}
 x_{j-i+1}/(j-i)!,&j\ge i,\\
 -(i-1),&j=i-1,\\0,&j<i-1.
 \end{cases}
\]
Thus $\BellHessenbergMatrixH{n}$ begins with
\[
\begin{bmatrix}
 x_1&\binom{n-1}{1}x_2&\binom{n-1}{2}x_3&\cdots&x_n\\
 -1&x_1&\binom{n-2}{1}x_2&\cdots&x_{n-1}\\
 0&-1&x_1&\cdots&x_{n-2}\\
 \vdots&&\ddots&\ddots&\vdots\\
 0&\cdots&0&-1&x_1
\end{bmatrix},
\]
while $\BellHessenbergMatrixK{n}$ has first row
$x_1,x_2/1!,\ldots,x_n/(n-1)!$ and subdiagonal
$-1,-2,\ldots,-(n-1)$.

\begin{theorem}[Bell determinants]\label{thm:bell-determinants}
\begin{equation}\label{eq:bell-determinants}
 \ExponentialCompleteBellPolynomial n(x_1,\ldots,x_n)=\det \BellHessenbergMatrixH{n}=\det \BellHessenbergMatrixK{n}.
\end{equation}
\end{theorem}
\begin{proof}
Expansion of the upper Hessenberg determinant $\det \BellHessenbergMatrixH{n}$ along its first row gives
\[
 D_n=\sum_{j=1}^{n}\binom{n-1}{j-1}x_jD_{n-j},\qquad D_0=1;
\]
the cofactor sign is cancelled by the product of the $j-1$ entries $-1$ on the
subdiagonal.  This is exactly \eqref{eq:complete-bell-recurrence}, so
$D_n=\ExponentialCompleteBellPolynomial n$.

Let $J$ reverse the coordinate order and let
$D=\DiagonalOperator(0!,1!,\ldots,(n-1)!)$.  A direct entrywise calculation gives
\[
 \BellHessenbergMatrixK{n}=D(J\BellHessenbergMatrixH{n}^{\mathsf T}J)D^{-1}.
\]
Indeed,
$\frac{(i-1)!}{(j-1)!}\binom{j-1}{i-1}=1/(j-i)!$ above the diagonal, and the
subdiagonal becomes $-(i-1)$.  Similarity, transposition, and reversal preserve
the determinant, proving the second equality.
\end{proof}

\section{The diamond convolution}
For sequences $x=(x_1,x_2,\ldots)$ and $y=(y_1,y_2,\ldots)$ define
\begin{equation}\label{eq:diamond}
 (x\mathbin\diamondsuit y)_n=\sum_{j=1}^{n-1}\binom njx_jy_{n-j}.
\end{equation}
This operation is commutative and associative because it corresponds to
multiplication of the EGFs $X(t)$ and $Y(t)$.

\begin{theorem}\label{thm:diamond-bell}
If $x^{k\diamondsuit}$ denotes the $k$-fold diamond power, then
\begin{equation}\label{eq:diamond-bell}
 \ExponentialPartialBellPolynomial nk(x_1,\ldots,x_{n-k+1})=\frac{(x^{k\diamondsuit})_n}{k!}.
\end{equation}
\end{theorem}
\begin{proof}
The EGF of $x^{k\diamondsuit}$ is $X(t)^k$.  Compare with
\eqref{eq:partial-bell-egf}.
\end{proof}
For example,
\[
 x\diamondsuit x=(0,2x_1^2,6x_1x_2,8x_1x_3+6x_2^2,\ldots)
\]
and
$(x^{3\diamondsuit})_4=36x_1^2x_2$, whence
$\ExponentialPartialBellPolynomial43=6x_1^2x_2$.

\chapter{Applications of Bell polynomials}

\section{Symmetric functions and Newton identities}
Let $e_n$ be the elementary symmetric function and
$p_n=\sum_i u_i^n$ the power sum in an alphabet $u_1,u_2,\ldots$.
\begin{theorem}\label{thm:bell-symmetric-functions}
\begin{align}
 e_n&=\frac1{n!}\ExponentialCompleteBellPolynomial n(p_1,-1!p_2,2!p_3,\ldots,(-1)^{n-1}(n-1)!p_n)
 \label{eq:elementary-via-bell}\\
 &=\frac{(-1)^n}{n!}\ExponentialCompleteBellPolynomial n(-p_1,-1!p_2,-2!p_3,\ldots,-(n-1)!p_n),
 \notag\\
 p_n&=\frac{(-1)^{n-1}}{(n-1)!}
 \sum_{k=1}^{n}(-1)^{k-1}(k-1)!
 \ExponentialPartialBellPolynomial nk(1!e_1,2!e_2,\ldots)
 \label{eq:powersum-via-bell}\\
 &=(-1)^n n\sum_{k=1}^{n}\frac1k
 \OrdinaryPartialBellPolynomial nk(-e_1,-e_2,\ldots).
 \notag
\end{align}
\end{theorem}
\begin{proof}
The classical product identity
\[
 E(t)=\sum_{n\ge0}e_nt^n=\prod_i(1+u_it)
 =\exp\!\left(\sum_{r\ge1}(-1)^{r-1}p_r\frac{t^r}{r}\right)
\]
and \eqref{eq:complete-bell-egf} give \eqref{eq:elementary-via-bell}.
Taking logarithms,
$\sum_{r\ge1}(-1)^{r-1}p_rt^r/r=\log(1+\sum_{j\ge1}e_jt^j)$, and expanding
the logarithm through ordinary Bell polynomials yields the second pair of forms.
The exponential form follows from \eqref{eq:ordinary-exponential-scaling}.
\end{proof}

\begin{corollary}[Determinant from traces]\label{cor:det-traces-bell}
For an $n\times n$ matrix $A$, set
$\TraceBellArgument{k}=-(k-1)!\TraceOperator(A^k)$.  Then
\begin{equation}\label{eq:det-traces-bell}
 \det A=\frac{(-1)^n}{n!}\ExponentialCompleteBellPolynomial n\bigl(
   \TraceBellArgument{1},\TraceBellArgument{2},\ldots,\TraceBellArgument{n}\bigr).
\end{equation}
\end{corollary}
\begin{proof}
Apply \eqref{eq:elementary-via-bell} to the eigenvalues of $A$, using
$e_n=\det A$ and $p_k=\TraceOperator(A^k)$.  Both sides are polynomial in the matrix entries,
so the argument remains valid for nondiagonalizable matrices.
\end{proof}

\section{Cycle indices}
Let $a_j$ mark cycles of length $j$ in a permutation.
\begin{theorem}\label{thm:cycle-index-bell}
The cycle index of the symmetric group is
\begin{equation}\label{eq:cycle-index-bell}
 Z(S_n)=\frac1{n!}\ExponentialCompleteBellPolynomial n(0!a_1,1!a_2,\ldots,(n-1)!a_n).
\end{equation}
\end{theorem}
\begin{proof}
A permutation is a set of cycles.  A labelled cycle of size $j$ has $(j-1)!$
linear realizations, so its component EGF is $(j-1)!a_jt^j/j!=a_jt^j/j$.
The exponential formula gives
$\exp(\sum_{j\ge1}a_jt^j/j)$; comparison with
\eqref{eq:complete-bell-egf} proves the identity after division by $n!$.
\end{proof}

\section{Moments and cumulants}
Let $M(t)=\Expectation \EulerE^{tX}=\sum_{n\ge0}\mu'_nt^n/n!$ and
$K(t)=\log M(t)=\sum_{n\ge1}\kappa_nt^n/n!$.
\begin{theorem}[Moment--cumulant relations]\label{thm:moment-cumulant}
\begin{align}
 \mu'_n&=\ExponentialCompleteBellPolynomial n(\kappa_1,\ldots,\kappa_n)
 =\sum_{k=1}^{n}\ExponentialPartialBellPolynomial nk(\kappa_1,\ldots),
 \label{eq:moments-from-cumulants}\\
 \kappa_n&=\sum_{k=1}^{n}(-1)^{k-1}(k-1)!
 \ExponentialPartialBellPolynomial nk(\mu'_1,\ldots,\mu'_{n-k+1}).
 \label{eq:cumulants-from-moments}
\end{align}
\end{theorem}
\begin{proof}
The identity $M=\EulerE^K$ is \eqref{eq:complete-bell-egf}; its logarithmic inverse is
\cref{thm:bell-transform-inverse}.
\end{proof}

\section{Hermite polynomials}
The probabilists' Hermite polynomials are defined by
\[
 \EulerE^{xt-t^2/2}=\sum_{n\ge0}\operatorname{He}_n(x)\frac{t^n}{n!}.
\]
Comparison with \eqref{eq:complete-bell-egf} gives and proves
\begin{equation}\label{eq:hermite-bell}
 \operatorname{He}_n(x)=\ExponentialCompleteBellPolynomial n(x,-1,0,\ldots,0).
\end{equation}
In particular, the usual recurrence and differential equation follow by
differentiating this generating function.

\section{Polynomial sequences of binomial type}
Let
$h(t)=\sum_{j\ge1}a_jt^j/j!$ with $a_1\ne0$, and define
\begin{equation}\label{eq:binomial-type-bell}
 p_n(x)=\ExponentialCompleteBellPolynomial n(a_1x,a_2x,\ldots,a_nx)
 =\sum_{k=0}^{n}\ExponentialPartialBellPolynomial nk(a_1,\ldots,a_{n-k+1})x^k.
\end{equation}

\begin{theorem}\label{thm:binomial-type-bell}
The polynomials $p_n$ satisfy
\begin{align}
 \sum_{n\ge0}p_n(x)\frac{t^n}{n!}&=\EulerE^{xh(t)},
 \label{eq:binomial-type-egf}\\
 p_n(x+y)&=\sum_{k=0}^{n}\binom nkp_k(x)p_{n-k}(y).
 \label{eq:binomial-type-identity}
\end{align}
If $Q=h^{-1}(D_x)$ is the associated delta operator, then
\begin{equation}\label{eq:delta-operator}
 Qp_n(x)=np_{n-1}(x).
\end{equation}
\end{theorem}
\begin{proof}
Substitute $x a_j$ in \eqref{eq:complete-bell-egf}.  The product
$\EulerE^{(x+y)h}=\EulerE^{xh}\EulerE^{yh}$ gives the binomial identity.  Finally
$h^{-1}(D_x)\EulerE^{xh(t)}=h^{-1}(h(t))\EulerE^{xh(t)}=t \EulerE^{xh(t)}$; coefficient comparison
gives \eqref{eq:delta-operator}.
\end{proof}

\begin{auditbox}
In some sources the symbol $B_n(x)$ denotes the Touchard polynomial, in others a
complete Bell polynomial, a Bernoulli polynomial, or a binary-tree sequence.  The
notations $\TouchardPolynomial{n}$, $\ExponentialCompleteBellPolynomial n$, $\beta_n(x)$, and $\CatalanNumber{n}$ are kept separate here; this
is essential when the objects occur in the same inversion formula.
\end{auditbox}

\part{Eulerian Numbers and Descent Calculus}

\chapter{Ordinary Eulerian numbers}
\index{Eulerian numbers}

\section{Descents, symmetry, and the basic recurrence}
For a permutation $\pi=\pi_1\cdots\pi_n$, a descent is an index
$i\in\{1,\ldots,n-1\}$ with $\pi_i>\pi_{i+1}$.
\begin{definition}
The Eulerian number $\TypeAEulerianNumber nk=A(n,k)$ counts permutations of $[n]$ with
exactly $k$ descents.  Put
\[
 \TypeAEulerianPolynomial{0}(t)=1,
 \qquad
 \TypeAEulerianPolynomial{n}(t)=\sum_{k=0}^{n-1}\TypeAEulerianNumber nk t^k\quad(n\ge1).
\]
Outside $0\le k<n$ the number is zero.
\end{definition}

\begin{theorem}[Eulerian recurrence]\label{thm:eulerian-recurrence}
\begin{equation}\label{eq:eulerian-recurrence}
 \TypeAEulerianNumber nk=(k+1)\TypeAEulerianNumber{n-1}{k}
 +(n-k)\TypeAEulerianNumber{n-1}{k-1}.
\end{equation}
Equivalently,
\begin{equation}\label{eq:eulerian-poly-recurrence}
 \TypeAEulerianPolynomial{n}(t)=\bigl(1+(n-1)t\bigr)\TypeAEulerianPolynomial{n-1}(t)
       +t(1-t)\TypeAEulerianPolynomial{n-1}'(t).
\end{equation}
\end{theorem}
\begin{proof}
Insert $n$ into a permutation of $[n-1]$.  If the old permutation has $k$
descents, insertion after one of those $k$ descents or at the end preserves the
number of descents, giving $k+1$ choices.  If it has $k-1$ descents, insertion in
any of the remaining $n-k$ slots creates one new descent.  This proves
\eqref{eq:eulerian-recurrence}.  Multiply by $t^k$, sum over $k$, and collect the
terms containing $k$ to obtain \eqref{eq:eulerian-poly-recurrence}.
\end{proof}

Complementation $\pi_i\mapsto n+1-\pi_i$ interchanges descents and ascents.
Therefore
\begin{equation}\label{eq:eulerian-symmetry}
 \TypeAEulerianNumber nk=\TypeAEulerianNumber n{n-1-k},
 \qquad
 \TypeAEulerianPolynomial{n}(t)=t^{n-1}\TypeAEulerianPolynomial{n}(t^{-1}).
\end{equation}
In particular $\TypeAEulerianNumber n0=\TypeAEulerianNumber n{n-1}=1$ and
$\sum_k\TypeAEulerianNumber nk=n!$.  The next diagonal is
\begin{equation}\label{eq:eulerian-k1}
 \TypeAEulerianNumber n1=2^n-(n+1),
\end{equation}
which follows either from the recurrence or from the explicit formula below.

\section{Exponential generating function and explicit formula}
\begin{theorem}[Eulerian EGF]\label{thm:eulerian-egf}
\begin{equation}\label{eq:eulerian-egf}
 \sum_{n\ge0}\TypeAEulerianPolynomial{n}(t)\frac{x^n}{n!}
 =\frac{t-1}{t-\EulerE^{(t-1)x}}
 =\frac{(1-t)\EulerE^{(1-t)x}}{1-te^{(1-t)x}}.
\end{equation}
\end{theorem}
\begin{proof}
Let $F(x,t)$ denote the left side.  Multiplying
\eqref{eq:eulerian-poly-recurrence} by $x^{n-1}/(n-1)!$ and summing gives the
first-order equation
\[
 (1-tx)F_x=t(1-t)F_t+F,
 \qquad F(0,t)=1.
\]
Direct substitution verifies that the displayed rational-exponential function
solves this equation and initial condition.  Formal uniqueness follows by
recursively comparing coefficients of $x$.
\end{proof}

Expanding the denominator geometrically gives the central power-series identity.
\begin{theorem}\label{thm:eulerian-power-series}
For $n\ge0$,
\begin{equation}\label{eq:eulerian-power-series}
 \sum_{m=0}^{\infty}(m+1)^n t^m
 =\frac{\TypeAEulerianPolynomial{n}(t)}{(1-t)^{n+1}}.
\end{equation}
Consequently,
\begin{equation}\label{eq:eulerian-explicit}
 \TypeAEulerianNumber nk=\sum_{j=0}^{k}(-1)^j\binom{n+1}{j}(k+1-j)^n.
\end{equation}
\end{theorem}
\begin{proof}
From \eqref{eq:eulerian-egf},
\[
 F(x,t)=(1-t)\sum_{m\ge0}t^m \EulerE^{(m+1)(1-t)x}.
\]
Extracting $x^n/n!$ proves \eqref{eq:eulerian-power-series}.  Multiply by
$(1-t)^{n+1}$ and extract $t^k$ to obtain
\eqref{eq:eulerian-explicit}.  Formula \eqref{eq:eulerian-k1} is its case $k=1$.
\end{proof}

\begin{theorem}[A second polynomial recurrence]\label{thm:eulerian-binomial-recurrence}
For $n>1$,
\begin{equation}\label{eq:eulerian-binomial-recurrence}
 \TypeAEulerianPolynomial{n}(t)=\sum_{k=0}^{n-1}\binom nk \TypeAEulerianPolynomial{k}(t)(t-1)^{n-1-k}.
\end{equation}
\end{theorem}
\begin{proof}
The EGF of the right sides, summed over $n\ge1$, is
\[
 F(x,t)\frac{\EulerE^{(t-1)x}-1}{t-1}.
\]
Using \eqref{eq:eulerian-egf}, this equals $F(x,t)-1$, which is the EGF of the
left sides.
\end{proof}

\section{Worpitzky's identity and sums of powers}
\begin{theorem}[Worpitzky]\label{thm:worpitzky}
For every indeterminate $x$ and $n\ge1$,
\begin{equation}\label{eq:worpitzky}
 x^n=\sum_{k=0}^{n-1}\TypeAEulerianNumber nk\binom{x+k}{n}.
\end{equation}
\end{theorem}
\begin{proof}
The coefficient of $t^{x-1}$ in
$t^{-k}(1-t)^{-n-1}$ is $\binom{x+k}{n}$ for positive integer $x$.
Equivalently, taking the coefficient of $t^{x-1}$ in
\eqref{eq:eulerian-power-series} after multiplying by the numerator gives
\eqref{eq:worpitzky}.  Since both sides are polynomials of degree $n$ in $x$ and
agree for all positive integers, they agree identically.

A direct combinatorial proof standardizes a word in $[x]^n$: its stable sorting
permutation records the relative order, and its $k$ descents specify $k$ mandatory
separations among the $n$ sorted positions.  Stars and bars gives
$\binom{x+k}{n}$ words for each permutation with $k$ descents.
\end{proof}

\begin{corollary}[Power sums]\label{cor:eulerian-power-sum}
For $m,n\ge1$,
\begin{equation}\label{eq:eulerian-power-sum}
 \sum_{r=1}^{m}r^n
 =\sum_{k=0}^{n-1}\TypeAEulerianNumber nk\binom{m+k+1}{n+1}.
\end{equation}
\end{corollary}
\begin{proof}
Sum \eqref{eq:worpitzky} over $x=1,\ldots,m$ and use the hockey-stick identity
$\sum_{x=1}^{m}\binom{x+k}{n}=\binom{m+k+1}{n+1}$; the omitted lower boundary
terms vanish for $0\le k<n$.
\end{proof}

\section{Negative polylogarithms}
For $|z|<1$, $\Polylogarithm_{-n}(z)=\sum_{m\ge1}m^nz^m$.  Shifting $m$ in
\eqref{eq:eulerian-power-series} gives
\begin{equation}\label{eq:negative-polylog-eulerian}
 \Polylogarithm_{-n}(z)=\frac{z\TypeAEulerianPolynomial{n}(z)}{(1-z)^{n+1}}
 =\frac{1}{(1-z)^{n+1}}
   \sum_{k=0}^{n-1}\TypeAEulerianNumber nk z^{n-k},
\end{equation}
where the second equality uses symmetry \eqref{eq:eulerian-symmetry}.  Since the
right side is rational, it also supplies the meromorphic continuation of the
negative polylogarithm.

\section{Eulerian and Stirling bases}
\begin{theorem}\label{thm:eulerian-stirling}
\begin{equation}\label{eq:eulerian-stirling}
 \TypeAEulerianPolynomial{n}(t)=\sum_{k=0}^{n}\StirlingSecondKind nk k!(t-1)^{n-k}.
\end{equation}
\end{theorem}
\begin{proof}
Use $m^n=\sum_k\StirlingSecondKind nk\FallingFactorial mk$ in
$\sum_{m\ge0}m^nt^m$.  Since
$\sum_{m\ge0}\FallingFactorial mk t^m=k!t^k/(1-t)^{k+1}$, multiplication by
$(1-t)^{n+1}/t$ yields
$\TypeAEulerianPolynomial{n}(t)=\sum_k\StirlingSecondKind nk k!t^{k-1}(1-t)^{n-k}$.  Replacing $t$ by $1/t$ and
using \eqref{eq:eulerian-symmetry}, or simply expanding around $t=1$, gives the
equivalent displayed form.
\end{proof}

\section{An Irwin--Hall integral and hypersimplex volumes}
\begin{theorem}\label{thm:eulerian-irwin-hall}
For every function $f:\RealNumbers\to\ComplexNumbers$ for which the following finite sum is defined,
\begin{equation}\label{eq:eulerian-irwin-hall}
 \sum_{k=0}^{n-1}\TypeAEulerianNumber nk f(k)
 =n!\int_{[0,1]^n}f\!\left(\Floor{x_1+\cdots+x_n}\right)
 \Differential x_1\cdots\Differential x_n.
\end{equation}
In particular, the slab
\[
 \Delta_{n,k}=\{x\in[0,1]^n:k\le x_1+\cdots+x_n<k+1\}
\]
has volume $\TypeAEulerianNumber nk/n!$.
\end{theorem}
\begin{proof}
The volume of $\{x\in[0,1]^n:\sum x_i<y\}$ is obtained from the simplex
$\{x_i\ge0:\sum x_i<y\}$ by inclusion--exclusion over coordinates exceeding $1$:
\[
 V_n(y)=\frac1{n!}\sum_{j=0}^{\Floor{y}}(-1)^j\binom nj(y-j)^n.
\]
Thus the volume of the $k$th slab is $V_n(k+1)-V_n(k)$.  Combining adjacent terms
by Pascal's identity gives
\[
 \frac1{n!}\sum_{j=0}^{k}(-1)^j\binom{n+1}{j}(k+1-j)^n
 =\frac1{n!}\TypeAEulerianNumber nk
\]
by \eqref{eq:eulerian-explicit}.  Summing $f(k)$ times these volumes proves the
integral identity.
\end{proof}

\section{Alternating sums and Bernoulli numbers}
\begin{theorem}\label{thm:eulerian-alternating}
For $n>0$,
\begin{equation}\label{eq:eulerian-alternating}
 \sum_{k=0}^{n-1}(-1)^k\TypeAEulerianNumber nk
 =2^{n+1}(2^{n+1}-1)\frac{\beta_{n+1}}{n+1}.
\end{equation}
For $n>1$,
\begin{align}
 \sum_{k=0}^{n-1}\frac{(-1)^k\TypeAEulerianNumber nk}{\binom{n-1}{k}}&=0,
 \label{eq:eulerian-reciprocal-zero}\\
 \sum_{k=0}^{n-1}\frac{(-1)^k\TypeAEulerianNumber nk}{\binom nk}&=(n+1)\beta_n.
 \label{eq:eulerian-reciprocal-bernoulli}
\end{align}
\end{theorem}
\begin{proof}
At $t=-1$, \eqref{eq:eulerian-egf} becomes
$2/(1+\EulerE^{-2x})=1+\tanh x$.  The Bernoulli generating function gives
\[
 \tanh x=\sum_{n\ge1}
 2^{n+1}(2^{n+1}-1)\frac{\beta_{n+1}}{n+1}\frac{x^n}{n!},
\]
which proves \eqref{eq:eulerian-alternating}.

For the remaining identities use
\[
 \binom{n-1}{k}^{-1}=n\int_0^1u^k(1-u)^{n-k-1}\Differential u,
 \quad
 \binom nk^{-1}=(n+1)\int_0^1u^k(1-u)^{n-k}\Differential u.
\]
Substitute $t=-u/(1-u)$ in \eqref{eq:eulerian-stirling}; then
\begin{equation}\label{eq:eulerian-beta-transform}
 (1-u)^n\TypeAEulerianPolynomial{n}\!\left(-\frac{u}{1-u}\right)
 =\sum_{j=0}^{n}(-1)^{n-j}\StirlingSecondKind njj!(1-u)^j.
\end{equation}
For the first integral divide by $1-u$ and integrate.  The result is
$n\sum_{j=1}^n(-1)^{n-j}\StirlingSecondKind nj(j-1)!=0$ for $n>1$, because it is the
coefficient of $x^n/n!$ in
$\log(1+(\EulerE^x-1))=x$.  For the second integral one obtains
$(n+1)\sum_j(-1)^{n-j}\StirlingSecondKind njj!/(j+1)$.  The standard Stirling formula
$\beta_n=\sum_j(-1)^jj!\StirlingSecondKind nj/(j+1)$ proves the claim; for odd $n>1$ both
sides vanish, and for even $n$ the signs coincide.
\end{proof}

\chapter{Geometric interpretations}

\section{The cube and its Ehrhart numerator}
For a lattice polytope $P$ of dimension $d$, its Ehrhart series has the form
$\sum_{m\ge0}|mP\cap\IntegerNumbers^d|t^m=h^*(t)/(1-t)^{d+1}$.
For the unit cube $C=[0,1]^n$ one has $|mC\cap\IntegerNumbers^n|=(m+1)^n$; hence
\eqref{eq:eulerian-power-series} proves:
\begin{theorem}
The $h^*$-polynomial of the standard $n$-cube is $\TypeAEulerianPolynomial{n}(t)$, so its $h^*$-vector is
the $n$th Eulerian row.
\end{theorem}
\begin{proof}
By the defining Ehrhart identity and the lattice-point count above,
\[
 \frac{h^*_{C}(t)}{(1-t)^{n+1}}
   =\sum_{m\ge0}(m+1)^nt^m.
\]
Equation~\eqref{eq:eulerian-power-series} identifies the right side with
$\TypeAEulerianPolynomial{n}(t)/(1-t)^{n+1}$.  Multiplication by $(1-t)^{n+1}$ therefore gives
$h^*_{C}(t)=\TypeAEulerianPolynomial{n}(t)$ coefficient by coefficient.
\end{proof}

The regions $\Delta_{n,k}$ in \cref{thm:eulerian-irwin-hall} are half-open
hypersimplices.  Their normalized volumes are $n!\operatorname{vol}\Delta_{n,k}
=\TypeAEulerianNumber nk$.

\section{The permutohedron}
The type-$A$ permutohedron is the convex hull of the permutations of
$(1,2,\ldots,n)$; its dimension is $n-1$.  Faces of its braid fan are indexed by
ordered set partitions.  There are $k!\StirlingSecondKind nk$ ordered partitions into $k$
blocks.

\begin{theorem}
The $h$-polynomial of the simple polytope dual to the type-$A$ permutohedron is
$\TypeAEulerianPolynomial{n}(t)$.
\end{theorem}
\begin{proof}
For a simple $(n-1)$-polytope the $h$-polynomial is the transform
$h(t)=\sum_F(t-1)^{\dim F}$ over faces of the dual simplicial complex.  Grouping
braid-fan faces by their $k$ ordered blocks gives
$\sum_k k!\StirlingSecondKind nk(t-1)^{n-k}$, which is $\TypeAEulerianPolynomial{n}(t)$ by
\eqref{eq:eulerian-stirling}.
\end{proof}

\chapter{Type-B and second-order Eulerian numbers}

\section{Signed permutations and the corrected type-B indexing}
A signed permutation is a word $\pi_1\cdots\pi_n$ whose absolute values form a
permutation of $[n]$.  Put $\pi_0=0$ and call
$i\in\{0,\ldots,n-1\}$ a type-$B$ descent when $\pi_i>\pi_{i+1}$.
Let $\TypeBEulerianNumber{n}{k}$ count signed permutations with $k$ such descents, and set
$\TypeBEulerianPolynomial{n}(t)=\sum_{k=0}^{n}\TypeBEulerianNumber{n}{k}t^k$.

\begin{auditbox}
The archived display reverses the two entries in the descent inequality and pairs
that convention with an incompatible explicit formula.  The definition above is
the standard Coxeter descent convention.  It gives the archived table
$1,6,1$ for $n=2$ and the midpoint-Eulerian generating function below.
\end{auditbox}

\begin{theorem}[Type-$B$ recurrence and Worpitzky series]\label{thm:typeB-eulerian}
\begin{align}
 \TypeBEulerianNumber{n}{k}&=(2k+1)\TypeBEulerianNumber{n-1}{k}+(2n-2k+1)\TypeBEulerianNumber{n-1}{k-1},
 \label{eq:typeB-recurrence}\\
 \sum_{m\ge0}(2m+1)^nt^m&=\frac{\TypeBEulerianPolynomial{n}(t)}{(1-t)^{n+1}},
 \label{eq:typeB-power-series}\\
 \TypeBEulerianNumber{n}{k}&=\sum_{j=0}^{k}(-1)^j\binom{n+1}{j}
             \bigl(2(k-j)+1\bigr)^n.
 \label{eq:typeB-explicit}
\end{align}
The first rows are
\[
1;\quad 1,1;\quad1,6,1;\quad1,23,23,1;\quad1,76,230,76,1.
\]
\end{theorem}
\begin{proof}
Insert either $n$ or $-n$ into a signed permutation of size $n-1$.  Marking the
slots according to whether insertion preserves or creates a descent gives
$2k+1$ preserving slots for an object with $k$ descents and
$2n-2k+1$ creating slots for one with $k-1$ descents; this proves the recurrence.

A type-$B$ standardization partitions the $(2m+1)^n$ words in the ordered alphabet
$-m<\cdots<-1<0<1<\cdots<m$ by their signed sorting permutation.  If that
permutation has $k$ type-$B$ descents, stars and bars gives
$\binom{m+n-k}{n}$ compatible weakly increasing absolute levels.  Hence
\[
 (2m+1)^n=\sum_{k=0}^{n}\TypeBEulerianNumber{n}{k}\binom{m+n-k}{n}.
\]
Multiply by $t^m$, sum in $m$, and use
$\sum_{m\ge0}\binom{m+n-k}{n}t^m=t^k/(1-t)^{n+1}$ to obtain
\eqref{eq:typeB-power-series}.  Multiplication by $(1-t)^{n+1}$ and coefficient
extraction gives \eqref{eq:typeB-explicit}.
\end{proof}

As in type $A$, these coefficients form the $h$-vector of the simple polytope dual
to the type-$B$ permutohedron; the proof repeats the braid-fan argument with signed
ordered partitions.

\section{Stirling permutations and second-order Eulerian numbers}
A \emph{Stirling permutation of order $n$} is a permutation of the multiset
$\{1,1,2,2,\ldots,n,n\}$ such that all entries between the two copies of $i$ are
larger than $i$.  Let $\SecondOrderEulerianNumber nk$ count those with $k$ ordinary descents and put
\[
 \SecondOrderEulerianPolynomial{n}(t)=\sum_{k=0}^{n-1}\SecondOrderEulerianNumber nk t^k,
 \qquad \SecondOrderEulerianPolynomial{0}(t)=1.
\]

\begin{theorem}[Second-order recurrence]\label{thm:second-eulerian-recurrence}
\begin{align}
 \SecondOrderEulerianNumber nk&=(2n-k-1)\SecondOrderEulerianNumber{n-1}{k-1}
 +(k+1)\SecondOrderEulerianNumber{n-1}{k},
 \label{eq:second-eulerian-recurrence}\\
 \SecondOrderEulerianPolynomial{n+1}(t)&=(2nt+1)\SecondOrderEulerianPolynomial{n}(t)-t(t-1)\SecondOrderEulerianPolynomial{n}'(t),
 \label{eq:second-eulerian-poly-recurrence}\\
 \SecondOrderEulerianPolynomial{n}(1)&=(2n-1)!!.
 \label{eq:second-eulerian-row-sum}
\end{align}
\end{theorem}
\begin{proof}
The two copies of the largest label $n$ must be adjacent.  Insert the pair $nn$
into one of the $2n-1$ gaps of a Stirling permutation of order $n-1$.  A descent
gap or the terminal gap preserves the number of descents, giving $k+1$ choices
from an object with $k$ descents.  Each of the remaining $2n-k-1$ gaps creates one
descent from an object with $k-1$ descents.  This is the first recurrence; summing
against $t^k$ gives the second.  At each stage there are $2n-1$ insertion gaps, so
the total count is $(2n-1)!!$.
\end{proof}

The differential recurrence has two useful equivalent forms:
\begin{align}
 (t-1)^{-2n-2}\SecondOrderEulerianPolynomial{n+1}(t)
 &=\left(t(1-t)^{-2n-1}\SecondOrderEulerianPolynomial{n}(t)\right)',
 \label{eq:second-eulerian-differential-form}\\
 u_{n+1}(t)&=\left(\frac{t}{1-t}u_n(t)\right)',
 \quad u_n(t)=(t-1)^{-2n}\SecondOrderEulerianPolynomial{n}(t),\quad u_0=1.
 \label{eq:second-eulerian-u}
\end{align}
They follow by the product rule from
\eqref{eq:second-eulerian-poly-recurrence}.

\begin{theorem}[Diagonal Stirling generating function]\label{thm:second-eulerian-stirling-gf}
For $n\ge1$,
\begin{equation}\label{eq:second-eulerian-stirling-gf}
 \sum_{m=0}^{\infty}\StirlingSecondKind{n+m}{m}t^m
 =\frac{t\SecondOrderEulerianPolynomial{n}(t)}{(1-t)^{2n+1}}.
\end{equation}
\end{theorem}
\begin{proof}
Let $F_n(t)$ denote the left side.  The second-kind recurrence gives
$(1-t)F_{n+1}=tF_n'$, with $F_0=(1-t)^{-1}$.  Hence
$F_1=t/(1-t)^3$.  If the displayed formula holds for $n$, differentiating
$t\SecondOrderEulerianPolynomial{n}/(1-t)^{2n+1}$ and using
\eqref{eq:second-eulerian-poly-recurrence} gives the formula for $n+1$.
\end{proof}


\section{Second-order Eulerian interpolation of first-kind diagonals}
The same second-order Eulerian row that controls
\eqref{eq:second-eulerian-stirling-gf} also gives a binomial-basis expansion for
the near-diagonals of the unsigned first-kind triangle.

\begin{theorem}[Second-order Eulerian diagonal formula]
\label{thm:second-eulerian-first-diagonal}
For $p,m\ge0$,
\begin{equation}
\label{eq:second-eulerian-first-diagonal}
 \UnsignedStirlingFirstKind m{m-p}
 =\sum_{j=0}^{p}\SecondOrderEulerianNumber pj\binom{m+j}{2p},
\end{equation}
where an out-of-range first-kind number is zero.  Equivalently, the polynomial
continuation in a variable $x$ satisfies
\[
 \UnsignedStirlingFirstKind{x}{x-p}
 =\sum_{j=0}^{p}\SecondOrderEulerianNumber pj\binom{x+j}{2p}.
\]
The first instances are
\begin{align*}
 \UnsignedStirlingFirstKind{x}{x-1}&=\binom{x}{2},\\
 \UnsignedStirlingFirstKind{x}{x-2}&=\binom{x}{4}+2\binom{x+1}{4},\\
 \UnsignedStirlingFirstKind{x}{x-3}&=\binom{x}{6}+8\binom{x+1}{6}
                    +6\binom{x+2}{6}.
\end{align*}
\end{theorem}
\begin{proof}
For fixed $p$ define
\[
 D_p(t)=\sum_{m\ge0}\UnsignedStirlingFirstKind m{m-p}t^m.
\]
Writing $a_{m,p}=\UnsignedStirlingFirstKind m{m-p}$, the first-kind recurrence becomes
$a_{m,p}=a_{m-1,p}+(m-1)a_{m-1,p-1}$.  Therefore
\begin{equation}
\label{eq:first-diagonal-ogf-recurrence}
 (1-t)D_p(t)=t^2D_{p-1}'(t),
 \qquad D_0(t)=\frac1{1-t}.
\end{equation}
Let
\[
 Q_p(t)=t^{2p}\SecondOrderEulerianPolynomial{p}(t^{-1}).
\]
The polynomial recurrence \eqref{eq:second-eulerian-poly-recurrence} is equivalent,
after the substitution $t\mapsto t^{-1}$, to
\[
 Q_p(t)=t^2\bigl((1-t)Q_{p-1}'(t)+(2p-1)Q_{p-1}(t)\bigr).
\]
It follows inductively from \eqref{eq:first-diagonal-ogf-recurrence} that
\begin{equation}
\label{eq:first-diagonal-second-eulerian-ogf}
 D_p(t)=\frac{Q_p(t)}{(1-t)^{2p+1}}
 =\frac{t^{2p}\SecondOrderEulerianPolynomial{p}(t^{-1})}{(1-t)^{2p+1}}.
\end{equation}
Since $\SecondOrderEulerianPolynomial{p}(u)=\sum_{j=0}^{p}\SecondOrderEulerianNumber pj u^j$, coefficient extraction gives
\[
 [t^m]D_p(t)
 =\sum_{j=0}^{p}\SecondOrderEulerianNumber pj
   [t^{m-2p+j}](1-t)^{-2p-1}
 =\sum_{j=0}^{p}\SecondOrderEulerianNumber pj\binom{m+j}{2p},
\]
which is \eqref{eq:second-eulerian-first-diagonal}.  Both sides are polynomials in
$m$ of degree at most $2p$, so the same identity gives the stated polynomial
continuation.
\end{proof}
\part{Differentiation, Composition, and Reversion}

\chapter{Fa\`a di Bruno's formula}
\index{Faà di Bruno formula@Fa\`a di Bruno formula}

\section{The set-partition form}
For a block $B\subseteq[n]$, write $|B|$ for its size.
\begin{theorem}[Fa\`a di Bruno, partition form]\label{thm:faa-partition}
For $n\ge0$ and sufficiently differentiable $f,g$,
\begin{equation}\label{eq:faa-partition}
 D^n(f\circ g)(x)=
 \sum_{\pi\in\SetPartitionOperator([n])}
 f^{(|\pi|)}(g(x))\prod_{B\in\pi}g^{(|B|)}(x).
\end{equation}
For $n=0$, the empty partition contributes $f(g(x))$.
\end{theorem}
\begin{proof}
Induct on $n$.  Differentiate a term indexed by a partition $\pi$ of $[n]$.
If the derivative hits $f^{(|\pi|)}(g)$, the chain rule contributes $g'$ and
creates the new singleton block $\{n+1\}$.  If it hits the factor associated with
a block $B$, it changes $g^{(|B|)}$ to $g^{(|B|+1)}$, which inserts $n+1$ into
$B$.  Every partition of $[n+1]$ arises uniquely in exactly one of these ways:
remove $n+1$, either deleting its singleton block or removing it from its
non-singleton block.  Hence all terms occur once.
\end{proof}

\section{Multiplicity and Bell-polynomial forms}
Grouping partitions by the number $m_j$ of blocks of size $j$ gives the classical
formula.
\begin{theorem}[Fa\`a di Bruno, multiplicity form]\label{thm:faa-multiplicity}
\begin{multline}\label{eq:faa-multiplicity}
 D^n f(g(x))=
 \sum_{\substack{m_1,\ldots,m_n\ge0\\
                  m_1+2m_2+\cdots+nm_n=n}}
 \frac{n!}{m_1!m_2!\cdots m_n!}\\
 {}\times f^{(m_1+\cdots+m_n)}(g(x))
 \prod_{j=1}^{n}\left(\frac{g^{(j)}(x)}{j!}\right)^{m_j}.
\end{multline}
Equivalently,
\begin{equation}\label{eq:faa-bell}
 D^n f(g(x))=
 \sum_{k=0}^{n}f^{(k)}(g(x))
 \ExponentialPartialBellPolynomial nk\bigl(g'(x),g''(x),\ldots,g^{(n-k+1)}(x)\bigr).
\end{equation}
\end{theorem}
\begin{proof}
For fixed multiplicities $(m_j)$, the number of set partitions of $[n]$ with
$m_j$ blocks of size $j$ is
$n!/\prod_j((j!)^{m_j}m_j!)$.  Substitution in
\eqref{eq:faa-partition} gives \eqref{eq:faa-multiplicity}.  Grouping further by
$k=\sum m_j$ and invoking \eqref{eq:partial-bell-definition} gives
\eqref{eq:faa-bell}.
\end{proof}

For $n=4$ this reads
\begin{align}
 (f\circ g)^{(4)}={}&f^{(4)}(g)(g')^4
 +6f^{(3)}(g)(g')^2g''+3f''(g)(g'')^2\notag\\
 &+4f''(g)g'g^{(3)}+f'(g)g^{(4)}.
 \label{eq:faa-n4}
\end{align}
The coefficients $1,6,3,4,1$ count partitions of types
$1+1+1+1$, $2+1+1$, $2+2$, $3+1$, and $4$ respectively.  This also proves the
expanded example in the archive without four separate differentiations.

\section{The exponential special case}
Because every derivative of $\EulerE^u$ equals $\EulerE^u$, summing the partial Bell
polynomials in \eqref{eq:faa-bell} yields
\begin{equation}\label{eq:exp-faa}
 D^n \EulerE^{g(x)}=\EulerE^{g(x)}
 \ExponentialCompleteBellPolynomial n\bigl(g'(x),g''(x),\ldots,g^{(n)}(x)\bigr).
\end{equation}
When $g$ is a cumulant generating function, this is precisely the
moment--cumulant relation of \cref{thm:moment-cumulant}.

\section{Mixed partial derivatives}
For $B\subseteq[n]$, write
$\partial_B=\prod_{j\in B}\partial/\partial x_j$.
\begin{theorem}[Multivariate Fa\`a di Bruno for distinct variables]
\label{thm:faa-multivariate}
If $y=g(x_1,\ldots,x_n)$, then
\begin{equation}\label{eq:faa-multivariate}
 \frac{\partial^n}{\partial x_1\cdots\partial x_n}f(y)
 =\sum_{\pi\in\SetPartitionOperator([n])}f^{(|\pi|)}(y)
   \prod_{B\in\pi}\partial_B y.
\end{equation}
\end{theorem}
\begin{proof}
The proof of \cref{thm:faa-partition} applies verbatim, except that the new
derivative is $\partial/\partial x_{n+1}$.  It either creates a singleton block or
joins the new index to one existing block.
\end{proof}

For three variables, \eqref{eq:faa-multivariate} becomes
\begin{align*}
 \partial_1\partial_2\partial_3f(y)
 ={}&f'(y)\partial_{123}y\\
 &+f''(y)\bigl((\partial_1y)(\partial_{23}y)
 +(\partial_2y)(\partial_{13}y)+(\partial_3y)(\partial_{12}y)\bigr)\\
 &+f^{(3)}(y)(\partial_1y)(\partial_2y)(\partial_3y).
\end{align*}
For repeated variables or vector-valued inner maps, one labels derivative slots
first and then contracts tensor indices; the same partition lattice remains the
combinatorial skeleton.

\chapter{Composition of ordinary and exponential series}

\section{Ordinary formal power series}
Let
\[
 F(u)=\sum_{k\ge0}a_ku^k,
 \qquad
 G(x)=\sum_{j\ge1}b_jx^j,
 \qquad
 F(G(x))=\sum_{n\ge0}c_nx^n.
\]
\begin{theorem}[Ordinary composition]\label{thm:ordinary-composition}
For $n\ge1$,
\begin{align}
 c_n&=\sum_{k=1}^{n}a_k
 \sum_{\substack{i_1+\cdots+i_k=n\\i_r\ge1}}
 b_{i_1}\cdots b_{i_k},
 \label{eq:ordinary-composition-compositions}\\
 &=\sum_{k=1}^{n}a_k
 \sum_{\substack{\pi_1+\cdots+\pi_n=k\\
                   \pi_1+2\pi_2+\cdots+n\pi_n=n}}
 \binom{k}{\pi_1,\ldots,\pi_n}
 b_1^{\pi_1}\cdots b_n^{\pi_n}
 \label{eq:ordinary-composition-multiplicities}\\
 &=\sum_{k=1}^{n}a_k\OrdinaryPartialBellPolynomial nk(b_1,\ldots,b_{n-k+1}),
 \label{eq:ordinary-composition-bell}
\end{align}
and $c_0=a_0$.
\end{theorem}
\begin{proof}
The coefficient of $x^n$ in $G(x)^k$ is obtained by choosing an ordered
composition $i_1+\cdots+i_k=n$, proving the first formula.  Collect equal part
sizes; the number of orderings with multiplicities $\pi_j$ is the multinomial
coefficient, proving the second.  The third is the definition of the ordinary
Bell polynomial.
\end{proof}

Two useful specializations are immediate.  If $F(u)=\EulerE^u$ and the series are
exponentially normalized, one obtains the exponential formula.  If
$A(x)=1+\sum_{n\ge1}a_nx^n$, then
\begin{equation}\label{eq:reciprocal-ordinary-bell}
 [x^n]\frac1{A(x)}
 =\sum_{k=1}^{n}(-1)^k\OrdinaryPartialBellPolynomial nk(a_1,\ldots,a_{n-k+1})
 \qquad(n\ge1),
\end{equation}
by taking $F(u)=1/(1-u)$ and $G(x)=-\sum_{n\ge1}a_nx^n$.

\section{Exponential formal power series}
Let
\[
 F(x)=\sum_{j\ge1}a_j\frac{x^j}{j!},\qquad
 G(u)=\sum_{k\ge0}b_k\frac{u^k}{k!},\qquad
 G(F(x))=\sum_{n\ge0}c_n\frac{x^n}{n!}.
\]
\begin{theorem}[Exponential composition]\label{thm:exponential-composition}
\begin{equation}\label{eq:exponential-composition}
 c_n=\sum_{k=0}^{n}b_k\ExponentialPartialBellPolynomial nk(a_1,\ldots,a_{n-k+1}).
\end{equation}
Equivalently,
\begin{equation}\label{eq:exponential-composition-partitions}
 c_n=\sum_{\pi\in\SetPartitionOperator([n])}b_{|\pi|}\prod_{B\in\pi}a_{|B|}.
\end{equation}
\end{theorem}
\begin{proof}
Substitute $F(x)$ into $G$ and apply
\eqref{eq:partial-bell-egf}.  The set-partition form is
\cref{thm:bell-poly-partitions} with an additional weight $b_k$ for the number of
blocks.
\end{proof}

Taking $b_k=1$ gives the complete Bell-polynomial exponential identity
\begin{equation}\label{eq:exp-complete-bell-again}
 \exp\!\left(\sum_{j\ge1}a_j\frac{x^j}{j!}\right)
 =\sum_{n\ge0}\ExponentialCompleteBellPolynomial n(a_1,\ldots,a_n)\frac{x^n}{n!}.
\end{equation}
If the series converge, \eqref{eq:exponential-composition} is exactly
\eqref{eq:faa-bell} evaluated at the expansion point; formally, no convergence
notion is needed.

\chapter{Lagrange inversion and Lagrange--Bürmann}
\index{Lagrange inversion}

\section{Formal theorem}
Let $R$ be a commutative $\RationalNumbers$-algebra and let $\phi(w)\in R[[w]]$ have invertible
constant term.  There is a unique $g(z)\in zR[[z]]$ satisfying
\begin{equation}\label{eq:lagrange-functional}
 g(z)=z\phi(g(z)).
\end{equation}
Uniqueness follows recursively because the coefficient of $z^n$ on the right
depends only on coefficients of $g$ below degree $n$.

\begin{theorem}[Lagrange--Bürmann]\label{thm:lagrange-burmann}
For $n\ge1$ and every formal series $H$,
\begin{align}
 \CoefficientExtraction zn H(g(z))
 &=\frac1n\CoefficientExtraction w{n-1}H'(w)\phi(w)^n,
 \label{eq:lagrange-burmann}\\
 &=\CoefficientExtraction wn H(w)\phi(w)^{n-1}\bigl(\phi(w)-w\phi'(w)\bigr).
 \label{eq:lagrange-burmann-alt}
\end{align}
In particular,
\begin{equation}\label{eq:lagrange-basic}
 \CoefficientExtraction zn g(z)=\frac1n\CoefficientExtraction w{n-1}\phi(w)^n.
\end{equation}
\end{theorem}
\begin{proof}
Put $f(w)=w/\phi(w)$, so that $f(g(z))=z$.  Formal residues and the change of
variables theorem \cref{thm:res-subst} give
\begin{align*}
 \CoefficientExtraction zn H(g(z))
 &=\ResidueOperator_z H(g(z))z^{-n-1}\Differential z\\
 &=\ResidueOperator_w H(w)f(w)^{-n-1}f'(w)\Differential w\\
 &=-\frac1n\ResidueOperator_w H(w)(f(w)^{-n})'\Differential w\\
 &=\frac1n\ResidueOperator_w H'(w)f(w)^{-n}\Differential w,
\end{align*}
where the residue of a derivative is zero.  Since
$f(w)^{-n}=w^{-n}\phi(w)^n$, this is
\eqref{eq:lagrange-burmann}.  Applying the same zero-residue rule to
$H(w)\phi(w)^nw^{-n}$ gives
\[
 \frac1n\ResidueOperator H'\phi^nw^{-n}
 =\ResidueOperator H\phi^nw^{-n-1}-\ResidueOperator H\phi^{n-1}\phi'w^{-n},
\]
which is \eqref{eq:lagrange-burmann-alt}.  Take $H(w)=w$ for
\eqref{eq:lagrange-basic}.
\end{proof}

\section{Analytic inverse functions}
\begin{theorem}[Analytic Lagrange inversion]\label{thm:analytic-lagrange}
Let $f$ be analytic near $a$, let $b=f(a)$, and suppose $f'(a)\ne0$.  Its local
inverse $g$ is analytic near $b$ and has
\begin{equation}\label{eq:analytic-lagrange-series}
 g(z)=a+\sum_{n\ge1}g_n\frac{(z-b)^n}{n!},
\end{equation}
where
\begin{equation}\label{eq:analytic-lagrange-coeff}
 g_n=\left.\frac{\Differential^{n-1}}{\Differential w^{n-1}}
 \left(\frac{w-a}{f(w)-f(a)}\right)^n\right|_{w=a}.
\end{equation}
The series has a positive radius of convergence.
\end{theorem}
\begin{proof}
The analytic inverse-function theorem gives an analytic local inverse and therefore
a convergent Taylor series.  Put $u=w-a$, $\LagrangeShiftCoordinate=z-b$, and
\[
 \phi(u)=\frac{u}{f(a+u)-f(a)};
\]
its constant term is $1/f'(a)$.  The inverse displacement satisfies
$u=\LagrangeShiftCoordinate\phi(u)$.  Formula \eqref{eq:lagrange-basic}, multiplied by $n!$, is
exactly \eqref{eq:analytic-lagrange-coeff}.
\end{proof}

If $f'(a)=0$ but $f(w)-f(a)=c(w-a)^q+O((w-a)^{q+1})$ with $c\ne0$, factor
$f(w)-f(a)=(w-a)^qU(w)$ with $U(a)=c$.  Choosing one of the $q$ values of
$c^{-1/q}$ reduces the equation to ordinary inversion in the variable
$(z-f(a))^{1/q}$.  Thus the inverse has $q$ local Puiseux branches.  This is the
precise sense in which Lagrange inversion extends to a critical point; a
single-valued Taylor inverse cannot exist there.

\section{Inverse coefficients in Bell-polynomial form}
Suppose
\[
 f(w)=\sum_{j\ge1}f_j\frac{w^j}{j!},\qquad f_1\ne0,
 \qquad
 g(z)=f^{\langle-1\rangle}(z)=\sum_{n\ge1}g_n\frac{z^n}{n!}.
\]
Define
\begin{equation}\label{eq:fhat}
 \NormalizedReversionCoefficient{f}{j}=\frac{f_{j+1}}{(j+1)f_1}.
\end{equation}
\begin{theorem}[Explicit inverse coefficients]\label{thm:inverse-bell-coeff}
One has $g_1=f_1^{-1}$ and, for $n\ge2$,
\begin{equation}\label{eq:inverse-bell-coeff}
 g_n=\frac1{f_1^n}\sum_{k=1}^{n-1}(-1)^k\RisingFactorial nk
 \ExponentialPartialBellPolynomial{n-1}{k}(\NormalizedReversionCoefficient{f}{1},\ldots,\NormalizedReversionCoefficient{f}{n-k}).
\end{equation}
\end{theorem}
\begin{proof}
Write
\[
 \frac{f(w)}{f_1w}=1+U(w),
 \qquad U(w)=\sum_{j\ge1}\NormalizedReversionCoefficient{f}{j}\frac{w^j}{j!}.
\]
By \eqref{eq:lagrange-basic},
$g_n=(n-1)!\CoefficientExtraction w{n-1}(w/f(w))^n$.  Now
\[
 (1+U)^{-n}=\sum_{k\ge0}(-1)^k\binom{n+k-1}{k}U^k
 =\sum_{k\ge0}(-1)^k\frac{\RisingFactorial nk}{k!}U^k.
\]
Apply \eqref{eq:partial-bell-egf} and extract degree $n-1$; the $k=0$ term
vanishes when $n\ge2$.
\end{proof}

\section{The associahedron interpretation, with normalized notation}
The formula in the archive mixes exponential coefficients with an ordinary-series
face formula.  The exact geometric statement is the following.

Let
\begin{equation}\label{eq:ordinary-series-associa}
 A(x)=x+\sum_{n\ge1}a_nx^{n+1},\qquad
 B(x)=A^{\langle-1\rangle}(x)=x+\sum_{n\ge1}b_nx^{n+1}.
\end{equation}
Use the indexing in which the associahedron $\AssociahedronByDimension{n-1}$ has dimension $n-1$ and its
faces are plane rooted trees with $n+1$ leaves and internal arities at least $2$.
If an internal vertex has arity $i+1$, give it weight $a_i$; let $a_F$ be the
product of all vertex weights.  Equivalently, if
$F\cong \AssociahedronByDimension{i_1-1}\times\cdots\times \AssociahedronByDimension{i_m-1}$, then
$a_F=a_{i_1}\cdots a_{i_m}$.

\begin{theorem}[Associahedral inversion]\label{thm:associahedral-inversion}
\begin{equation}\label{eq:associahedral-inversion}
 b_n=\sum_{F\text{ a face of }\AssociahedronByDimension{n-1}}(-1)^{n-\dim F}a_F.
\end{equation}
\end{theorem}
\begin{proof}
The equation $A(B)=x$ is
\[
 B=x-\sum_{i\ge1}a_iB^{i+1}.
\]
Iterating it recursively expands $b_n$ as a sum over plane rooted trees with
$n+1$ leaves: every internal vertex of arity $i+1$ contributes $-a_i$.
A face $F$ of $\AssociahedronByDimension{n-1}$ is exactly such a tree after contracting a chosen set of
internal edges.  If it has $m$ internal vertices, then
$\dim F=n-m$, so its sign is $(-1)^m=(-1)^{n-\dim F}$ and its weight is $a_F$.
Thus the tree expansion is the alternating face sum.
\end{proof}

For example,
\begin{align*}
 b_1&=-a_1,\\
 b_2&=-a_2+2a_1^2,\\
 b_3&=-a_3+5a_2a_1-5a_1^3,\\
 b_4&=-a_4+6a_3a_1+3a_2^2-21a_2a_1^2+14a_1^4.
\end{align*}
For $\AssociahedronByDimension{3}$, the coefficients $1,6,3,21,14$ count respectively the whole
3-polytope, its six pentagonal and three square facets, its edges, and its
vertices.

\chapter{Algebraic and transcendental examples}

\section{The Fuss--Catalan algebraic equation}
Let $p\ge2$ be an integer and consider
\begin{equation}\label{eq:fuss-equation}
 x^p-x+z=0,
 \qquad x(0)=0.
\end{equation}
It is equivalent to
$x=z\phi(x)$ with $\phi(x)=(1-x^{p-1})^{-1}$.

\begin{theorem}\label{thm:fuss-series}
The distinguished solution is
\begin{equation}\label{eq:fuss-series}
 x(z)=\sum_{k=0}^{\infty}\binom{pk}{k}
       \frac{z^{(p-1)k+1}}{(p-1)k+1}.
\end{equation}
Its radius of convergence is
\begin{equation}\label{eq:fuss-radius}
 R_p=(p-1)p^{-p/(p-1)},
\end{equation}
and the series converges absolutely also on $|z|=R_p$.
\end{theorem}
\begin{proof}
By \eqref{eq:lagrange-basic},
\[
 \CoefficientExtraction znx(z)=\frac1n\CoefficientExtraction w{n-1}(1-w^{p-1})^{-n}.
\]
This vanishes unless $n-1=(p-1)k$.  In that case the coefficient is
$\frac1n\binom{n+k-1}{k}=\binom{pk}{k}/((p-1)k+1)$, proving the series.

The inverse map is $z=f(x)=x-x^p$.  Its nearest critical points satisfy
$f'(x)=1-px^{p-1}=0$ and all have modulus $p^{-1/(p-1)}$; their critical values
have modulus $(p-1)p^{-p/(p-1)}$.  No local inverse can pass analytically through
a critical value, while the inverse-function theorem continues the branch up to
one, so this is the radius.  Stirling's formula gives coefficients of order
$Ck^{-3/2}R_p^{-((p-1)k+1)}$, which proves absolute convergence on the boundary.
\end{proof}

\section{Lambert's \texorpdfstring{$W$}{W} function at the origin}
The Lambert function satisfies $W(z)\EulerE^{W(z)}=z$.
\begin{theorem}\label{thm:lambert-W-zero}
On its principal branch near $0$,
\begin{equation}\label{eq:lambert-W-zero}
 W(z)=\sum_{n=1}^{\infty}(-n)^{n-1}\frac{z^n}{n!}
 =z-z^2+\frac32z^3-\frac83z^4+O(z^5),
\end{equation}
with radius of convergence $\EulerE^{-1}$.
\end{theorem}
\begin{proof}
Here $f(w)=we^w$, so $w/f(w)=\EulerE^{-w}$.  Formula
\eqref{eq:analytic-lagrange-coeff} gives
$g_n=D^{n-1}\EulerE^{-nw}|_{w=0}=(-n)^{n-1}$.  The inverse map has its nearest critical
point at $w=-1$, where $f(-1)=-\EulerE^{-1}$; this square-root branch point determines
the radius.
\end{proof}

\section{A rapidly useful expansion around \texorpdfstring{$W(e)=1$}{W(e)=1}}
Put $\xi=\log z-1$ and $u=W(z)-1$.  Then
\begin{equation}\label{eq:lambert-near-e-implicit}
 u+\log(1+u)=\xi.
\end{equation}
Reversion gives
\begin{equation}\label{eq:lambert-near-e}
 W(\EulerE^{1+\xi})=1+\frac{\xi}{2}+\frac{\xi^2}{16}
 -\frac{\xi^3}{192}-\frac{\xi^4}{3072}
 +\frac{13\xi^5}{61440}+O(\xi^6).
\end{equation}
\begin{proof}
The derivative at $u=0$ of the left side of
\eqref{eq:lambert-near-e-implicit} is $2$, so it has an ordinary inverse.
Substitute $u=c_1\xi+\cdots+c_5\xi^5$ into
$u+\log(1+u)=\xi$ and compare successive powers.  The triangular equations give
$c_1=1/2,c_2=1/16,c_3=-1/192,c_4=-1/3072,c_5=13/61440$.
\end{proof}
The closest critical points of $u+\log(1+u)$ satisfy
$1+(1+u)^{-1}=0$, hence $u=-2$; using the two adjacent logarithm values gives
images $-2\pm \ImaginaryUnit\pi$.  Therefore the radius in the $\xi$-plane is
$\sqrt{4+\pi^2}$.  For positive real $z$, the expansion converges whenever
$|\log z-1|<\sqrt{4+\pi^2}$, approximately
$0.0655<z<112.63$.  Taking $\xi=-1$ yields a convergent evaluation of
$W(1)=0.567143\ldots$.

\section{Binary trees and Catalan numbers}
Let $\CatalanNumber{n}$ count rooted ordered full binary trees with $n$ internal vertices, and
let $B(z)=\sum_{n\ge0}\CatalanNumber{n}z^n$.  Removing the root gives
\begin{equation}\label{eq:binary-tree-functional}
 B(z)=1+zB(z)^2.
\end{equation}
Put $C(z)=B(z)-1$; then $C=z(1+C)^2$.  Lagrange inversion gives
\begin{equation}\label{eq:catalan-lagrange}
 \CatalanNumber{n}=\CoefficientExtraction znC(z)=\frac1n\CoefficientExtraction w{n-1}(1+w)^{2n}
 =\frac1n\binom{2n}{n-1}=\frac1{n+1}\binom{2n}{n}.
\end{equation}
This simultaneously proves the functional equation, its unique formal solution,
and the Catalan closed form.  The symbol $\CatalanNumber{n}$ is used here to avoid confusing
these binary-tree numbers with Bell numbers.

\chapter{Lagrange's reversion theorem in shifted form}

\section{The two-parameter implicit equation}
Let $v=v(x,y)$ be the unique formal solution of
\begin{equation}\label{eq:shifted-reversion-equation}
 v=x+yf(v),
 \qquad v(x,0)=x.
\end{equation}

\begin{theorem}[Lagrange reversion]\label{thm:shifted-lagrange-reversion}
For every formal or analytic $g$,
\begin{equation}\label{eq:shifted-lagrange-reversion}
 g(v)=g(x)+\sum_{k=1}^{\infty}\frac{y^k}{k!}
 \left(\frac{\partial}{\partial x}\right)^{k-1}
 \bigl(f(x)^kg'(x)\bigr).
\end{equation}
In particular,
\begin{equation}\label{eq:shifted-lagrange-v}
 v=x+\sum_{k=1}^{\infty}\frac{y^k}{k!}
 \left(\frac{\partial}{\partial x}\right)^{k-1}f(x)^k.
\end{equation}
\end{theorem}
\begin{proof}
Put $u=v-x$.  Then $u=y\phi(u)$ with $\phi(u)=f(x+u)$, treating $x$ as a
coefficient parameter.  Apply \eqref{eq:lagrange-burmann} to
$H(u)=g(x+u)$:
\[
 [y^k]g(v)=\frac1k[u^{k-1}]g'(x+u)f(x+u)^k
 =\frac1{k!}D_x^{k-1}(g'(x)f(x)^k).
\]
Summation proves the first formula; take $g(z)=z$ for the second.
\end{proof}

\section{The delta-function proof decoded}
The archive also gives a distributional derivation.  Assume temporarily that
$x,y\in\RealNumbers$, the equation $z-x-yf(z)=0$ has one simple root $v$, and
$1-yf'(v)>0$.  The elementary change-of-variables rule for the Dirac delta gives
\begin{equation}\label{eq:delta-root}
 g(v)=\int_{\RealNumbers}\delta(yf(z)-z+x)g(z)(1-yf'(z))\,\Differential z.
\end{equation}
Using
$\delta(s)=\int_{\RealNumbers}\EulerE^{iks}\Differential k/(2\pi)$ and expanding only in the formal
parameter $y$,
\begin{align}
 g(v)
 &=\sum_{n\ge0}D_x^n\left[
 \frac{y^nf(x)^ng(x)}{n!}(1-yf'(x))\right]\notag\\
 &=\sum_{n\ge0}D_x^n\left[
 \frac{y^nf(x)^ng(x)}{n!}
 -\frac{y^{n+1}}{(n+1)!}\bigl((gf^{n+1})'-g'f^{n+1}\bigr)
 \right].
 \label{eq:delta-reversion-middle}
\end{align}
The first term of level $n+1$ cancels the total derivative in the second term of
level $n$; what remains is exactly
\eqref{eq:shifted-lagrange-reversion}.  Thus the Fourier--delta calculation is a
compact analytic avatar of the formal residue proof.  Formula
\eqref{eq:delta-root} requires an absolute Jacobian in general; the displayed sign
is valid in the stated orientation-preserving neighbourhood.

\chapter{Laplace-type asymptotic integrals}
\index{Laplace asymptotics}

\section{Endpoint expansion}
Consider
\begin{equation}\label{eq:laplace-integral}
 I(\lambda)=\int_a^b \EulerE^{-\lambda f(x)}g(x)\,\Differential x,
 \qquad \lambda\to+\infty,
\end{equation}
where the endpoint $a$ is the unique relevant minimum and, with $t=x-a$,
\begin{align}
 f(a+t)&\sim f(a)+\sum_{k\ge0}a_kt^{k+\alpha},
 &a_0>0,\quad\alpha>0,
 \label{eq:laplace-f-expansion}\\
 g(a+t)&\sim\sum_{k\ge0}b_kt^{k+\beta-1},
 &\beta>0.
 \label{eq:laplace-g-expansion}
\end{align}
Let $\OrdinaryPartialBellPolynomial{k}{j}$ denote the ordinary Bell polynomial.

\begin{theorem}[Laplace--Erd\'elyi coefficient formula]
\label{thm:laplace-bell}
Under the standard differentiability and remainder hypotheses that permit Watson's
lemma,
\begin{equation}\label{eq:laplace-asymptotic}
 I(\lambda)\sim \EulerE^{-\lambda f(a)}
 \sum_{n=0}^{\infty}
 \Gamma\!\left(\frac{n+\beta}{\alpha}\right)
 \frac{c_n}{\lambda^{(n+\beta)/\alpha}},
\end{equation}
where
\begin{equation}\label{eq:laplace-cn}
 c_n=\frac1{\alpha a_0^{(n+\beta)/\alpha}}
 \sum_{k=0}^{n}b_{n-k}
 \sum_{j=0}^{k}
 \binom{-\frac{n+\beta}{\alpha}}{j}
 \frac{\OrdinaryPartialBellPolynomial{k}{j}(a_1,\ldots,a_{k-j+1})}{a_0^j}.
\end{equation}
\end{theorem}
\begin{proof}
Write
$f(a+t)-f(a)=t^\alpha A(t)$ with
$A(t)=a_0+a_1t+\cdots$, and set
$s=(f(a+t)-f(a))^{1/\alpha}=tA(t)^{1/\alpha}$.  This has an inverse $t=t(s)$.
Expand
\[
 g(a+t(s))t'(s)=\sum_{n\ge0}d_ns^{n+\beta-1}.
\]
A generalized Lagrange--Bürmann calculation, first for integral $\beta$ and then
by formal continuation in the exponent, gives
\[
 d_n=\sum_{k=0}^{n}b_{n-k}[t^k]A(t)^{-(n+\beta)/\alpha}.
\]
Indeed, apply \eqref{eq:lagrange-burmann} to an antiderivative of
$t^{m+\beta-1}$ and differentiate back after substitution.

Put $\rho=(n+\beta)/\alpha$.  The ordinary Bell expansion yields
\begin{align*}
 [t^k]A(t)^{-\rho}
 &=a_0^{-\rho}[t^k]
 \sum_{j\ge0}\binom{-\rho}{j}
 \left(\sum_{r\ge1}\frac{a_r}{a_0}t^r\right)^j\\
 &=a_0^{-\rho}\sum_{j=0}^{k}\binom{-\rho}{j}a_0^{-j}
 \OrdinaryPartialBellPolynomial{k}{j}(a_1,\ldots,a_{k-j+1}).
\end{align*}
Thus $c_n=d_n/\alpha$ is \eqref{eq:laplace-cn}.  Finally
\[
 \int_0^\infty \EulerE^{-\lambda s^\alpha}s^{n+\beta-1}\Differential s
 =\frac1\alpha\Gamma\!\left(\frac{n+\beta}{\alpha}\right)
 \lambda^{-(n+\beta)/\alpha},
\]
and Watson's lemma controls the replacement of the finite endpoint interval by
the asymptotic expansion.
\end{proof}

\begin{remark}
The appearance of ordinary rather than exponential Bell polynomials is forced by
the ordinary Taylor expansion of $A(t)^{-\rho}$.  Lagrange inversion enters in the
change from $t$ to the natural Laplace coordinate $s$; the gamma factor enters only
in the final monomial integral.
\end{remark}

\part{Incidence Algebra, Umbral Calculus, and Multivariate Inversion}

\chapter{Binomial inversion, partition lattices, and Sheffer--Riordan calculus}
\label{ch:merged-umbral}

The concrete triangular arrays developed earlier are instances of a larger
calculus.  This chapter collects the genuinely nonduplicate material from the
five drafts: incidence-algebra inversion, the partition-lattice Möbius
function, Sheffer sequences, and exponential Riordan arrays.  These viewpoints
explain why changes of basis, labelled decomposition, and formal substitution
repeatedly lead to the same matrices.

\section{Binomial inversion}

\begin{theorem}[Binomial inversion]
\label{thm:merged-binomial-inversion}
For sequences $(a_n)$ and $(b_n)$ in a commutative $\RationalNumbers$-algebra, the formulas
\begin{align}
 b_n&=\sum_{k=0}^{n}\binom nk a_k,                                  \label{eq:merged-binomial-forward}\\
 a_n&=\sum_{k=0}^{n}(-1)^{n-k}\binom nk b_k                         \label{eq:merged-binomial-backward}
\end{align}
are equivalent.  In exponential generating functions,
\begin{equation}
 B(z)=\EulerE^zA(z),\qquad A(z)=\EulerE^{-z}B(z).                               \label{eq:merged-binomial-egf}
\end{equation}
\end{theorem}

\begin{proof}
Multiplication of EGFs gives
\[
 \EulerE^zA(z)=\sum_{n\geq0}\left(\sum_{k=0}^{n}\binom nk a_k\right)
             \frac{z^n}{n!},
\]
so \eqref{eq:merged-binomial-forward} is equivalent to the first identity in
\eqref{eq:merged-binomial-egf}.  Multiplication by $\EulerE^{-z}$ and coefficient
extraction gives \eqref{eq:merged-binomial-backward}.  Directly, substitution
of the forward formula into the backward one makes the coefficient of $a_j$
\[
 \sum_{k=j}^{n}(-1)^{n-k}\binom nk\binom kj
 =\binom nj\sum_{r=0}^{n-j}(-1)^{n-j-r}\binom{n-j}{r}
 =\delta_{nj},
\]
which independently verifies the inverse matrices.
\end{proof}

\section{The partition lattice and cumulants}

Let $\Pi_n$ be the lattice of partitions of $[n]$, ordered by refinement,
with bottom $\PartitionLatticeMinimum$ and top $\PartitionLatticeMaximum$.

\begin{theorem}[Möbius function of $\Pi_n$]
\label{thm:merged-partition-mobius}
If $\sigma\leq\pi$ and $m_C$ is the number of blocks of $\sigma$ contained in
a block $C$ of $\pi$, then
\begin{equation}
 \mu(\sigma,\pi)=\prod_{C\in\pi}(-1)^{m_C-1}(m_C-1)!.              \label{eq:merged-partition-mobius}
\end{equation}
In particular,
\begin{align}
 \mu(\PartitionLatticeMinimum,\pi)&=\prod_{B\in\pi}(-1)^{|B|-1}(|B|-1)!,         \label{eq:merged-partition-mobius-bottom}\\
 \mu(\pi,\PartitionLatticeMaximum)&=(-1)^{|\pi|-1}(|\pi|-1)!.                    \label{eq:merged-partition-mobius-top}
\end{align}
\end{theorem}

\begin{proof}
The interval $[\sigma,\pi]$ is the direct product, over $C\in\pi$, of
partition lattices on the $m_C$ blocks of $\sigma$ inside $C$.  Möbius
functions multiply on direct products, so it remains to compute
$u_m=\mu_{\Pi_m}(\PartitionLatticeMinimum,\PartitionLatticeMaximum)$.  The defining relation gives, for
$m>1$,
\[
 0=\sum_{\rho\in\Pi_m}\mu(\PartitionLatticeMinimum,\rho)
  =\sum_{\rho\in\Pi_m}\prod_{B\in\rho}u_{|B|}.
\]
Let $U(z)=\sum_{m\geq1}u_mz^m/m!$.  By the labelled exponential formula,
\[
 \exp U(z)=1+z,
\]
because the inner sum is $1$ for $m=0,1$ and $0$ thereafter.  Hence
$U(z)=\log(1+z)$ and $u_m=(-1)^{m-1}(m-1)!$.  Product decomposition now gives
the general formula.
\end{proof}

\begin{theorem}[Moment--cumulant inversion]
\label{thm:merged-moment-cumulant}
The relations
\begin{align}
 m_n&=\sum_{\pi\in\Pi_n}\prod_{B\in\pi}\kappa_{|B|},               \label{eq:merged-moment-cumulant-forward}\\
 \kappa_n&=\sum_{\pi\in\Pi_n}(-1)^{|\pi|-1}(|\pi|-1)!
                 \prod_{B\in\pi}m_{|B|}                            \label{eq:merged-moment-cumulant-backward}
\end{align}
are inverse.  Equivalently, if
$M(z)=1+\sum_{n\geq1}m_nz^n/n!$ and
$K(z)=\sum_{n\geq1}\kappa_nz^n/n!$, then $M=\EulerE^K$ and $K=\log M$.
\end{theorem}

\begin{proof}
The forward relation is the exponential formula for a set of blocks weighted
by cumulants.  Möbius inversion on $\Pi_n$, using
\eqref{eq:merged-partition-mobius-top}, gives the backward relation.  Grouping
partitions by block sizes is exactly the coefficient expansion of $\EulerE^K$ and
$\log M$.
\end{proof}

\section{Sheffer sequences}

\begin{definition}
A polynomial sequence $(s_n(x))_{n\geq0}$ is Sheffer for $(A,B)$ when
\begin{equation}
 \sum_{n\geq0}s_n(x)\frac{t^n}{n!}=A(t)\EulerE^{xB(t)},                   \label{eq:merged-sheffer-egf}
\end{equation}
where $A(0)\ne0$, $B(0)=0$, and $B'(0)\ne0$.  It is of binomial type when
$A=1$, and Appell when $B(t)=t$.
\end{definition}

\begin{theorem}[Addition and lowering]
\label{thm:merged-sheffer}
Let $p_n$ be the binomial-type sequence with EGF $\EulerE^{yB(t)}$.  Then
\begin{equation}
 s_n(x+y)=\sum_{k=0}^{n}\binom nk s_k(x)p_{n-k}(y).                 \label{eq:merged-sheffer-addition}
\end{equation}
If $\overline B$ is the compositional inverse of $B$ and
$Q=\overline B(D_x)$, then
\begin{equation}
 Qs_n(x)=ns_{n-1}(x).                                               \label{eq:merged-sheffer-lowering}
\end{equation}
\end{theorem}

\begin{proof}
Multiplying the two EGFs at $x$ and $y$ proves the addition law.  For every
formal series $F$,
$F(D_x)\EulerE^{xB(t)}=F(B(t))\EulerE^{xB(t)}$.  Taking $F=\overline B$ gives the
eigenvalue $t$; multiplication of the EGF by $t$ and coefficient comparison
gives the lowering law.
\end{proof}

\begin{proposition}[Abel polynomials]
\label{prop:merged-abel}
Let $T(t)$ solve $T=t \EulerE^{-aT}$.  Then
\begin{equation}
 \EulerE^{xT(t)}=1+\sum_{n\geq1}x(x-na)^{n-1}\frac{t^n}{n!}.              \label{eq:merged-abel-egf}
\end{equation}
Thus $x(x-na)^{n-1}$ is a sequence of binomial type.
\end{proposition}

\begin{proof}
Lagrange--Bürmann gives
\[
 [t^n]\EulerE^{xT(t)}=\frac{x}{n}[w^{n-1}]\EulerE^{xw}\EulerE^{-anw}
 =\frac{x(x-na)^{n-1}}{n!}.
\]
An EGF of the form $\EulerE^{xB(t)}$ satisfies the binomial identity by
multiplication at $x$ and $y$.
\end{proof}

\section{Exponential Riordan arrays}

\begin{definition}
For $g(0)\ne0$, $f(0)=0$, and $f'(0)\ne0$, define
\begin{equation}
 R_{n,k}=\frac{n!}{k!}[t^n]g(t)f(t)^k.                              \label{eq:merged-riordan-entry}
\end{equation}
The resulting lower-triangular matrix is denoted $[g,f]$.
\end{definition}

\begin{theorem}[Riordan action, product, and inverse]
\label{thm:merged-riordan}
If $(a_k)$ has EGF $A(t)$, then $[g,f]a$ has EGF
\begin{equation}
 g(t)A(f(t)).                                                        \label{eq:merged-riordan-action}
\end{equation}
Consequently,
\begin{align}
 [g,f][h,\ell]&=[g\,(h\circ f),\ell\circ f],                          \label{eq:merged-riordan-product}\\
 [g,f]^{-1}&=\left[\frac1{g\circ\overline f},\overline f\right].   \label{eq:merged-riordan-inverse}
\end{align}
\end{theorem}

\begin{proof}
By definition,
\[
 \sum_n\left(\sum_kR_{n,k}a_k\right)\frac{t^n}{n!}
 =g(t)\sum_k a_k\frac{f(t)^k}{k!}=g(t)A(f(t)).
\]
Applying this action twice gives the product law, and substituting
$\overline f$ gives the inverse.
\end{proof}

\begin{example}[Stirling matrices]
The second-kind matrix is $[1,\EulerE^t-1]$, while the signed first-kind matrix is
$[1,\log(1+t)]$.  Since the two substitution series are compositional
inverses, Stirling inversion is the Riordan inverse law.  The Bell-polynomial
matrix is obtained by replacing $\EulerE^t-1$ with
$\sum_{j\geq1}x_jt^j/j!$.
\end{example}

\chapter{Complementary Bell numbers and signed partition calculus}
\label{ch:merged-complementary-bell}

The MathWorld material includes the complementary Bell numbers, which are the
set-partition transform at block weight $-1$.  They are especially useful as a
clean example of signed labelled enumeration.

\begin{definition}
Define
\begin{equation}
 \ComplementaryBellNumber{n}=\sum_{k=0}^{n}(-1)^k\StirlingSecondKind nk.               \label{eq:merged-complementary-bell}
\end{equation}
Thus
$\ComplementaryBellNumber{0}=1$ and the next values are
$-1,0,1,1,-2,-9,-9,50,267,413,\ldots$.
\end{definition}

\begin{theorem}[EGF, recurrence, and Dobiński-type series]
\label{thm:merged-complementary-bell}
The complementary Bell numbers satisfy
\begin{align}
 \sum_{n\geq0}\ComplementaryBellNumber{n}\frac{z^n}{n!}
   &=\exp(1-\EulerE^z),                                                    \label{eq:merged-complementary-egf}\\
 \ComplementaryBellNumber{n+1}
   &=-\sum_{k=0}^{n}\binom nk\ComplementaryBellNumber{k},                  \label{eq:merged-complementary-recurrence}\\
 \ComplementaryBellNumber{n}
   &=\EulerE\sum_{m=0}^{\infty}\frac{(-1)^m m^n}{m!}.                    \label{eq:merged-complementary-dobinski}
\end{align}
The last series is absolutely convergent.
\end{theorem}

\begin{proof}
Using the EGF of the Stirling numbers and summing over $k$ gives
\[
 \sum_{n\geq0}\ComplementaryBellNumber{n}\frac{z^n}{n!}
 =\sum_{k\geq0}\frac{(-1)^k}{k!}(\EulerE^z-1)^k=\EulerE^{1-\EulerE^z}.
\]
Differentiating yields $C'(z)=-\EulerE^zC(z)$; the exponential Cauchy product gives
the recurrence.  Finally expand $\EulerE^{-\EulerE^z}$:
\[
 \EulerE^{1-\EulerE^z}=\EulerE\sum_{m\geq0}\frac{(-1)^m}{m!}\EulerE^{mz}
 =\sum_{n\geq0}\left(\EulerE\sum_{m\geq0}\frac{(-1)^mm^n}{m!}\right)
                  \frac{z^n}{n!}.
\]
Absolute convergence follows from the ratio test for $m^n/m!$ at fixed $n$.
\end{proof}

\begin{corollary}[Convolution inverse of the Bell sequence]
\label{cor:merged-bell-convolution-inverse}
For every $n\geq0$,
\begin{equation}
 \sum_{k=0}^{n}\binom nk\BellNumber k\,\ComplementaryBellNumber{n-k}=\delta_{n0}.\label{eq:merged-bell-convolution-inverse}
\end{equation}
\end{corollary}

\begin{proof}
The two EGFs are $\EulerE^{\EulerE^z-1}$ and $\EulerE^{1-\EulerE^z}$, whose product is $1$.
\end{proof}

\begin{remark}[Signed combinatorial meaning]
Equation \eqref{eq:merged-complementary-bell} is the signed sum of all
partitions of $[n]$, with weight $(-1)^{\#\text{blocks}}$.  It is not a count
of an ordinary unweighted class; it is an Euler characteristic in the
incidence-algebra sense.  The convolution inverse is the exponential formula
applied to component weights $+1$ and $-1$.
\end{remark}

\chapter{Coordinate-free composition and derivatives of inverse maps}
\label{ch:merged-frechet-inverse}

\section{Fréchet Faà di Bruno formula}

\begin{theorem}[Coordinate-free Faà di Bruno]
\label{thm:merged-frechet-faa}
Let $E,F,G$ be Banach spaces, $g:E\to F$, $f:F\to G$, and let $D^r$ denote
the symmetric $r$-linear Fréchet derivative.  Then
\begin{equation}
 D^n(f\circ g)(x)[h_1,\ldots,h_n]
 =\sum_{\pi\in\Pi_n}D^{|\pi|}f(g(x))
 \Bigl[D^{|B|}g(x)[h_i:i\in B]\Bigr]_{B\in\pi}.                   \label{eq:merged-frechet-faa}
\end{equation}
\end{theorem}

\begin{proof}
Induct on $n$.  For $n=1$ this is the chain rule.  Differentiate a term
indexed by $\pi\in\Pi_n$ in direction $h_{n+1}$.  Either the derivative hits
the outer derivative, creating the new singleton block $\{n+1\}$, or it hits
one inner derivative associated with a block $B$, adjoining $n+1$ to $B$.
Every partition of $[n+1]$ is obtained uniquely in one of these two ways,
according as $\{n+1\}$ is a singleton or not.  The product and chain rules for
multilinear maps supply no further multiplicities because blocks are already
labelled subsets.
\end{proof}

\begin{corollary}
When all directions are equal to $h$,
\begin{equation}
 D^n(f\circ g)(x)[h^n]
 =\sum_{k=1}^{n}D^kf(g(x))
 \left[\ExponentialPartialBellPolynomial nk\bigl(Dg[h],D^2g[h^2],\ldots\bigr)\right],          \label{eq:merged-frechet-bell}
\end{equation}
where the Bell polynomial is interpreted as the symmetric tensor sum over
blocks.  In one dimension this is the usual scalar Faà di Bruno formula.
\end{corollary}

\begin{proof}
Group the set partitions in \eqref{eq:merged-frechet-faa} by their number of
blocks and then by block-size multiplicities.  The number of partitions of a
fixed profile is exactly the coefficient defining $\ExponentialPartialBellPolynomial nk$.
\end{proof}

\section{Recursive derivatives of an inverse function}

\begin{theorem}[Inverse derivative recursion]
\label{thm:merged-inverse-derivative}
Let $g=f^{-1}$ locally and $f'(g(y))\ne0$.  For $n\geq2$,
\begin{equation}
 g^{(n)}(y)=-\frac1{f'(g(y))}
 \sum_{k=2}^{n}f^{(k)}(g(y))
 \ExponentialPartialBellPolynomial nk\bigl(g'(y),g''(y),\ldots,g^{(n-k+1)}(y)\bigr).           \label{eq:merged-inverse-derivative}
\end{equation}
In particular, writing all derivatives of $f$ at $x=g(y)$,
\begin{align}
 g'&=\frac1{f'},\\
 g''&=-\frac{f''}{(f')^3},\\
 g'''&=\frac{3(f'')^2-f'f'''}{(f')^5},\\
 g^{(4)}&=\frac{-15(f'')^3+10f'f''f'''-(f')^2f^{(4)}}{(f')^7}.      \label{eq:merged-inverse-first-four}
\end{align}
\end{theorem}

\begin{proof}
Differentiate $f(g(y))=y$ $n$ times.  Faà di Bruno gives zero for $n\geq2$.
The term with one outer derivative is
$f'(g)\ExponentialPartialBellPolynomial n1=f'(g)g^{(n)}$; every other term involves only derivatives of
$g$ of order at most $n-1$.  Isolating the first term proves the recursion.
The displayed cases follow by substituting
$\ExponentialPartialBellPolynomial22=x_1^2$, $\ExponentialPartialBellPolynomial32=3x_1x_2$,
$\ExponentialPartialBellPolynomial42=4x_1x_3+3x_2^2$, and
$\ExponentialPartialBellPolynomial43=6x_1^2x_2$, then simplifying.
\end{proof}

\section{Several nested functions}

For a chain $f_1\circ\cdots\circ f_r$, repeated use of
Theorem~\ref{thm:merged-frechet-faa} produces a levelled hierarchy of set
partitions.  At the innermost level, derivative labels that strike the same
occurrence of $f_r$ form blocks; at the next level, those blocks are grouped
according to occurrences of $f_{r-1}$, and so on.  Flattening the hierarchy in
either association order is canonical, giving a combinatorial proof of
associativity of formal composition.  For three scalar functions,
\begin{align}
 D^n f(g(h(x)))
 &=\sum_{k=0}^{n}f^{(k)}(g(h(x)))
   \ExponentialPartialBellPolynomial nk\bigl((g\circ h)',\ldots,(g\circ h)^{(n-k+1)}\bigr),    \label{eq:merged-three-composition}\\
 (g\circ h)^{(j)}
 &=\sum_{\ell=0}^{j}g^{(\ell)}(h(x))
   \ExponentialPartialBellPolynomial j\ell\bigl(h',\ldots,h^{(j-\ell+1)}\bigr).              \label{eq:merged-three-composition-inner}
\end{align}
Substitution of the second formula into the first is a finite closed formula;
the partition hierarchy explains its multiplicities without additional index
machinery.

\chapter{Good's multivariate Lagrange inversion formula}
\label{ch:merged-good}

Let $\bm w=(w_1,\ldots,w_d)$ be determined formally by
\begin{equation}
 w_i=t_i\phi_i(\bm w),\qquad \phi_i(\bm0)\ne0.                     \label{eq:merged-good-system}
\end{equation}
For a multi-index $\bm n$, write
$\bm t^{\bm n}=\prod_it_i^{n_i}$.

\begin{lemma}[Multivariate residue substitution]
\label{lem:merged-multiresidue}
Let
\[
 T_i(\bm x)=x_iU_i(\bm x),\qquad U_i(\bm0)\ne0\quad(1\leq i\leq d).
\]
Then, for every formal Laurent series for which either side is defined
coefficientwise,
\begin{equation}
 \ResidueOperator_{\bm t}H(\bm t)\Differential t_1\cdots\Differential t_d
 =\ResidueOperator_{\bm x}H(\bm T(\bm x))
   \det\!\left(\frac{\partial T_i}{\partial x_j}\right)
   \Differential x_1\cdots\Differential x_d.                                          \label{eq:merged-multiresidue}
\end{equation}
\end{lemma}

\begin{proof}
Interpret the determinant times $\Differential x_1\cdots\Differential x_d$ as the pullback of
the top differential form $\Differential t_1\wedge\cdots\wedge\Differential t_d$.  By linearity
it is enough to take
$H(\bm t)=t_1^{m_1}\cdots t_d^{m_d}$.  If $m_j\ne-1$ for some $j$, then
\[
 \bm t^{\bm m}\,\Differential t_1\wedge\cdots\wedge\Differential t_d
 =\Differential\!\left(
   \frac{(-1)^{j-1}}{m_j+1}
   t_j^{m_j+1}\prod_{i\ne j}t_i^{m_i}
   \bigwedge_{\substack{1\le r\le d\\ r\ne j}}\Differential t_r\right).
\]
Pullback commutes with the exterior derivative.  The residue of an exact top
form is zero: after expansion it is a sum of formal partial derivatives, and
$\partial/\partial x_j$ cannot create an $x_j^{-1}$ term.  Thus both sides of
\eqref{eq:merged-multiresidue} vanish for this monomial.

It remains to take $m_1=\cdots=m_d=-1$.  Since every $U_i$ is a unit,
\[
 \frac{\Differential T_i}{T_i}=\frac{\Differential x_i}{x_i}+\frac{\Differential U_i}{U_i}.
\]
In
$\bigwedge_{i=1}^{d}\Differential T_i/T_i$, the product
$\bigwedge_i\Differential x_i/x_i$ has residue $1$.  Every other exterior-product term
contains at least one logarithmic differential $\Differential U_i/U_i$, whose
coefficients are ordinary power series, and consequently lacks at least one
of the negative powers $x_j^{-1}$ needed for the multivariate residue.  Its
residue is therefore zero.  The transformed residue is $1$, equal to the
residue of $\prod_i t_i^{-1}\Differential t_i$, and the proof is complete.
\end{proof}

\begin{theorem}[Lagrange--Good inversion]
\label{thm:merged-good}
For every formal series $F$,
\begin{align}
 [\bm t^{\bm n}]F(\bm w(\bm t))
 ={}&[\bm x^{\bm n}]F(\bm x)\prod_{i=1}^{d}\phi_i(\bm x)^{n_i}       \notag\\
 &\times\det\!\left(
 \delta_{ij}-\frac{x_j}{\phi_i(\bm x)}
                  \frac{\partial\phi_i}{\partial x_j}(\bm x)
 \right)_{i,j=1}^{d}.                                               \label{eq:merged-good}
\end{align}
\end{theorem}

\begin{proof}
Represent the left side as
\[
 \ResidueOperator_{\bm t}F(\bm w(\bm t))
 \prod_i t_i^{-n_i-1}\Differential t_1\cdots\Differential t_d.
\]
Use the inverse substitution $t_i=x_i/\phi_i(\bm x)$.  Its Jacobian entries are
\[
 \frac{\partial t_i}{\partial x_j}
 =\frac1{\phi_i}\left(\delta_{ij}-
   \frac{x_i}{\phi_i}\frac{\partial\phi_i}{\partial x_j}\right).
\]
After multiplying by $\prod_i t_i^{-n_i-1}$, the row factors
$1/\phi_i$ cancel one power from each $t_i^{-n_i-1}$, leaving
$\prod_i\phi_i^{n_i}$.  The remaining determinant initially has $x_i$ in row
$i$.  Write it as $\det(I-AB)$ with
$A=\DiagonalOperator(x_1,\ldots,x_d)$ and
$B=(\partial_j\log\phi_i)_{i,j}$; Sylvester's identity
$\det(I-AB)=\det(I-BA)$ moves the factor $x_i$ from row $i$ to
column $j$, producing the
displayed determinant.  The remaining residue is exactly the coefficient on
the right.
\end{proof}

\begin{example}[A coupled Fuss--Catalan system]
Let
\[
 u=x(1+u+v)^a,\qquad v=y(1+u+v)^b,
\]
and write $S=1+u+v$.  For integers $m,n\geq0$ and arbitrary $c$, the
constant coefficient is $[x^0y^0]S^c=1$, while for $m+n\geq1$,
\begin{equation}
 [x^my^n]S^c
 =\frac{c}{m!n!}(am+bn+c-1)^{\underline{m+n-1}}.                    \label{eq:merged-bivariate-fuss}
\end{equation}
When $am+bn+c\ne0$, this is equivalently
\begin{equation}
 [x^my^n]S^c
 =\frac{c}{am+bn+c}\,
   \frac{(am+bn+c)^{\underline{m+n}}}{m!n!}.                        \label{eq:merged-bivariate-fuss-rational}
\end{equation}
\end{example}

\begin{proof}
Apply Theorem~\ref{thm:merged-good} with
$\phi_1=S^a$, $\phi_2=S^b$, and $F=S^c$.  Put $A=am+bn+c$.  The determinant is
\[
 \det\begin{pmatrix}
 1-au/S&-av/S\\-bu/S&1-bv/S
 \end{pmatrix}=1-\frac{au+bv}{S}.
\]
Thus the desired coefficient is
\[
 [u^mv^n]\left(S^A-a uS^{A-1}-b vS^{A-1}\right).
\]
For $m+n\geq1$, multinomial extraction (with a term containing $u$ or
$v$ understood as zero when its requested exponent is zero) gives, after
putting all three terms over the common factor
$(A-1)^{\underline{m+n-1}}/(m!n!)$,
\[
 \frac{(A-1)^{\underline{m+n-1}}}{m!n!}
 \bigl(A-am-bn\bigr)
 =\frac{c}{m!n!}(A-1)^{\underline{m+n-1}}.
\]
This proves \eqref{eq:merged-bivariate-fuss}, including values where $A=0$ by
polynomial continuation.  Multiplying and dividing by $A$
when $A\ne0$ gives \eqref{eq:merged-bivariate-fuss-rational}.  The case
$m=n=0$ is immediate from $S(0,0)=1$.
\end{proof}

\begin{proposition}[Characteristic determinant]
\label{prop:merged-characteristic}
For the analytic system $H_i(\bm w,\bm t)=w_i-t_i\phi_i(\bm w)$, the Jacobian
with respect to $\bm w$, evaluated on a solution, has determinant equal to the
Good determinant.  Therefore an analytic branch can fail the implicit-function
test only where some $\phi_i$ is singular or this determinant vanishes.
\end{proposition}

\begin{proof}
Differentiate:
$\partial H_i/\partial w_j=\delta_{ij}-t_i\partial_j\phi_i$.
On a solution $t_i=w_i/\phi_i(\bm w)$, giving the row form of the determinant;
Sylvester's identity converts it to the column form in
\eqref{eq:merged-good}.
\end{proof}

\begin{remark}
The determinant identifies candidate branch points, not automatically the
radius of every derived series.  One must also establish accessibility from
the selected branch, exclude earlier coefficient singularities, and check
that the outer function $F$ does not cancel the local singular term.  This is
the multivariate analogue of examining critical values of a one-variable
inverse map; see also Bender and Richmond's asymptotic formulation
\cite{BenderRichmond1998}.
\end{remark}

\part{Bernoulli--Euler Calculus and Special Functions}

\chapter{Bernoulli, Euler, Genocchi, and Cauchy--Gregory families}
\label{ch:merged-bernoulli-euler}

Bernoulli numbers already appeared in the Eulerian and power-sum chapters.
The MathWorld material adds the surrounding polynomial calculus, Euler and
Genocchi families, Bernoulli polynomials of the second kind, and the
Euler--Maclaurin formula.  These families also illustrate the Sheffer
framework: ordinary Bernoulli and Euler polynomials are Appell sequences,
whereas the second-kind Bernoulli polynomials are Sheffer for the forward
difference operator.  Our convention is $\beta_1=-1/2$, in agreement with the
DLMF \cite{NISTDLMF}.

\section{Generating functions and Appell structure}

\begin{definition}
The Bernoulli numbers $\beta_n$, Bernoulli polynomials $\beta_n(x)$, and Euler
polynomials $\mathsf E_n(x)$ are defined by
\begin{align}
 \frac{t}{\EulerE^t-1}&=\sum_{n\geq0}\beta_n\frac{t^n}{n!},               \label{eq:merged-bernoulli-egf}\\
 \frac{te^{xt}}{\EulerE^t-1}&=\sum_{n\geq0}\beta_n(x)\frac{t^n}{n!},      \label{eq:merged-bernoulli-poly-egf}\\
 \frac{2\EulerE^{xt}}{\EulerE^t+1}&=\sum_{n\geq0}\mathsf E_n(x)\frac{t^n}{n!}.\label{eq:merged-euler-poly-egf}
\end{align}
\end{definition}

\begin{theorem}[Appell, translation, and explicit formulas]
\label{thm:merged-appell}
For $n\geq1$,
\begin{align}
 \beta_n'(x)&=n\beta_{n-1}(x),&
 \mathsf E_n'(x)&=n\mathsf E_{n-1}(x),                              \label{eq:merged-appell-derivative}\\
 \beta_n(x+y)&=\sum_{k=0}^{n}\binom nk\beta_k(x)y^{n-k},&
 \mathsf E_n(x+y)&=\sum_{k=0}^{n}\binom nk\mathsf E_k(x)y^{n-k},  \label{eq:merged-appell-translation}\\
 \beta_n(x)&=\sum_{k=0}^{n}\binom nk\beta_kx^{n-k}.                \label{eq:merged-bernoulli-explicit}
\end{align}
\end{theorem}

\begin{proof}
Differentiating either polynomial EGF with respect to $x$ multiplies it by
$t$, which proves the derivative identities.  Replacing $x$ by $x+y$
multiplies by $\EulerE^{yt}$, and the exponential Cauchy product gives translation.
Setting $x=0$ in the Bernoulli translation formula gives the explicit
polynomial expansion.
\end{proof}

\begin{theorem}[Difference and reflection identities]
\label{thm:merged-bernoulli-euler-basic}
For all relevant $n$,
\begin{align}
 \beta_n(x+1)-\beta_n(x)&=nx^{n-1} &&(n\geq1),                      \label{eq:merged-bernoulli-difference}\\
 \mathsf E_n(x+1)+\mathsf E_n(x)&=2x^n,                             \label{eq:merged-euler-difference}\\
 \beta_n(1-x)&=(-1)^n\beta_n(x),                                   \label{eq:merged-bernoulli-reflection}\\
 \mathsf E_n(1-x)&=(-1)^n\mathsf E_n(x).                           \label{eq:merged-euler-reflection}
\end{align}
Consequently $\beta_{2m+1}=0$ for $m\geq1$ and
$\mathsf E_{2m+1}(1/2)=0$.
\end{theorem}

\begin{proof}
Subtract the Bernoulli EGFs at $x+1$ and $x$:
\[
 \frac{te^{(x+1)t}-te^{xt}}{\EulerE^t-1}=te^{xt}.
\]
Add the two Euler EGFs to obtain $2\EulerE^{xt}$; coefficient comparison proves the
difference equations.  Replacing $t$ by $-t$ in the Bernoulli EGF gives
\[
 \frac{(-t)\EulerE^{-xt}}{\EulerE^{-t}-1}=\frac{te^{(1-x)t}}{\EulerE^t-1},
\]
and the analogous Euler calculation proves both reflections.  At $x=0$,
$\beta_n(1)=\beta_n$ for $n>1$ by the difference equation; reflection then
forces every odd $\beta_n$ beyond $\beta_1$ to vanish.  Euler reflection at
$x=1/2$ proves the last assertion.
\end{proof}

\section{Power sums and multiplication formulas}

\begin{theorem}[Faulhaber with an offset]
\label{thm:merged-faulhaber}
For $N,p\in\NaturalNumbers$,
\begin{equation}
 \sum_{j=0}^{N-1}(x+j)^p
 =\frac{\beta_{p+1}(x+N)-\beta_{p+1}(x)}{p+1}.                      \label{eq:merged-faulhaber}
\end{equation}
\end{theorem}

\begin{proof}
Apply \eqref{eq:merged-bernoulli-difference} to $\beta_{p+1}(x+j)$ and sum;
the left side telescopes.
\end{proof}

\begin{theorem}[Alternating power sums]
\label{thm:merged-alternating-sums}
For $N,p\in\NaturalNumbers$,
\begin{equation}
 \sum_{j=0}^{N-1}(-1)^j(x+j)^p
 =\frac{\mathsf E_p(x)+(-1)^{N-1}\mathsf E_p(x+N)}2.               \label{eq:merged-alternating-sums}
\end{equation}
The formula gives zero when $N=0$.
\end{theorem}

\begin{proof}
Multiply the Euler difference equation at $x+j$ by $(-1)^j$ and sum.  All
interior Euler terms cancel in pairs, leaving the two endpoints.
\end{proof}

\begin{theorem}[Raabe multiplication]
\label{thm:merged-raabe}
For every positive integer $q$,
\begin{equation}
 \sum_{r=0}^{q-1}\beta_n\!\left(x+\frac rq\right)
 =q^{1-n}\beta_n(qx).                                               \label{eq:merged-raabe}
\end{equation}
In particular,
\begin{equation}
 \beta_n\!\left(\frac12\right)=(2^{1-n}-1)\beta_n.                 \label{eq:merged-bernoulli-half}
\end{equation}
\end{theorem}

\begin{proof}
Sum the generating functions and use the finite geometric series:
\[
 \sum_{r=0}^{q-1}\frac{te^{(x+r/q)t}}{\EulerE^t-1}
 =\frac{te^{xt}}{\EulerE^{t/q}-1}
 =q\frac{(t/q)\EulerE^{(qx)(t/q)}}{\EulerE^{t/q}-1}.
\]
Coefficient comparison gives the multiplication formula.  Set $q=2,x=0$ and
solve for the half-value.
\end{proof}

\section{Genocchi numbers}

\begin{definition}
The Genocchi numbers $\GenocchiNumber{n}$ are defined by
\begin{equation}
 \frac{2t}{\EulerE^t+1}=\sum_{n\geq1}\GenocchiNumber{n}\frac{t^n}{n!}.                  \label{eq:merged-genocchi-egf}
\end{equation}
\end{definition}

\begin{theorem}[Bernoulli--Euler--Genocchi relations]
\label{thm:merged-genocchi}
For $n\geq1$,
\begin{equation}
 \GenocchiNumber{n}=2(1-2^n)\beta_n=n\mathsf E_{n-1}(0).                           \label{eq:merged-genocchi}
\end{equation}
Thus $\GenocchiNumber{1}=1$, $\GenocchiNumber{2m+1}=0$ for $m\geq1$, and
$\GenocchiNumber{2}=-1,\GenocchiNumber{4}=1,\GenocchiNumber{6}=-3,\GenocchiNumber{8}=17,\GenocchiNumber{10}=-155$.
\end{theorem}

\begin{proof}
Use
\[
 \frac{2t}{\EulerE^t+1}=\frac{2t}{\EulerE^t-1}-\frac{4t}{\EulerE^{2t}-1}
\]
and compare Bernoulli coefficients.  Multiplying the Euler EGF at $x=0$ by
$t$ gives the second equality.  Parity and initial values follow immediately.
\end{proof}

\section{Bernoulli polynomials of the second kind}

The historical terminology is not uniform.  We use the Roman--MathWorld
normalization, also called Cauchy polynomials of the first kind.

\begin{definition}
Define $b_n(x)$ by
\begin{equation}
 \frac{t(1+t)^x}{\log(1+t)}
 =\sum_{n\geq0}b_n(x)\frac{t^n}{n!}.                               \label{eq:merged-bernoulli-second-egf}
\end{equation}
\end{definition}

\begin{theorem}[Integral representation and finite difference]
\label{thm:merged-cauchy-polynomials}
For $n\geq0$,
\begin{align}
 b_n(x)&=\int_0^1\FallingFactorial{x+u}{n}\Differential u,                             \label{eq:merged-cauchy-integral}\\
 b_n(x+1)-b_n(x)&=n b_{n-1}(x)\qquad(n\geq1).                       \label{eq:merged-cauchy-difference}
\end{align}
\end{theorem}

\begin{proof}
Integrate the falling-factorial binomial expansion:
\[
 \int_0^1(1+t)^{x+u}\Differential u
 =(1+t)^x\frac{t}{\log(1+t)}.
\]
Coefficient comparison proves the integral.  Replacing $x$ by $x+1$ in the
EGF and subtracting multiplies it by $t$, proving the difference law.
\end{proof}

\begin{remark}
The EGF has Sheffer form with
$A(t)=t/\log(1+t)$ and $B(t)=\log(1+t)$.  Hence the lowering operator is the
forward difference $\EulerE^D-1$.  The first numbers $b_n=b_n(0)$ are
$1,1/2,-1/6,1/4,-19/30,9/4,\ldots$.
\end{remark}

\section{A modified Bernoulli normalization}

Define the even ordinary-series coefficients $\beta_{2m}^{\sharp}$ by
\begin{equation}
 \frac12\log\!\left(\frac{\EulerE^{x/2}-\EulerE^{-x/2}}{x/2}\right)
 =\sum_{m\geq0}\beta_{2m}^{\sharp}x^{2m},
 \qquad \beta_0^{\sharp}=\frac12\log2.                              \label{eq:merged-modified-bernoulli-def}
\end{equation}

\begin{proposition}
For even $n\geq2$,
\begin{equation}
 \beta_n^{\sharp}=\frac{\beta_n}{2n^2\Gamma(n)},                    \label{eq:merged-modified-bernoulli}
\end{equation}
and odd coefficients vanish.  In particular,
$\beta_2^{\sharp}=1/48$, $\beta_4^{\sharp}=-1/5760$, and
$\beta_6^{\sharp}=1/362880$.
\end{proposition}

\begin{proof}
Differentiate the defining series:
\[
 \frac14\coth\frac x2-\frac1{2x}
 =\frac1{2x}\left(\frac x2\coth\frac x2-1\right).
\]
The Bernoulli generating function implies
\[
 \frac x2\coth\frac x2=\sum_{m\geq0}\beta_{2m}\frac{x^{2m}}{(2m)!}.
\]
Integrating termwise gives
$\beta_{2m}^{\sharp}=\beta_{2m}/(4m(2m)!)$, which is the displayed formula.
\end{proof}

\section{Euler--Maclaurin summation}

Put $P_r(x)=\beta_r(\{x\})/r!$, with $\{x\}$ the fractional part.

\begin{theorem}[Euler--Maclaurin with remainder]
\label{thm:merged-euler-maclaurin}
Let $a<b$ be integers and $f\in C^{2m}([a,b])$.  Then
\begin{align}
 \sum_{k=a}^{b}f(k)
 ={}&\int_a^b f(x)\Differential x+\frac{f(a)+f(b)}2                           \notag\\
 &+\sum_{r=1}^{m}\frac{\beta_{2r}}{(2r)!}
 \bigl(f^{(2r-1)}(b)-f^{(2r-1)}(a)\bigr)+R_m,                      \label{eq:merged-euler-maclaurin}
\end{align}
where
\begin{equation}
 R_m=-\frac1{(2m)!}\int_a^b\beta_{2m}(\{x\})f^{(2m)}(x)\Differential x       \label{eq:merged-euler-maclaurin-rem}
\end{equation}
and
\begin{equation}
 |R_m|\leq\frac{2\RiemannZeta(2m)}{(2\pi)^{2m}}
                  \int_a^b|f^{(2m)}(x)|\Differential x.                     \label{eq:merged-euler-maclaurin-bound}
\end{equation}
\end{theorem}

\begin{proof}
On $[k,k+1]$, integration by parts gives
\[
 \int_k^{k+1}\left(x-k-\frac12\right)f'(x)\Differential x
 =\frac{f(k)+f(k+1)}2-\int_k^{k+1}f(x)\Differential x.
\]
Summing yields the trapezoidal term plus $\int_a^bP_1f'$.  Integrate by parts
repeatedly, using $P_r'=P_{r-1}$ away from integers and continuity of $P_r$ at
integers for $r\geq2$.  Boundary values are $\beta_r/r!$; odd Bernoulli
numbers beyond $\beta_1$ vanish.  The surviving boundary terms and final
integral are exactly \eqref{eq:merged-euler-maclaurin} and
\eqref{eq:merged-euler-maclaurin-rem}.

The periodic Bernoulli function has absolutely convergent Fourier series
\[
 \frac{\beta_{2m}(\{x\})}{(2m)!}
 =(-1)^{m+1}\frac{2}{(2\pi)^{2m}}
   \sum_{q\geq1}\frac{\cos(2\pi qx)}{q^{2m}}.
\]
Taking absolute values gives the supremum bound used in
\eqref{eq:merged-euler-maclaurin-bound}.
\end{proof}

\begin{corollary}[Stirling's expansion]
For fixed $M$,
\begin{equation}
 \log(n!)=\left(n+\frac12\right)\log n-n+\frac12\log(2\pi)
 +\sum_{r=1}^{M-1}\frac{\beta_{2r}}{2r(2r-1)n^{2r-1}}
 +O(n^{-(2M-1)}).                                                    \label{eq:merged-stirling-em}
\end{equation}
\end{corollary}

\begin{proof}
Apply Euler--Maclaurin to $f(x)=\log x$.  Its derivatives give the displayed
inverse powers.  The terms at the lower endpoint and the limiting remainder
combine into a constant, evaluated as $\frac12\log(2\pi)$ by Wallis' product.
The remainder bound one order farther gives the stated error.
\end{proof}

\chapter{Pochhammer symbols, polygamma functions, and Barnes' \texorpdfstring{$G$}{G}-function}
\label{ch:merged-gamma-barnes}

\section{Pochhammer algebra}

\begin{definition}
The rising Pochhammer symbol is
\begin{equation}
 (a)_n=\RisingFactorial an=a(a+1)\cdots(a+n-1),\qquad(a)_0=1.                \label{eq:merged-pochhammer}
\end{equation}
\end{definition}

\begin{theorem}[Gamma quotient and binomial series]
\label{thm:merged-pochhammer}
Where defined meromorphically,
\begin{equation}
 (a)_n=\frac{\Gamma(a+n)}{\Gamma(a)},                               \label{eq:merged-pochhammer-gamma}
\end{equation}
and, formally in $z$ and analytically for $|z|<1$,
\begin{equation}
 (1-z)^{-a}=\sum_{n\geq0}\frac{(a)_n}{n!}z^n.                      \label{eq:merged-binomial-series}
\end{equation}
\end{theorem}

\begin{proof}
Iterate $\Gamma(w+1)=w\Gamma(w)$ to obtain the quotient.  For the series, let
$F=\sum(a)_nz^n/n!$.  The recurrence $(a)_{n+1}=(a+n)(a)_n$ gives
$(1-z)F'=aF$ and $F(0)=1$, whose unique formal solution is $(1-z)^{-a}$.
The ratio test gives analytic convergence in the unit disk.
\end{proof}

\begin{corollary}[Chu--Vandermonde]
\begin{equation}
 \frac{(a+b)_n}{n!}=\sum_{k=0}^{n}\frac{(a)_k}{k!}
                         \frac{(b)_{n-k}}{(n-k)!}.                  \label{eq:merged-vandermonde-rising}
\end{equation}
\end{corollary}

\begin{proof}
Multiply the binomial series for exponents $a$ and $b$ and compare with the
series for exponent $a+b$.
\end{proof}

For $r\geq1$, define shifted generalized harmonic sums
\begin{equation}
 H_n^{(r)}(a)=\sum_{j=0}^{n-1}(a+j)^{-r},\qquad H_n(a)=H_n^{(1)}(a). \label{eq:merged-shifted-harmonic}
\end{equation}

\begin{theorem}[All parameter derivatives]
\label{thm:merged-pochhammer-derivatives}
For $m\geq1$ away from the apparent poles,
\begin{equation}
 \frac{\partial^m}{\partial a^m}(a)_n
 =(a)_n\ExponentialCompleteBellPolynomial m\!\left(H_n(a),-1!H_n^{(2)}(a),2!H_n^{(3)}(a),\ldots,
 (-1)^{m-1}(m-1)!H_n^{(m)}(a)\right).                              \label{eq:merged-pochhammer-derivatives}
\end{equation}
For $m=0$ the corresponding statement is the identity
$(a)_n=(a)_n\ExponentialCompleteBellPolynomial0$, with $\ExponentialCompleteBellPolynomial0=1$.  The right side extends
polynomially across all removable singularities.
\end{theorem}

\begin{proof}
Write $L(a)=\log(a)_n$.  Then
$L^{(r)}(a)=(-1)^{r-1}(r-1)!H_n^{(r)}(a)$.  Apply the exponential case of
Faà di Bruno to $\EulerE^{L(a)}$.  Since the left side is a polynomial in $a$,
agreement off the zeros supplies the polynomial continuation.
\end{proof}

\section{Digamma and polygamma series}

Define $\psi=\Gamma'/\Gamma$ and $\psi^{(m)}=D^m\psi$.

\begin{theorem}[Partial fractions]
\label{thm:merged-polygamma-series}
Away from nonpositive integers,
\begin{align}
 \psi(z)&=-\gamma+\sum_{k=0}^{\infty}
     \left(\frac1{k+1}-\frac1{k+z}\right),                          \label{eq:merged-digamma-series}\\
 \psi^{(m)}(z)&=(-1)^{m+1}m!\sum_{k=0}^{\infty}(k+z)^{-m-1}
 \quad(m\geq1).                                                     \label{eq:merged-polygamma-series}
\end{align}
\end{theorem}

\begin{proof}
Logarithmically differentiate Weierstrass' product
\[
 \frac1{\Gamma(z)}=ze^{\gamma z}
 \prod_{k\geq1}\left(1+\frac zk\right)\EulerE^{-z/k}.
\]
Normal convergence on compact sets away from the poles justifies the first
series and every further derivative.
\end{proof}

\begin{corollary}[Recurrences and special values]
\begin{align}
 \psi(z+1)&=\psi(z)+z^{-1},                                         \label{eq:merged-digamma-rec}\\
 \psi^{(m)}(z+1)&=\psi^{(m)}(z)+(-1)^mm!z^{-m-1},                   \label{eq:merged-polygamma-rec}\\
 \psi(n)&=H_{n-1}-\gamma,                                           \label{eq:merged-digamma-int}\\
 \psi^{(m)}(1)&=(-1)^{m+1}m!\RiemannZeta(m+1).                     \label{eq:merged-polygamma-one}
\end{align}
\end{corollary}

\begin{proof}
Subtract the partial-fraction series at consecutive arguments; the series
telescopes.  Iterate from $\psi(1)=-\gamma$, and set $z=1$ in the polygamma
series.
\end{proof}

\begin{theorem}[Reflection]
\label{thm:merged-polygamma-reflection}
For $z\notin\IntegerNumbers$,
\begin{align}
 \psi(1-z)-\psi(z)&=\pi\cot(\pi z),                                 \label{eq:merged-digamma-reflection}\\
 (-1)^m\psi^{(m)}(1-z)-\psi^{(m)}(z)
 &=\pi\frac{\Differential^m}{\Differential z^m}\cot(\pi z).                            \label{eq:merged-polygamma-reflection}
\end{align}
\end{theorem}

\begin{proof}
Multiplying the Weierstrass products for $1/\Gamma(1+z)$ and
$1/\Gamma(1-z)$ and using Euler's sine product gives
$\Gamma(z)\Gamma(1-z)=\pi/\sin(\pi z)$.  Logarithmic differentiation proves
the first formula; differentiating it $m$ more times proves the second.
\end{proof}

\begin{theorem}[Gauss multiplication]
\label{thm:merged-gauss-multiplication}
For $q\in\PositiveIntegers$,
\begin{equation}
 \Gamma(qz)=(2\pi)^{(1-q)/2}q^{qz-1/2}
 \prod_{r=0}^{q-1}\Gamma\!\left(z+\frac rq\right).                 \label{eq:merged-gauss-multiplication}
\end{equation}
Hence
\begin{align}
 \sum_{r=0}^{q-1}\psi\!\left(z+\frac rq\right)
 &=q\psi(qz)-q\log q,                                              \label{eq:merged-digamma-mult}\\
 \sum_{r=0}^{q-1}\psi^{(m)}\!\left(z+\frac rq\right)
 &=q^{m+1}\psi^{(m)}(qz).                                          \label{eq:merged-polygamma-mult}
\end{align}
\end{theorem}

\begin{proof}
Define
\[
 Q(z)=\frac{q^{qz-1/2}\prod_{r=0}^{q-1}\Gamma(z+r/q)}{\Gamma(qz)}.
\]
The gamma recurrence shows $Q(z+1/q)=Q(z)$.  Stirling's formula in a right
half-plane gives $Q(z)\to(2\pi)^{(q-1)/2}$ as $\Re z\to+\infty$ in a fixed
horizontal strip.  Periodicity therefore makes $Q$ equal to that constant;
meromorphic continuation extends the identity throughout its domain.  Solving
for $\Gamma(qz)$ proves \eqref{eq:merged-gauss-multiplication}.  Logarithmic
differentiation gives the digamma identity and further differentiation gives
the polygamma identity.
\end{proof}

\begin{corollary}
For every $m\geq0$,
\begin{equation}
 \frac{\partial^{m+1}}{\partial a^{m+1}}\log(a)_n
 =\psi^{(m)}(a+n)-\psi^{(m)}(a)
 =(-1)^m m!H_n^{(m+1)}(a).                                        \label{eq:merged-pochhammer-polygamma}
\end{equation}
\end{corollary}

\begin{proof}
Logarithmically differentiate the gamma quotient and compare with direct
differentiation of the finite product.
\end{proof}

\section{Barnes' function}

\begin{theorem}[Canonical product]
\label{thm:merged-barnes-product}
The normally convergent product
\begin{align}
 \BarnesGFunction(z+1)={}&(2\pi)^{z/2}
 \exp\!\left(-\frac{z(z+1)}2-\frac{\gamma z^2}{2}\right)             \notag\\
 &\times\prod_{k=1}^{\infty}\left(1+\frac zk\right)^k
 \exp\!\left(-z+\frac{z^2}{2k}\right)                             \label{eq:merged-barnes-product}
\end{align}
defines an entire function satisfying
\begin{equation}
 \BarnesGFunction(1)=1,\qquad \BarnesGFunction(z+1)=\Gamma(z)\BarnesGFunction(z).                                \label{eq:merged-barnes-recurrence}
\end{equation}
Consequently $\BarnesGFunction(n)=1!2!\cdots(n-2)!$ for $n\geq2$.
\end{theorem}

\begin{proof}
The logarithm of the $k$th factor is
$k\log(1+z/k)-z+z^2/(2k)=z^3/(3k^2)+O_K(k^{-3})$ on a compact set $K$,
so the product converges normally and defines an entire function with the
prescribed integral-multiplicity zeros.  At $z=0$ it equals $1$.  Divide the
finite products at $z$ and $z-1$, rearrange before taking the limit, and use
Weierstrass' gamma product; all quadratic convergence factors cancel and the
limit is $\Gamma(z)$.  Normal convergence justifies passage to the limit,
proving the recurrence.  Iteration gives the integer values.
\end{proof}

\begin{theorem}[Alexeiewsky identity]
\label{thm:merged-alexeiewsky}
With compatible logarithm branches,
\begin{equation}
 \log \BarnesGFunction(z+1)=\frac z2\log(2\pi)-\frac{z(z+1)}2
 +z\log\Gamma(z+1)-\int_0^z\log\Gamma(t+1)\Differential t.                   \label{eq:merged-alexeiewsky}
\end{equation}
\end{theorem}

\begin{proof}
Logarithmic differentiation of the product gives
\[
 \frac{\BarnesGFunction'(z+1)}{\BarnesGFunction(z+1)}
 =\frac12\log(2\pi)-z-\frac12+z\psi(z+1),
\]
where the simplification uses the digamma partial fractions.  The derivative
of the right side of \eqref{eq:merged-alexeiewsky} is identical, and both
sides vanish at $z=0$.
\end{proof}

\begin{corollary}
For $n\geq1$,
\begin{equation}
 \prod_{k=1}^{n}k^k=\frac{(n!)^n}{\BarnesGFunction(n+1)}.                          \label{eq:merged-hyperfactorial}
\end{equation}
\end{corollary}

\begin{proof}
In $(n!)^n/\BarnesGFunction(n+1)$, the exponent of $k$ is $n-(n-k)=k$.
\end{proof}

\begin{theorem}[Barnes asymptotics]
\label{thm:merged-barnes-asymptotic}
As $z\to\infty$ in $|\arg z|\leq\pi-\delta$,
\begin{align}
 \log \BarnesGFunction(z+1)\sim{}&\frac14z^2+z\log\Gamma(z+1)
 -\left(\frac12z(z+1)+\frac1{12}\right)\log z-\log A              \notag\\
 &+\sum_{k\geq1}\frac{\beta_{2k+2}}
 {2k(2k+1)(2k+2)z^{2k}},                                           \label{eq:merged-barnes-asymptotic}
\end{align}
where $A$ is the Glaisher--Kinkelin constant.  Equivalently,
\begin{equation}
 \log A=\lim_{n\to\infty}\left\{
 \sum_{k=1}^{n}k\log k-
 \left(\frac{n^2}{2}+\frac n2+\frac1{12}\right)\log n
 +\frac{n^2}{4}\right\}.                                         \label{eq:merged-glaisher}
\end{equation}
\end{theorem}

\begin{proof}
Insert Stirling's Euler--Maclaurin expansion for $\log\Gamma(t+1)$ into
\eqref{eq:merged-alexeiewsky} and integrate termwise.  The integrated
Bernoulli terms, after shifting the index, are exactly the displayed inverse
even powers; the integrated remainder has the next Poincaré order uniformly
in each closed sector.  The remaining integration constant is independent of
$z$.  Evaluate at positive integers using
$\log \BarnesGFunction(n+1)=\sum_{j=1}^{n-1}\log(j!)$ and summation by parts; comparison with
\eqref{eq:merged-glaisher} identifies it as $-\log A$.  This is the standard
DLMF normalization of Barnes' $G$-function \cite{NISTDLMF}.
\end{proof}

\begin{remark}
Polynomial identities such as the Pochhammer derivative formula are formal
and require no convergence.  The partial-fraction series and Barnes product
are convergent analytic identities, while
\eqref{eq:merged-barnes-asymptotic} is a Poincaré expansion.  Keeping these
three meanings of ``equals'' separate is essential in a reliable formula
compendium.
\end{remark}

\part{Catalan Refinements and Singularity Transfer}

\chapter{Narayana numbers, polygon dissections, associahedra, and the Tamari order}
\label{ch:merged-catalan-refinements}

The earlier Lagrange chapter derived the Catalan and Fuss--Catalan formulas.
The Wikipedia and MathWorld dossiers add refinements and geometric models.
We fix the convention that $\AssociahedronByDimension{N-3}$ is the associahedron attached to a
convex $N$-gon; hence $\dim\AssociahedronByDimension{N-3}=N-3$.  This eliminates the two-unit
index shifts common in sources that index instead by dimension or number of
binary-tree leaves.

\section{Dyck paths and reflection}

A Dyck path of semilength $n$ uses $n$ up-steps $U=(1,1)$ and $n$ down-steps
$D=(1,-1)$ and never falls below height $0$.

\begin{theorem}[Catalan formula]
\label{thm:merged-catalan-reflection}
The number of Dyck paths of semilength $n$ is
\begin{equation}
 \CatalanNumber{n}=\binom{2n}{n}-\binom{2n}{n-1}
 =\frac1{n+1}\binom{2n}{n}.                                       \label{eq:merged-catalan-reflection}
\end{equation}
\end{theorem}

\begin{proof}
Among all $\binom{2n}{n}$ balanced paths, reflect the prefix through the first
step that reaches height $-1$.  This is a bijection from bad paths to words
with $n+1$ up-steps and $n-1$ down-steps, counted by
$\binom{2n}{n-1}$.  Subtraction and simplification prove the formula.
\end{proof}

\begin{theorem}[First return]
\label{thm:merged-catalan-first-return}
The generating function $C(z)=\sum \CatalanNumber{n}z^n$ satisfies
\begin{equation}
 C(z)=1+zC(z)^2,
 \qquad C(z)=\frac{1-\sqrt{1-4z}}{2z}.                              \label{eq:merged-catalan-gf}
\end{equation}
\end{theorem}

\begin{proof}
Every nonempty path has a unique decomposition $UPDQ$, where $D$ is the first
return to height zero and $P,Q$ are Dyck paths.  This gives the functional
equation.  Of the two quadratic roots, only the displayed branch has constant
term $1$.
\end{proof}

\section{Narayana refinement}

Let $N(n,k)$ count Dyck paths with $k$ peaks, where a peak is an occurrence of
$UD$, and let $C(z,u)$ mark semilength by $z$ and peaks by $u$.

\begin{theorem}[Narayana formula]
\label{thm:merged-narayana}
For $1\leq k\leq n$,
\begin{equation}
 N(n,k)=\frac1n\binom nk\binom n{k-1},                              \label{eq:merged-narayana}
\end{equation}
and
\begin{equation}
 C(z,u)=1+zC(z,u)\bigl(C(z,u)+u-1\bigr).                            \label{eq:merged-narayana-gf}
\end{equation}
Consequently $\sum_kN(n,k)=\CatalanNumber{n}$ and $N(n,k)=N(n,n+1-k)$.
\end{theorem}

\begin{proof}
Put $H=C-1$.  In the first-return decomposition $UPDQ$, the initial arch is a
peak exactly when $P$ is empty.  Hence
$H=z(uC+(C-1)C)=z(1+H)(u+H)$.  Lagrange inversion gives
\[
 [z^nu^k]H=\frac1n[w^{n-1}u^k](1+w)^n(u+w)^n
 =\frac1n\binom n{k-1}\binom nk.
\]
Setting $u=1$ gives the Catalan equation, and the closed formula is symmetric
under $k\leftrightarrow n+1-k$.
\end{proof}

\begin{corollary}[Peak mean and variance]
For a uniformly random Dyck path of semilength $n\geq2$, the number $K_n$ of
peaks satisfies
\begin{equation}
 \Expectation K_n=\frac{n+1}{2},
 \qquad \Variance(K_n)=\frac{n^2-1}{4(2n-1)}.                            \label{eq:merged-narayana-moments}
\end{equation}
\end{corollary}

\begin{proof}
Symmetry gives the mean.  Differentiate the Narayana polynomial at $u=1$ and
use Vandermonde's convolution:
\begin{align*}
 \sum_k k\binom nk\binom n{k-1}&=n\binom{2n-1}{n},\\
 \sum_k k(k-1)\binom nk\binom n{k-1}
 &=n(n-1)\binom{2n-2}{n-1}.
\end{align*}
Divide by $nC_n$ and convert the second factorial moment to the variance.
\end{proof}

\section{Explicit Catalan bijections}

\begin{theorem}[Binary trees and parenthesizations]
\label{thm:merged-catalan-bijection}
Dyck paths of semilength $n$, rooted plane full binary trees with $n$ internal
vertices, and full parenthesizations of $n+1$ ordered factors are in mutually
explicit bijection.
\end{theorem}

\begin{proof}
Map the empty path to a leaf.  Recursively map $UPDQ$ to an internal root whose
left and right subtrees are the images of $P,Q$.  The construction is
reversible.  A binary tree becomes a parenthesization by labelling its leaves
from left to right and interpreting each internal node as the product of its
two children; the final multiplication recovers the root split.
\end{proof}

\begin{theorem}[Triangulations]
\label{thm:merged-triangulations}
Triangulations of a convex $(n+2)$-gon are counted by $\CatalanNumber{n}$.
\end{theorem}

\begin{proof}
Fix a root boundary edge.  The planar dual of a triangulation has one vertex
per triangle and one edge across each internal diagonal.  Attach a leaf across
every nonroot boundary edge and root at the triangle incident with the root
edge.  Each triangle has two remaining sides, so the result is a full plane
binary tree.  Conversely, glue a triangle at each internal vertex along its
parent side.  The constructions are inverse, and there are $n$ triangles.
\end{proof}

\begin{theorem}[Noncrossing partitions]
\label{thm:merged-noncrossing}
Noncrossing partitions of $[n]$ in circular order are counted by $\CatalanNumber{n}$.
\end{theorem}

\begin{proof}
Let $a_n$ count them.  Decompose a partition at the block containing $1$.
The gaps between successive elements of that block carry independent
noncrossing partitions.  If the block has $r+1$ elements, the gap sizes sum to
$n-r-1$.  Therefore the OGF $A(z)=\sum a_nz^n$ satisfies
\[
 A-1=\sum_{r\geq0}z^{r+1}A^{r+1}=\frac{zA}{1-zA},
\]
which rearranges to $A=1+zA^2$.  The decomposition is reversible, so this is
also a recursive bijection.
\end{proof}

\section{Dissections and the associahedron}

Let $D(N,d)$ count sets of $d$ pairwise noncrossing diagonals of a convex
$N$-gon.

\begin{theorem}[Kirkman--Cayley dissection formula]
\label{thm:merged-kirkman-cayley}
For $0\leq d\leq N-3$,
\begin{equation}
 D(N,d)=\frac1{d+1}\binom{N+d-1}{d}\binom{N-3}{d}.                 \label{eq:merged-kirkman-cayley}
\end{equation}
\end{theorem}

\begin{proof}
Root one boundary edge and take the planar dual of a dissection.  Its
$I=d+1$ regions become internal vertices and its $L=N-1$ other boundary edges
become leaves.  This is a rooted plane tree with no unary internal vertices,
and the construction is reversible.

In preorder, record outdegree minus one.  Internal entries are positive,
leaves contribute $-1$, and the total is $-1$.  Choose the $I$ internal
positions in $\binom{I+L}{I}$ ways and their positive values, summing to
$L-1$, in $\binom{L-2}{I-1}$ ways.  The cycle lemma says exactly one of the
$I+L$ cyclic shifts has nonnegative proper partial sums and hence encodes a
plane tree.  Thus the number is
\[
 \frac1{I+L}\binom{I+L}{I}\binom{L-2}{I-1}
 =\frac1I\binom{L+I-1}{I-1}\binom{L-2}{I-1},
\]
which becomes \eqref{eq:merged-kirkman-cayley} after substitution.
\end{proof}

\begin{corollary}[Associahedral face vector]
\label{cor:merged-associahedron-f}
A $j$-face of $\AssociahedronByDimension{N-3}$ corresponds to $N-3-j$ fixed noncrossing
diagonals, and
\begin{equation}
 f_j(\AssociahedronByDimension{N-3})=\frac1{N-2-j}
 \binom{2N-4-j}{N-3-j}\binom{N-3}{j}.                              \label{eq:merged-associahedron-f}
\end{equation}
Thus $f_0=\CatalanNumber{N-2}$, $f_{N-4}=N(N-3)/2$, and $f_{N-3}=1$.
\end{corollary}

\begin{proof}
A face is a partial triangulation, equivalently a set of fixed noncrossing
diagonals.  Its dimension is the number of diagonals still free, so substitute
$d=N-3-j$ in \eqref{eq:merged-kirkman-cayley} and simplify.
\end{proof}

\section{The Tamari order}

Orient the elementary rotation by
\begin{equation}
 ((AB)C)\longrightarrow(A(BC)).                                    \label{eq:merged-tamari-rotation}
\end{equation}

\begin{theorem}[The rotation relation is a partial order]
\label{thm:merged-tamari-order}
Reachability under \eqref{eq:merged-tamari-rotation} is antisymmetric and
therefore defines the Tamari partial order.
\end{theorem}

\begin{proof}
Label the leaves $1,\ldots,N$ from left to right and let $d_T(i)$ be the
depth of leaf $i$.  Define the integer potential
\[
 \Phi(T)=\sum_{i=1}^{N}2^i d_T(i).
\]
Under a rotation $((AB)C)\to(A(BC))$, the depths of leaves in $A$ decrease by
one, those in $C$ increase by one, and those in $B$ are unchanged.  Hence
\[
 \Phi(A(BC))-\Phi((AB)C)
 =\sum_{i\in C}2^i-\sum_{i\in A}2^i>0.
\]
Indeed every index in $C$ is larger than every index in $A$, and the smallest
power of two occurring in the first sum exceeds the sum of all preceding
powers.  Thus the potential strictly increases along each directed rotation,
so no directed cycle exists.  Reachability is reflexive and transitive by
definition, and acyclicity gives antisymmetry.
\end{proof}

\begin{theorem}[Raney numbers]
\label{thm:merged-raney}
If $T=1+zT^p$, then for integers $p,r\geq1$ and $n\geq0$,
\begin{equation}
 [z^n]T(z)^r=\frac{r}{pn+r}\binom{pn+r}{n}.                         \label{eq:merged-raney}
\end{equation}
\end{theorem}

\begin{proof}
The case $n=0$ is immediate.  For $n\geq1$, put
$Y=T-1=z(1+Y)^p$.  Lagrange--Bürmann applied to $(1+Y)^r$ gives
\[
 [z^n]T^r=\frac rn[w^{n-1}](1+w)^{pn+r-1}
 =\frac r{pn+r}\binom{pn+r}{n}.
\]
\end{proof}

\chapter{Darboux's finite formula and transfer from singularities}
\label{ch:merged-darboux}

The MathWorld entry emphasizes an exact finite integration identity, while
analytic combinatorics uses Darboux's principle to transfer local singular
expansions to coefficient asymptotics.  Both are proved here; the latter is
stated in a form sufficient for the Catalan and inverse-function examples
\cite{FlajoletSedgewick2009}.

\section{Finite Darboux identity}

\begin{theorem}[Darboux formula]
\label{thm:merged-darboux-finite}
Let $f$ have $n+1$ continuous derivatives on the segment from $a$ to $z$, and
let $\phi$ be a polynomial of degree at most $n$.  Then
\begin{align}
 \phi^{(n)}(0)(f(z)-f(a))
={}&\sum_{m=1}^{n}(-1)^{m-1}(z-a)^m
 \left(\phi^{(n-m)}(1)f^{(m)}(z)
      -\phi^{(n-m)}(0)f^{(m)}(a)\right)                             \notag\\
 &+(-1)^n(z-a)^{n+1}\int_0^1\phi(t)
 f^{(n+1)}(a+t(z-a))\Differential t.                                         \label{eq:merged-darboux-finite}
\end{align}
\end{theorem}

\begin{proof}
Put $d=z-a$, $F_m(t)=f^{(m)}(a+td)$, and
\[
 S(t)=\sum_{m=1}^{n}(-1)^{m-1}d^m\phi^{(n-m)}(t)F_m(t).
\]
Differentiate and reindex the terms in which the derivative hits $F_m$.
Every interior term cancels, leaving
\[
 S'(t)=d\phi^{(n)}(0)F_1(t)
       +(-1)^{n-1}d^{n+1}\phi(t)F_{n+1}(t).
\]
Integrate from $0$ to $1$, move the last integral to the other side, and
expand $S(1)-S(0)$.
\end{proof}

\begin{corollary}[Taylor integral remainder]
\begin{equation}
 f(z)=\sum_{m=0}^{n}\frac{f^{(m)}(a)}{m!}(z-a)^m
 +\frac{(z-a)^{n+1}}{n!}\int_0^1(1-t)^n
 f^{(n+1)}(a+t(z-a))\Differential t.                                        \label{eq:merged-taylor-remainder}
\end{equation}
\end{corollary}

\begin{proof}
Take $\phi(t)=(1-t)^n$ in the theorem.  All derivatives below order $n$
vanish at $t=1$, and
$\phi^{(n-m)}(0)=(-1)^{n-m}n!/m!$.
\end{proof}

\section{Algebraic singularities}

\begin{lemma}[Exact binomial coefficient]
\label{lem:merged-singular-binomial}
For $R\ne0$,
\begin{equation}
 [z^n](1-z/R)^{-\alpha}
 =R^{-n}\frac{\Gamma(n+\alpha)}{\Gamma(\alpha)\Gamma(n+1)},         \label{eq:merged-singular-binomial}
\end{equation}
where $1/\Gamma(\alpha)$ supplies continuation at nonpositive integers.  For
fixed admissible $\alpha$,
\begin{equation}
 [z^n](1-z/R)^{-\alpha}
 =\frac{R^{-n}n^{\alpha-1}}{\Gamma(\alpha)}
 \left(1+\frac{\alpha(\alpha-1)}{2n}+O(n^{-2})\right).              \label{eq:merged-singular-binomial-asymp}
\end{equation}
\end{lemma}

\begin{proof}
The exact formula is the generalized binomial series.  Subtract the Stirling
expansions of $\log\Gamma(n+\alpha)$ and $\log\Gamma(n+1)$ and exponentiate
to obtain the ratio expansion.
\end{proof}

\begin{theorem}[Darboux transfer]
\label{thm:merged-darboux-transfer}
Let $0<R<\rho$, let $h$ be analytic in $|z|<\rho$, and fix $J\geq1$.  Then
\begin{align}
 [z^n]h(z)(1-z/R)^{-\alpha}
={}&R^{-n}\sum_{j=0}^{J-1}\frac{(-R)^jh^{(j)}(R)}{j!}
 \frac{\Gamma(n+\alpha-j)}{\Gamma(\alpha-j)\Gamma(n+1)}            \notag\\
 &+O\!\left(R^{-n}n^{\Re\alpha-J-1}\right).                       \label{eq:merged-darboux-transfer}
\end{align}
In particular, if $h(R)\ne0$,
\begin{equation}
 [z^n]h(z)(1-z/R)^{-\alpha}
 \sim\frac{h(R)}{\Gamma(\alpha)}R^{-n}n^{\alpha-1}.                \label{eq:merged-basic-transfer}
\end{equation}
\end{theorem}

\begin{proof}
Taylor-expand
$h(z)=\sum_{j<J}h^{(j)}(R)(z-R)^j/j!+(z-R)^Jq(z)$.  Since
$(z-R)^j=(-R)^j(1-z/R)^j$, the finite terms give the displayed gamma ratios.
Choose $r$ with $R<r<\rho$; Cauchy's estimate gives $[z^k]q=O(r^{-k})$.
Convolving this geometric bound with
$[z^m](1-z/R)^{J-\alpha}=O(R^{-m}(m+1)^{\Re\alpha-J-1})$ gives the remainder,
because $R/r<1$.  The leading case follows from the gamma-ratio asymptotic.
\end{proof}

\begin{corollary}[Multiple dominant points]
If finitely many local singular models occur at points $\rho_j$ of common
modulus $R$ and the remainder is analytic past that circle, their coefficient
contributions add:
\begin{equation}
 [z^n]F(z)=\sum_j\frac{h_j(\rho_j)}{\Gamma(\alpha_j)}
 \rho_j^{-n}n^{\alpha_j-1}(1+O(n^{-1})).                            \label{eq:merged-multiple-transfer}
\end{equation}
Complex conjugate points therefore produce periodic or oscillatory terms.
\end{corollary}

\begin{proof}
Subtract the finite singular models and apply Cauchy's estimate to the analytic
remainder; then apply Theorem~\ref{thm:merged-darboux-transfer} term by term.
\end{proof}

\begin{proposition}[Logarithmic and meromorphic cases]
For $n\geq1$,
\begin{equation}
 [z^n]\log\frac1{1-z/R}=\frac{R^{-n}}n.                             \label{eq:merged-log-transfer}
\end{equation}
If a meromorphic function has principal parts
$c_{j,m}(1-z/\rho_j)^{-m}$ at its poles of smallest modulus, then its
coefficients equal
\begin{equation}
 \sum_{j,m}c_{j,m}\rho_j^{-n}\binom{n+m-1}{m-1}
\end{equation}
up to an exponentially smaller contribution from the analytic remainder.
\end{proposition}

\begin{proof}
The logarithmic formula is its Taylor series.  For poles, subtract all
principal parts, use the integer binomial coefficient formula, and apply
Cauchy's estimate to what remains.
\end{proof}

\section{Catalan asymptotics}

\begin{theorem}
As $n\to\infty$,
\begin{equation}
 \CatalanNumber{n}=\frac{4^n}{\sqrt\pi\,n^{3/2}}
 \left(1-\frac9{8n}+\frac{145}{128n^2}+O(n^{-3})\right).            \label{eq:merged-catalan-asymp}
\end{equation}
\end{theorem}

\begin{proof}
Use $\CatalanNumber{n}=(n+1)^{-1}\binom{2n}{n}$ and
\[
 \binom{2n}{n}=\frac{4^n}{\sqrt{\pi n}}
 \left(1-\frac1{8n}+\frac1{128n^2}+O(n^{-3})\right),
\]
which follows from Stirling's expansion.  Multiply by
$(n+1)^{-1}=n^{-1}(1-n^{-1}+n^{-2}+O(n^{-3}))$.
Alternatively, $C(z)=2-2\sqrt{1-4z}+O(1-4z)$ near $z=1/4$, and Darboux
transfer with $\alpha=-1/2$ gives the leading term directly.
\end{proof}

\begin{remark}[Choosing the asymptotic method]
Darboux transfer applies to fixed finite singularities.  The Bell EGF
$\exp(\EulerE^z-1)$ is entire, so its coefficients require the moving saddle point
and Lambert $W$ developed earlier.  Implicit algebraic functions are normally
handled by Lagrange inversion followed by singularity transfer; multivariate
systems require the characteristic determinant and often a saddle analysis.
\end{remark}

\part{Further Differential and Reversion Examples}

\chapter{Higher product rules, quotients, and elementary inverse series}
\label{ch:merged-further-calculus}

The five drafts contain several useful consequences of the two master
composition theorems that do not belong to a single classical triangle.  They
are collected here rather than repeated in the Bell-polynomial and inversion
chapters.

\section{Products, reciprocals, and quotients}

\begin{theorem}[Multinomial Leibniz rule]
\label{thm:merged-multinomial-leibniz}
For functions $f_1,\ldots,f_q$ possessing $n$ derivatives,
\begin{equation}
 D^n\!\left(\prod_{r=1}^{q}f_r(x)\right)
 =\sum_{j_1+\cdots+j_q=n}
   \binom{n}{j_1,\ldots,j_q}
   \prod_{r=1}^{q}f_r^{(j_r)}(x).                                  \label{eq:merged-multinomial-leibniz}
\end{equation}
\end{theorem}

\begin{proof}
Taylor's formula, interpreted either analytically near $x$ or formally in a
new variable $h$, gives
\[
 \prod_{r=1}^{q}f_r(x+h)
 =\prod_{r=1}^{q}\left(\sum_{j_r\geq0}
       f_r^{(j_r)}(x)\frac{h^{j_r}}{j_r!}\right).
\]
The coefficient of $h^n$ is the sum over $j_1+\cdots+j_q=n$ of
$\prod_r f_r^{(j_r)}(x)/j_r!$.  Multiplication by $n!$ proves
\eqref{eq:merged-multinomial-leibniz}.
\end{proof}

\begin{theorem}[Higher reciprocal and quotient rules]
\label{thm:merged-higher-quotient}
If $g(x)\ne0$, then
\begin{equation}
 \left(\frac1g\right)^{(n)}
 =\sum_{k=0}^{n}(-1)^k k!\,
   \frac{\ExponentialPartialBellPolynomial nk(g',g'',\ldots,g^{(n-k+1)})}{g^{k+1}}.           \label{eq:merged-reciprocal-derivative}
\end{equation}
Consequently,
\begin{equation}
 D^n\!\left(\frac fg\right)
 =\sum_{r=0}^{n}\binom nr f^{(n-r)}
   \sum_{k=0}^{r}(-1)^k k!\,
   \frac{\ExponentialPartialBellPolynomial rk(g',g'',\ldots,g^{(r-k+1)})}{g^{k+1}}.          \label{eq:merged-higher-quotient}
\end{equation}
For $n>0$ the $k=0$ term in
\eqref{eq:merged-reciprocal-derivative} vanishes.
\end{theorem}

\begin{proof}
Apply Fa\`a di Bruno's formula to $F(g(x))$ with $F(u)=u^{-1}$.  Since
$F^{(k)}(u)=(-1)^k k!u^{-k-1}$, the Bell-polynomial form gives
\eqref{eq:merged-reciprocal-derivative}.  Formula
\eqref{eq:merged-higher-quotient} follows by applying the ordinary two-factor
Leibniz rule to $f\,g^{-1}$ and substituting the reciprocal formula.
\end{proof}

\begin{corollary}[Trigonometric outer functions]
For $n\geq0$,
\begin{align}
 D^n\sin(g(x))
 &=\sum_{k=0}^{n}\sin\!\left(g(x)+\frac{k\pi}{2}\right)
   \ExponentialPartialBellPolynomial nk(g',g'',\ldots),                                      \label{eq:merged-sin-composite}\\
 D^n\cos(g(x))
 &=\sum_{k=0}^{n}\cos\!\left(g(x)+\frac{k\pi}{2}\right)
   \ExponentialPartialBellPolynomial nk(g',g'',\ldots).                                      \label{eq:merged-cos-composite}
\end{align}
\end{corollary}

\begin{proof}
Use $D^k\sin u=\sin(u+k\pi/2)$ and
$D^k\cos u=\cos(u+k\pi/2)$ in Fa\`a di Bruno's formula.
\end{proof}

\section{Diagonal identities furnished by inverse functions}

\begin{proposition}[Inverse sine and tangent diagonals]
For $m\geq0$,
\begin{align}
 \arcsin y
 &=\sum_{m\geq0}\frac1{2m+1}\binom{2m}{m}
      \frac{y^{2m+1}}{4^m},                                       \label{eq:merged-arcsin-series}\\
 [x^{2m}]\left(\frac{x}{\sin x}\right)^{2m+1}
 &=\frac1{4^m}\binom{2m}{m},                                      \label{eq:merged-sine-diagonal}\\
 [x^{2m}]\left(\frac{x}{\tan x}\right)^{2m+1}
 &=(-1)^m.                                                         \label{eq:merged-tangent-diagonal}
\end{align}
\end{proposition}

\begin{proof}
The binomial theorem gives
$(1-y^2)^{-1/2}=\sum_{m\geq0}\binom{2m}{m}y^{2m}/4^m$; integrating
from $0$ to $y$ proves \eqref{eq:merged-arcsin-series}.  Apply analytic
Lagrange inversion to the inverse pair $y=\sin x$, $x=\arcsin y$:
\[
 [y^{2m+1}]\arcsin y
 =\frac1{2m+1}[x^{2m}]
   \left(\frac{x}{\sin x}\right)^{2m+1}.
\]
Comparison with \eqref{eq:merged-arcsin-series} proves
\eqref{eq:merged-sine-diagonal}.  Similarly,
$\arctan y=\sum_{m\geq0}(-1)^my^{2m+1}/(2m+1)$ follows by integrating
$(1+y^2)^{-1}$.  Lagrange inversion for $y=\tan x$ then gives
\eqref{eq:merged-tangent-diagonal}.
\end{proof}

\begin{proposition}[Exponential--logarithmic diagonals]
For $n\geq1$,
\begin{align}
 [x^{n-1}]\left(\frac{x}{\EulerE^x-1}\right)^n&=(-1)^{n-1},              \label{eq:merged-norlund-diagonal}\\
 [x^{n-1}]\left(\frac{x}{\log(1+x)}\right)^n&=\frac1{(n-1)!}.     \label{eq:merged-cauchy-diagonal}
\end{align}
\end{proposition}

\begin{proof}
The inverse of $\EulerE^x-1$ is $\log(1+y)$, whose coefficient of $y^n$ is
$(-1)^{n-1}/n$.  Formula \eqref{eq:lagrange-basic} therefore gives
\eqref{eq:merged-norlund-diagonal}.  Conversely, the inverse of
$\log(1+x)$ is $\EulerE^y-1$, whose coefficient of $y^n$ is $1/n!$; the same
formula gives \eqref{eq:merged-cauchy-diagonal}.
\end{proof}

\section{Expansion about a nonzero equilibrium}

\begin{theorem}[Offset Lagrange--B\"urmann formula]
\label{thm:merged-offset-lagrange}
Let $u(t)$ be the unique formal solution of
\begin{equation}
 u=a+t\Phi(u),\qquad \Phi(a)\ne0.                                  \label{eq:merged-offset-equation}
\end{equation}
Then for every formal $H$ and $n\geq1$,
\begin{equation}
 [t^n]H(u(t))
 =\frac1n[z^{n-1}]H'(a+z)\Phi(a+z)^n.                              \label{eq:merged-offset-lagrange}
\end{equation}
The same formula holds analytically in a sufficiently small neighbourhood
when $H$ and $\Phi$ are analytic.
\end{theorem}

\begin{proof}
Put $w=u-a$ and $\phi(w)=\Phi(a+w)$.  Then $w=t\phi(w)$, so
Theorem~\ref{thm:lagrange-burmann}, applied to the outer series
$K(w)=H(a+w)$, gives
\[
 [t^n]K(w(t))=\frac1n[z^{n-1}]K'(z)\phi(z)^n,
\]
which is \eqref{eq:merged-offset-lagrange}.  The analytic statement follows
from the analytic inverse-function theorem exactly as in
Theorem~\ref{thm:analytic-lagrange}.
\end{proof}

\chapter{Kepler's equation as a complete reversion example}
\label{ch:merged-kepler}

A classical nonlinear equation with both a formal coefficient expansion and
a globally convergent Fourier representation is
\begin{equation}
 E=M+e\sin E,\qquad 0\leq e<1.                                    \label{eq:merged-kepler}
\end{equation}
Here $M$ is the mean anomaly, $E$ the eccentric anomaly, and $e$ is treated
mathematically as a parameter.

\begin{theorem}[Lagrange and Fourier--Bessel forms]
\label{thm:merged-kepler}
The unique real solution of \eqref{eq:merged-kepler} satisfies
\begin{align}
 E
 &=M+\sum_{k\geq1}\frac{\EulerE^k}{k!}
   \left(\frac{\Differential}{\Differential M}\right)^{k-1}\sin^k M,                  \label{eq:merged-kepler-reversion}\\
 &=M+2\sum_{n\geq1}\frac{J_n(ne)}{n}\sin(nM),                    \label{eq:merged-kepler-bessel}
\end{align}
where, for integral $n\geq0$,
\begin{equation}
 J_n(z)=\frac1{2\pi}\int_0^{2\pi}\cos(n\theta-z\sin\theta)\Differential\theta. \label{eq:merged-bessel-integral}
\end{equation}
The first formula is an identity of formal series in $e$; for every fixed
$0\leq e<1$, the second series converges absolutely and uniformly in real
$M$.
\end{theorem}

\begin{proof}
Formula \eqref{eq:merged-kepler-reversion} is the shifted reversion formula
\eqref{eq:shifted-lagrange-v} with $x=M$, $y=e$, and $f(x)=\sin x$.

For the Fourier form, the function $M(E)=E-e\sin E$ has derivative
$1-e\cos E\geq1-e>0$.  Hence it has a smooth inverse satisfying
$E(M+2\pi)=E(M)+2\pi$, and $E(M)-M$ is a smooth odd $2\pi$-periodic
function.  Write
\[
 E(M)-M=\sum_{n\geq1}c_n\sin(nM).
\]
Its sine coefficient, after changing variables from $M$ to $E$, is
\begin{align*}
 c_n
 &=\frac1\pi\int_0^{2\pi}(E-M)\sin(nM)\Differential M\\
 &=\frac e\pi\int_0^{2\pi}\sin E\,
   \sin\!\bigl(n(E-e\sin E)\bigr)(1-e\cos E)\Differential E.
\end{align*}
Since
\[
 \frac{\Differential}{\Differential E}\cos\!\bigl(n(E-e\sin E)\bigr)
 =-n(1-e\cos E)\sin\!\bigl(n(E-e\sin E)\bigr),
\]
integration by parts, with vanishing boundary term, gives
\[
 c_n=\frac e{\pi n}\int_0^{2\pi}\cos E\,
       \cos(nE-ne\sin E)\Differential E
 =\frac e n\bigl(J_{n-1}(ne)+J_{n+1}(ne)\bigr).
\]
It remains to prove the recurrence used to simplify this expression.  The
integral over a full period of
$\frac{\Differential}{\Differential\theta}\sin(n\theta-z\sin\theta)$ is zero; therefore
\[
 n\int_0^{2\pi}\cos(n\theta-z\sin\theta)\Differential\theta
 =z\int_0^{2\pi}\cos\theta\cos(n\theta-z\sin\theta)\Differential\theta.
\]
The product-to-sum identity and \eqref{eq:merged-bessel-integral} yield
$J_{n-1}(z)+J_{n+1}(z)=2nJ_n(z)/z$, with the value at $z=0$ understood by
continuity.  Thus $c_n=2J_n(ne)/n$, proving
\eqref{eq:merged-kepler-bessel}.  Finally, the inverse $E(M)$ is smooth for
$e<1$; repeated integration by parts makes its Fourier coefficients decay
faster than every fixed inverse power of $n$, which proves absolute and
uniform convergence.
\end{proof}

\begin{remark}
The two expansions organize the same solution in different coordinates.  The
Lagrange series is local in the perturbation parameter $e$ and exposes finite
derivative polynomials at every order.  The Fourier--Bessel series is global
in the angular variable $M$ and diagonalizes the periodic translation
symmetry.  Their equality is an instructive example of coefficient calculus
as a change of basis rather than a collection of unrelated formulas.
\end{remark}

\part{Algorithms, Normalization, and Reproducible Certificates}

\chapter{Exact coefficient algorithms and consistency checks}
\label{ch:merged-algorithms}

A formula compendium is reliable only when its normalizations can be checked
and its coefficients reproduced.  This chapter merges the algorithmic and
proof-audit appendices of the drafts.  Every recurrence is triangular and can
therefore be implemented over exact integers, rationals, finite fields, or
symbolic coefficient rings.

\section{Normalization dictionary}

\begin{auditbox}[title={Ordinary versus exponential coefficients}]
If
\[
 f(z)=\sum_{n\geq0}\alpha_nz^n
     =\sum_{n\geq0}f_n\frac{z^n}{n!},
\]
then $f_n=n!\alpha_n$.  With the conventions of this monograph,
\begin{equation}
 \OrdinaryPartialBellPolynomial nk(x_1,x_2,\ldots)
 =\frac{k!}{n!}\ExponentialPartialBellPolynomial nk(1!x_1,2!x_2,\ldots).                      \label{eq:merged-bell-normalization}
\end{equation}
Signed first-kind Stirling numbers are
$\SignedStirlingFirstKind{n}{k}=(-1)^{n-k}\UnsignedStirlingFirstKind nk$; inverse-matrix identities always use the
signed array.
\end{auditbox}

\begin{proof}[Proof of \eqref{eq:merged-bell-normalization}]
The exponential definition gives
\[
 \frac1{k!}\left(\sum_{j\geq1}x_jt^j\right)^k
 =\sum_{n\geq k}\ExponentialPartialBellPolynomial nk(1!x_1,2!x_2,\ldots)\frac{t^n}{n!}.
\]
The coefficient on the left is $\OrdinaryPartialBellPolynomial nk(x)/k!$ in ordinary normalization.
Comparison proves the conversion.
\end{proof}

\begin{center}
\renewcommand{\arraystretch}{1.18}
\begin{tabularx}{\textwidth}{@{}>{\bfseries}l X X@{}}
\toprule
Object & Defining expansion & Primary role\\
\midrule
$\SignedStirlingFirstKind{n}{k}$ & $\FallingFactorial{x}{n}=\sum_k\SignedStirlingFirstKind{n}{k}x^k$ & inverse factorial basis\\
$\UnsignedStirlingFirstKind nk$ & $\RisingFactorial{x}{n}=\sum_k\UnsignedStirlingFirstKind nkx^k$ & permutation cycles\\
$\StirlingSecondKind nk$ & $x^n=\sum_k\StirlingSecondKind nk\FallingFactorial{x}{k}$ & set partitions\\
$\ExponentialPartialBellPolynomial nk$ & exponential partial Bell polynomial & derivatives and labelled composition\\
$\OrdinaryPartialBellPolynomial nk$ & ordinary partial Bell polynomial & ordinary series composition\\
$\BellNumber n$ & $\sum_k\StirlingSecondKind nk$ & unweighted partitions\\
$\beta_n$ & $t/(\EulerE^t-1)=\sum\beta_nt^n/n!$ & power sums and Euler--Maclaurin\\
\bottomrule
\end{tabularx}
\renewcommand{\arraystretch}{1}
\end{center}

\section{Elementary series operations}

Let $A(z)=\sum a_nz^n$, $B(z)=\sum b_nz^n$, and $C(z)=\sum c_nz^n$.

\begin{algorithm}[Product, reciprocal, and quotient]
\label{alg:merged-product-quotient}
For $C=AB$,
\begin{equation}
 c_n=\sum_{j=0}^{n}a_jb_{n-j}.                                     \label{eq:merged-alg-product}
\end{equation}
If $a_0$ is invertible and $C=A^{-1}$,
\begin{equation}
 c_0=a_0^{-1},\qquad
 c_n=-a_0^{-1}\sum_{j=1}^{n}a_jc_{n-j}.                            \label{eq:merged-alg-reciprocal}
\end{equation}
For $C=B/A$,
\begin{equation}
 c_n=a_0^{-1}\left(b_n-\sum_{j=1}^{n}a_jc_{n-j}\right).            \label{eq:merged-alg-quotient}
\end{equation}
\end{algorithm}

\begin{proof}
The product formula is Cauchy's rule.  Equate coefficients in $AC=1$ or
$AC=B$ and isolate the term $a_0c_n$.
\end{proof}

\begin{algorithm}[Logarithm, exponential, and powers]
\label{alg:merged-exp-log-power}
Assume $a_0=1$.  If $L=\log A=\sum_{n\geq1}\ell_nz^n$, then
\begin{equation}
 \ell_n=a_n-\frac1n\sum_{j=1}^{n-1}j\ell_ja_{n-j}.                 \label{eq:merged-alg-log}
\end{equation}
If $A=\EulerE^L$, then
\begin{equation}
 a_n=\frac1n\sum_{j=1}^{n}j\ell_ja_{n-j}.                         \label{eq:merged-alg-exp}
\end{equation}
If $C=A^\alpha$, then
\begin{equation}
 c_0=1,\qquad
 c_n=\frac1n\sum_{j=1}^{n}\bigl((\alpha+1)j-n\bigr)a_jc_{n-j}.    \label{eq:merged-alg-power}
\end{equation}
\end{algorithm}

\begin{proof}
Compare coefficients in $A'=AL'$ to obtain the first two recurrences.  For
powers, compare degree $n-1$ in $AC'=\alpha A'C$ and isolate $nc_n$.
\end{proof}

\begin{algorithm}[Differentiation and integration]
If $C=A'$, then $c_n=(n+1)a_{n+1}$.  If
$C(z)=c_0+\int_0^zA(t)\Differential t$, then $c_{n+1}=a_n/(n+1)$.
\end{algorithm}

\section{Composition and reversion}

\begin{algorithm}[Composition power table]
\label{alg:merged-composition}
Assume $B(0)=0$ and set $p_{k,n}=[z^n]B(z)^k$.  Initialize
$p_{0,0}=1$, $p_{0,n}=0$ for $n>0$, and use
\begin{equation}
 p_{k,n}=\sum_{j=1}^{n}b_jp_{k-1,n-j},\qquad
 [z^n]A(B(z))=\sum_{k=0}^{n}a_kp_{k,n}.                            \label{eq:merged-alg-composition}
\end{equation}
\end{algorithm}

\begin{proof}
The first recurrence is $B^k=B\,B^{k-1}$; the second is
$A(B)=\sum_ka_kB^k$.  Since $B(0)=0$, only finitely many powers contribute to
each degree.
\end{proof}

\begin{algorithm}[Lagrange reversion]
\label{alg:merged-lagrange-reversion}
If $F(z)=a_1z+a_2z^2+\cdots$ with invertible $a_1$, then its inverse
$G(w)=\sum_{n\geq1}g_nw^n$ has
\begin{equation}
 g_n=\frac1n[z^{n-1}]\left(\frac z{F(z)}\right)^n.                 \label{eq:merged-alg-lagrange-reversion}
\end{equation}
\end{algorithm}

\begin{proof}
This is Lagrange--Bürmann with outer function $H(z)=z$.  The coefficient
$g_n$ depends only on $a_1,\ldots,a_n$, so every truncation is exact through
its declared order.
\end{proof}

\begin{algorithm}[Newton reversion]
\label{alg:merged-newton-reversion}
If $F(G_m(w))\equiv w\pmod{w^m}$, define
\begin{equation}
 G_{2m}=G_m-\frac{F(G_m)-w}{F'(G_m)}\pmod{w^{2m}}.                  \label{eq:merged-newton-reversion}
\end{equation}
Then $F(G_{2m})\equiv w\pmod{w^{2m}}$.
\end{algorithm}

\begin{proof}
Put $E=F(G_m)-w=O(w^m)$ and $\Delta=-E/F'(G_m)=O(w^m)$.  Taylor's formula
in the formal ring gives
$F(G_m+\Delta)=F(G_m)+F'(G_m)\Delta+O(\Delta^2)=w+O(w^{2m})$.
The denominator is invertible because its constant term is $a_1$.
\end{proof}

\begin{remark}[Arithmetic complexity]
The triangular product, reciprocal, exponential, and logarithm recurrences use
$O(N^2)$ scalar operations through degree $N$; a direct composition table uses
$O(N^3)$.  Faster multiplication and Newton iteration reduce asymptotic cost,
but the triangular forms are preferable for symbolic audits because every
coefficient carries a short dependency certificate.
\end{remark}

\section{Small exact cross-checks}

\[
\begin{array}{c|rrrrrrr}
 n\backslash k&0&1&2&3&4&5&6\\ \hline
0&1&0&0&0&0&0&0\\1&0&1&0&0&0&0&0\\2&0&1&1&0&0&0&0\\
3&0&2&3&1&0&0&0\\4&0&6&11&6&1&0&0\\5&0&24&50&35&10&1&0\\
6&0&120&274&225&85&15&1
\end{array}
\qquad
\begin{array}{c|rrrrrrr}
 n\backslash k&0&1&2&3&4&5&6\\ \hline
0&1&0&0&0&0&0&0\\1&0&1&0&0&0&0&0\\2&0&1&1&0&0&0&0\\
3&0&1&3&1&0&0&0\\4&0&1&7&6&1&0&0\\5&0&1&15&25&10&1&0\\
6&0&1&31&90&65&15&1
\end{array}                                                               \tag{\(*\)}
\]
The left table contains the unsigned first-kind Stirling numbers; its
$n$th row sums to $n!$.  The right table contains the second-kind Stirling
numbers, whose row sums through $n=6$ are
\[
 1,1,2,5,15,52,203.
\]
After the first-kind array is signed, its product with the second-kind array
must be the identity through this order.

\begin{proposition}[Recurrence uniqueness]
If two triangular arrays have the same boundary data and obey the same
recurrence from row $n-1$ to row $n$, then they agree identically.
\end{proposition}

\begin{proof}
Induct on the row number; every new entry is determined by already equal
entries in the preceding row.
\end{proof}

\section{Deterministic and modular certificates}

\begin{theorem}[Polynomial grid certificate]
\label{thm:merged-grid-certificate}
Let $P,Q\in K[x_1,\ldots,x_d]$ and suppose
$\deg_{x_i}(P-Q)\leq D_i$.  If $P=Q$ on a Cartesian grid
$S_1\times\cdots\times S_d$ with $D_i+1$ distinct values in $S_i$, then
$P=Q$ identically.
\end{theorem}

\begin{proof}
For one variable, a nonzero degree-$D$ polynomial has at most $D$ roots.
Induct on $d$, viewing $P-Q$ as a polynomial in $x_d$.  At each grid point in
the first $d-1$ coordinates all its coefficient polynomials vanish; the
induction hypothesis makes each coefficient identically zero.
\end{proof}

\begin{theorem}[Chinese-remainder certificate]
\label{thm:merged-crt-certificate}
If an integer $N$ has a proved bound $|N|<M/2$, where
$M=p_1\cdots p_s$ and the $p_i$ are pairwise coprime, and if
$N\equiv0\pmod{p_i}$ for every $i$, then $N=0$.
\end{theorem}

\begin{proof}
The congruences give $M\mid N$.  The only multiple of $M$ in
$(-M/2,M/2)$ is zero.
\end{proof}

\begin{remark}
Random evaluations are excellent error detectors but become proofs only after
a degree bound (Schwartz--Zippel or a full deterministic grid) is supplied.
Modular computations become exact certificates only after an integer-size
bound is recorded.  This distinction prevents numerical evidence from being
silently promoted to an identity.
\end{remark}

\section{Certificate for the Bell-number constant}

Put
\begin{equation}
 R_n=\frac{\BellNumber n^{1/n}\log(n+1)}{n}.                              \label{eq:merged-bell-ratio-cert}
\end{equation}
The following entries are upper endpoints rounded upward to six decimals,
computed from exact Bell integers and directed interval evaluation.

\begin{center}
\small
\begin{tabular}{rr@{\qquad}rr@{\qquad}rr}
\toprule
$n$&upper $R_n$&$n$&upper $R_n$&$n$&upper $R_n$\\
\midrule
1&0.693148&19&0.740692&37&0.707939\\
2&0.776837&20&0.738200&38&0.706638\\
3&0.790177&21&0.735818&39&0.705373\\
4&0.791840&22&0.733538&40&0.704143\\
5&0.789788&23&0.731353&41&0.702946\\
6&0.786239&24&0.729257&42&0.701780\\
7&0.782127&25&0.727243&43&0.700644\\
8&0.777875&26&0.725307&44&0.699537\\
9&0.773672&27&0.723442&45&0.698457\\
10&0.769603&28&0.721646&46&0.697403\\
11&0.765707&29&0.719913&47&0.696375\\
12&0.761995&30&0.718240&48&0.695370\\
13&0.758468&31&0.716623&49&0.694388\\
14&0.755118&32&0.715059&50&0.693429\\
15&0.751938&33&0.713545&51&0.692491\\
16&0.748915&34&0.712079&52&0.691573\\
17&0.746040&35&0.710657&53&0.690675\\
18&0.743302&36&0.709278&54&0.689795\\
\bottomrule
\end{tabular}
\end{center}

\begin{proposition}[Finite Bell certificate]
The maximum in the table occurs at $n=4$:
\begin{equation}
 R_4=\frac{15^{1/4}\log5}{4}=0.7918392970\ldots<0.792.              \label{eq:merged-bell-ratio-max}
\end{equation}
Together with the analytic estimate already proved for $n\geq39$, this gives
an overlapping certificate of
$\BellNumber n<(0.792n/\log(n+1))^n$ for every $n\geq1$.
\end{proposition}

\begin{proof}
Generate exact Bell integers by
$\BellNumber0=1$ and
$\BellNumber{n+1}=\sum_{k=0}^{n}\binom nk\BellNumber k$.  Directed interval arithmetic
for logarithm and exponentiation proves each upward-rounded entry.  All are
below $0.792$, and exact interval comparison locates the maximum at $n=4$.
The table covers $1\leq n\leq54$ and the analytic estimate covers
$n\geq39$, leaving no gap.
\end{proof}

\section{Verification protocol}

A robust workflow for a new identity is: normalize all signed/unsigned and
ordinary/exponential conventions; derive the identity by coefficient
extraction, a bijection, recurrence uniqueness, or incidence inversion;
expand through at least fourth order, where the Faà di Bruno multiplicities
$1,4,3,6,1$ expose most missing factorials; and retain exact or bounded modular
certificates for every computational step.  The derivation is the proof; the
other stages protect it against transcription and implementation errors.

\part{Source Concordance, Deduplication Record, and Verification}

\chapter{How the five drafts were unified}
\label{ch:merged-concordance}

The supplied \nolinkurl{articles/} directory contains five substantial drafts,
not five independent subjects.  Their main theorems overlap, but each draft also
contains material absent from the others.  The merger therefore used theorem-level
rather than paragraph-level deduplication.  For every repeated topic, one statement
and one complete proof were retained; genuinely different proofs were kept when
they reveal different structures, such as a bijection versus a generating-function
argument, or a formal-residue proof versus an analytic contour proof.

\section{Draft-by-draft contribution map}

\renewcommand{\arraystretch}{1.15}
\begin{longtable}{>{\raggedright\arraybackslash}p{0.27\textwidth}
                  >{\raggedright\arraybackslash}p{0.29\textwidth}
                  >{\raggedright\arraybackslash}p{0.34\textwidth}}
\toprule
\textbf{Draft} & \textbf{Distinct strengths} & \textbf{Use in the unified article}\\
\midrule
\endfirsthead
\toprule
\textbf{Draft} & \textbf{Distinct strengths} & \textbf{Use in the unified article}\\
\midrule
\endhead
\midrule
\multicolumn{3}{r}{\small Continued on the next page}\\
\endfoot
\bottomrule
\endlastfoot

\nolinkurl{articles/combinatorial_coefficient_calculus_bundle/}
\nolinkurl{combinatorial_coefficient_calculus.tex}
& The broadest proof-oriented core: restricted Stirling families, Bell-number
  interpretations and bounds, saddle methods, geometric Eulerian theory,
  normalized associahedral inversion, and Laplace asymptotics.
& Used as the structural spine.  Its collision-free notation and separation of
  formal from analytic claims were retained; duplicated introductions and ledgers
  were replaced by the present concordance.\\[1ex]

\nolinkurl{articles/Combinatorial_Formulae_and_Inversion_Theorems/}
\nolinkurl{Combinatorial_Formulae_and_Inversion_Theorems.tex}
& Especially detailed change-of-basis proofs, probability interpretations,
  zeta and gamma series, two-parameter reversion, and a compact collection of
  reverse recurrences and boundary conventions.
& Its nonduplicate derivations were reconciled with the Stirling, Bell, and
  shifted-reversion chapters.  Equivalent displays were collapsed to one theorem
  with signed and unsigned conventions stated explicitly.\\[1ex]

\nolinkurl{articles/Formulae_Compendium/Formulae_Compendium.tex}
& Incidence algebra, partition-lattice M\"obius inversion, Sheffer sequences,
  exponential Riordan arrays, coordinate-free Fa\`a di Bruno, recursive inverse
  derivatives, Good inversion, elementary inverse examples, and coefficient
  algorithms.
& Supplies the main new material in
  Chapters~\ref{ch:merged-umbral}, \ref{ch:merged-frechet-inverse},
  \ref{ch:merged-good}, \ref{ch:merged-further-calculus}, and
  \ref{ch:merged-algorithms}.  Its multivariate determinant was rederived from
  the Jacobian so that row and column conventions cannot be confused.\\[1ex]

\nolinkurl{articles/Formulae_Monograph_Package/Formulae_Monograph.tex}
& Strong normalization tables, low-order cross-checks, exact finite certificates,
  a Bell-number constant audit, and concise formula-family dependency maps.
& Supplies the reproducibility material in Chapter~\ref{ch:merged-algorithms}.
  Numerical assertions are presented with exact recurrences and interval or
  modular certificate procedures rather than as unsupported decimal evidence.\\[1ex]

\nolinkurl{articles/Formulae_Monograph_Source_and_PDF/Formulae_Monograph.tex}
& A polished synthesis of the composition and inversion chapters, including
  Fr\'echet derivatives, the delta-function interpretation of reversion, and
  Kepler's equation.
& Its distinct higher-dimensional and physical examples are retained in
  Chapters~\ref{ch:merged-frechet-inverse} and \ref{ch:merged-kepler}; repeated
  Stirling, Bell, Eulerian, and Lagrange material was absorbed into the spine.\\
\end{longtable}
\renewcommand{\arraystretch}{1}

\section{Deduplication rules}

The merger applied the following rules consistently.
\begin{enumerate}
\item \emph{Definitions are unique.}  Every array or polynomial family has one
      canonical defining equation.  Alternate normalizations appear only in an
      explicit conversion formula.
\item \emph{A theorem is stated once.}  Repeated formulas for Stirling transforms,
      Bell-polynomial composition, Fa\`a di Bruno, and Lagrange--B\"urmann were
      merged at the theorem level.  Independent proofs are retained only when they
      expose different mechanisms.
\item \emph{Boundary values are part of the theorem.}  Empty sums, $0^0$ in
      coefficient formulas, signed versus unsigned first-kind numbers, and the
      indexing of type-$B$ Eulerian numbers are fixed globally.
\item \emph{Formal and analytic assertions are separated.}  Algebraic identities
      in formal series require no convergence; radii, branches, termwise
      differentiation, and asymptotic remainders are asserted only after suitable
      analytic hypotheses are supplied.
\item \emph{Computations carry certificates.}  A finite table, primality report,
      modular period, or numerical constant is not called a symbolic theorem.
      The recurrence, exact integer data, interval enclosure, or modular witness
      needed to reproduce it is recorded.
\end{enumerate}

\chapter{Supplemental Wikipedia and MathWorld material}
\label{ch:merged-reference-concordance}

The encyclopedic files in \nolinkurl{Wikipedia/} and \nolinkurl{MathWorld/}
were used as topic checklists and formula-discovery aids, not as substitutes for
proofs.  Their formulas were normalized to the notation of this monograph and
proved from the common generating-function, residue, incidence-algebra, or
analytic framework.

\section{Wikipedia dossier}

\renewcommand{\arraystretch}{1.15}
\begin{longtable}{>{\raggedright\arraybackslash}p{0.26\textwidth}
                  >{\raggedright\arraybackslash}p{0.31\textwidth}
                  >{\raggedright\arraybackslash}p{0.33\textwidth}}
\toprule
\textbf{Files} & \textbf{Material checked} & \textbf{Integrated location}\\
\midrule
\endfirsthead
\toprule
\textbf{Files} & \textbf{Material checked} & \textbf{Integrated location}\\
\midrule
\endhead
\midrule
\multicolumn{3}{r}{\small Continued on the next page}\\
\endfoot
\bottomrule
\endlastfoot

\nolinkurl{Stirling_number.txt},
\nolinkurl{Stirling_numbers_of_the_first_kind.txt},
\nolinkurl{Stirling_numbers_of_the_second_kind.txt}
& Factorial-basis definitions, cycle and partition interpretations, recurrences,
  inclusion--exclusion, generating functions, congruences, asymptotics, negative
  indices, and restricted variants.
& Parts ``Stirling Numbers, Lah Numbers, and Basis Changes'' and ``Coefficient
  Calculus and Labelled Structures.''  Conflicting signs and edge conventions are
  resolved by the inverse-matrix definitions.\\[1ex]

\nolinkurl{Bell_number.txt}, \nolinkurl{Bell_polynomials.txt}
& Set partitions, Bell--Aitken triangle, Dobi\'nski and Touchard formulas,
  congruences, saddle growth, partial and complete Bell polynomials, composition,
  moments, cumulants, and determinants.
& Parts ``Bell Numbers and Set Partitions'' and ``Bell Polynomials and
  Composition,'' together with Chapter~\ref{ch:merged-complementary-bell} for the
  signed complementary transform.\\[1ex]

\nolinkurl{Eulerian_number.txt}
& Descent recurrences, Eulerian polynomials, Worpitzky interpolation, power sums,
  polylogarithms, real-rootedness, type $B$, second order, and geometric models.
& Part ``Eulerian Numbers and Descent Calculus.''  The zero-position descent in
  type $B$ is fixed before any recurrence is stated.\\[1ex]

Fa\`a di Bruno formula file (whose archive filename uses an escaped character)
& Set-partition, multiplicity-vector, Bell-polynomial, and multivariate forms of
  the higher chain rule.
& The original Fa\`a di Bruno chapter and
  Chapter~\ref{ch:merged-frechet-inverse}.  The filename is rendered in standard
  spelling while all coefficient normalizations are derived.\\[1ex]

\nolinkurl{Lagrange_inversion_theorem.txt},
\nolinkurl{Lagrange_reversion_theorem.txt}
& Formal residues, analytic inversion, shifted reversion, implicit equations,
  tree series, Lambert $W$, and coefficient forms.
& The Lagrange and shifted-reversion chapters,
  Chapter~\ref{ch:merged-good}, and the examples in
  Chapters~\ref{ch:merged-further-calculus}--\ref{ch:merged-kepler}.\\[1ex]

\nolinkurl{Catalan_number.txt}, \nolinkurl{Associahedron.txt}
& Catalan models, Narayana refinements, polygon triangulations, dissections,
  associahedral faces, and the Tamari order.
& The base inversion examples and Chapter~\ref{ch:merged-catalan-refinements},
  where the associahedron index is tied explicitly to its dimension and number of
  leaves.\\
\end{longtable}
\renewcommand{\arraystretch}{1}

\section{MathWorld dossier}

\renewcommand{\arraystretch}{1.15}
\begin{longtable}{>{\raggedright\arraybackslash}p{0.29\textwidth}
                  >{\raggedright\arraybackslash}p{0.30\textwidth}
                  >{\raggedright\arraybackslash}p{0.31\textwidth}}
\toprule
\textbf{Files or family} & \textbf{Added material} & \textbf{Integrated location}\\
\midrule
\endfirsthead
\toprule
\textbf{Files or family} & \textbf{Added material} & \textbf{Integrated location}\\
\midrule
\endhead
\midrule
\multicolumn{3}{r}{\small Continued on the next page}\\
\endfoot
\bottomrule
\endlastfoot

Bell, Bell polynomial, Dobi\'nski, and complementary-Bell files
& Signed block weights, the complementary Bell EGF, recurrence, Dobi\'nski-type
  sum, and convolution inverse of the ordinary Bell sequence.
& Chapter~\ref{ch:merged-complementary-bell}, with the ordinary Bell theory in the
  base chapters.\\[1ex]

Bernoulli number and polynomial, Bernoulli polynomial of the second kind,
modified Bernoulli, Euler polynomial, Genocchi, and Euler--Maclaurin files
& A unified Appell-family treatment; reflection, difference, multiplication and
  power-sum formulas; Genocchi conversion; Cauchy--Gregory coefficients; modified
  normalization; Euler--Maclaurin with a periodic remainder.
& Chapter~\ref{ch:merged-bernoulli-euler}.  The symbol $\beta_n$ prevents collision
  with Bell numbers.\\[1ex]

Pochhammer, digamma, polygamma, and Barnes $G$ files
& Rising-factorial parameter derivatives, logarithmic Bell-polynomial formulas,
  polygamma sums and multiplication, the Weierstrass product for $G$, Alexeiewsky's
  integral, and its asymptotic expansion.
& Chapter~\ref{ch:merged-gamma-barnes}; constants and branch conventions are
  matched to the NIST normalization \cite{NISTDLMF}.\\[1ex]

Catalan-number files
& Narayana numbers, explicit Catalan bijections, polygon dissections, associahedral
  face counts, Tamari order, and Raney numbers.
& Chapter~\ref{ch:merged-catalan-refinements}.\\[1ex]

Darboux-formula files
& The finite Darboux identity, its Taylor-remainder proof, algebraic and logarithmic
  singularity transfer, and the derivation of Catalan asymptotics.
& Chapter~\ref{ch:merged-darboux}, aligned with the transfer framework of
  analytic combinatorics \cite{FlajoletSedgewick2009}.\\[1ex]

First- and second-kind Stirling files and notebooks
& Low-order tables, explicit sums, recurrences, and normalization checks.
& The main Stirling parts and Chapter~\ref{ch:merged-algorithms}.  Notebook output
  was used only as an independent low-order check; the proofs do not depend on a
  computer algebra system.\\
\end{longtable}
\renewcommand{\arraystretch}{1}

The paired PDF and notebook files in the MathWorld directory serve different
purposes: the PDFs preserve the human-readable reference pages, while the notebooks
provide exact symbolic expressions suitable for spot checks.  Neither representation
is treated as a proof certificate by itself.

\chapter{Formula-status ledger and proof audit}
\label{ch:merged-proof-audit}

\section{Classes of source statements}
Every substantive source display was assigned to one of four classes before the
merger.
\begin{enumerate}
\item \emph{Exact identities.}  These are stated in normalized notation and proved
      by a bijection, coefficient extraction, recurrence uniqueness, finite
      differences, incidence inversion, formal residues, or analytic estimates.
\item \emph{Equivalent coordinate forms.}\par
      Examples include the partition, frequency-vector, and Bell polynomial forms
      of Fa\`a di Bruno; the matrix,
      basis-change, and generating-function forms of Stirling inversion; and the
      coefficient, derivative, and residue forms of Lagrange inversion.  They are
      presented as one theorem with explicit conversions.
\item \emph{Corrected formulas.}  An index shift, omitted range, absent Jacobian,
      ordinary/exponential factorial mismatch, or unstated branch hypothesis is
      repaired before proof.  The type-$B$ Eulerian convention, associahedron
      indexing, multivariate Good determinant, and delta-function Jacobian are the
      principal examples.
\item \emph{Schematic or computational claims.}  Undefined asymptotic placeholders
      are replaced by recursively defined coefficients.  Finite searches and
      numerical bounds are accompanied by exact recurrences, directed intervals,
      or modular certificates.
\end{enumerate}

\section{Reusable proof engines}
The long list of formulas reduces to a small number of mechanisms.
\begin{description}[leftmargin=34mm,style=nextline]
\item[Change of polynomial basis]
Falling and rising factorials diagonalize finite differences; Stirling and Lah
numbers are transition matrices.  Eulerian numbers provide the parallel transition
between powers and shifted binomial coefficients.

\item[Labelled decomposition]
Exponential generating functions turn disjoint labelled components into products
and sets of components into exponentials.  Bell polynomials retain the complete
block-size profile.

\item[Incidence inversion]
M\"obius inversion on the Boolean lattice gives binomial inversion; on the partition
lattice it gives moments versus cumulants.  Sheffer and Riordan arrays package the
same substitutions as lower-triangular operators.

\item[Formal residues]
The residue change-of-variables rule proves one-variable Lagrange--B\"urmann and,
with a Jacobian determinant, Good's multivariate formula.

\item[Singularity localization]
Cauchy's formula, saddle points, Darboux's identity, and transfer theorems convert
local analytic structure into coefficient asymptotics.  The assumptions needed for
uniformity and error terms are stated at the theorem where they are used.
\end{description}

\section{Independent verification protocol}
The exact algorithms in Chapter~\ref{ch:merged-algorithms} provide a second layer of
protection against transcription errors.  A new identity should be checked by:
\begin{enumerate}
\item converting every series to one declared ordinary or exponential
      normalization;
\item expanding far enough to encounter nontrivial multiplicities---fourth order
      already tests the Fa\`a di Bruno pattern $1,4,3,6,1$;
\item verifying triangular recurrences and inverse-matrix products over exact
      integers or rationals;
\item using a degree-bounded grid or a size-bounded Chinese-remainder certificate
      when the symbolic expression is too large for direct comparison; and
\item keeping analytic remainder estimates separate from finite numerical checks.
\end{enumerate}
This protocol does not replace derivation.  It makes the derivation reproducible and
catches the factorial, sign, and index errors most likely to survive visual review.

\section{External cross-checks}
The special-function normalizations and principal identities were cross-checked
against the current NIST Digital Library of Mathematical Functions
\cite{NISTDLMF}.  The multivariate inversion statement and its asymptotic role were
checked against Bender and Richmond \cite{BenderRichmond1998}; the one-variable
inversion framework was compared with Gessel's proof survey \cite{Gessel2016}.
These references are used for provenance and normalization.  The proofs in the body
are self-contained at the level of formal coefficient calculus developed here.

\section{Dependency map}
The mathematical dependence of the unified article is summarized by
\begin{center}
\renewcommand{\arraystretch}{1.15}
\begin{tabular}{c}
formal series, residues, and factorial bases\\
$\Downarrow$\\
Stirling and Lah transforms\\[1mm]
labelled products and partition lattices\\
$\Downarrow$\\
Bell polynomials, moments, cumulants, and composition\\[1mm]
Bell composition calculus and residue substitution\\
$\Downarrow$\\
Fa\`a di Bruno, inverse derivatives, and Lagrange--Good inversion\\[1mm]
analytic continuation and local singular expansions\\
$\Downarrow$\\
Darboux, saddle-point, and Laplace asymptotics
\end{tabular}
\renewcommand{\arraystretch}{1}
\end{center}
Eulerian and Catalan families form two transverse strands: the first translates
powers into descent and polyhedral data, while the second translates reversion into
trees, paths, polygon dissections, and associahedra.  The resulting structure is a
coefficient calculus rather than a catalogue of isolated identities.

\backmatter

\chapter*{Selected bibliography and provenance}
\addcontentsline{toc}{chapter}{Selected bibliography and provenance}

\begin{thebibliography}{99}

\bibitem{Bell1934}
E.~T. Bell,
``Exponential polynomials,''
\emph{Annals of Mathematics}, second series, vol.~35 (1934), pp.~258--277.

\bibitem{BenderRichmond1998}
E.~A. Bender and L.~B. Richmond,
``A multivariate Lagrange inversion formula for asymptotic calculations,''
\emph{Electronic Journal of Combinatorics}, vol.~5 (1998), Research Paper~33.

\bibitem{Broder1984}
A.~Z. Broder,
``The $r$-Stirling numbers,''
\emph{Discrete Mathematics}, vol.~49 (1984), pp.~241--259.

\bibitem{Charalambides2002}
C.~A. Charalambides,
\emph{Enumerative Combinatorics},
Chapman \& Hall/CRC, 2002.

\bibitem{Claesson2001}
A. Claesson,
``Generalized pattern avoidance,''
\emph{European Journal of Combinatorics}, vol.~22 (2001), pp.~961--971.

\bibitem{Comtet1974}
L. Comtet,
\emph{Advanced Combinatorics: The Art of Finite and Infinite Expansions},
translated by J.~W. Nienhuys, D. Reidel, 1974.

\bibitem{FlajoletSedgewick2009}
P. Flajolet and R. Sedgewick,
\emph{Analytic Combinatorics},
Cambridge University Press, 2009.

\bibitem{Gessel2016}
I.~M. Gessel,
``Lagrange inversion,''
\emph{Journal of Combinatorial Theory, Series A}, vol.~144 (2016),
pp.~212--249.

\bibitem{GKP1994}
R.~L. Graham, D.~E. Knuth, and O. Patashnik,
\emph{Concrete Mathematics}, second edition,
Addison--Wesley, 1994.

\bibitem{Good1960}
I.~J. Good,
``Generalizations to several variables of Lagrange's expansion, with applications
to stochastic processes,''
\emph{Proceedings of the Cambridge Philosophical Society}, vol.~56 (1960),
pp.~367--380.

\bibitem{Harper1967}
L.~H. Harper,
``Stirling behavior is asymptotically normal,''
\emph{Annals of Mathematical Statistics}, vol.~38 (1967), pp.~410--414.

\bibitem{Henrici1974}
P. Henrici,
\emph{Applied and Computational Complex Analysis}, vol.~1,
Wiley, 1974.

\bibitem{HsuShiue1998}
L.~C. Hsu and P.~J.-S. Shiue,
``A unified approach to generalized Stirling numbers,''
\emph{Advances in Applied Mathematics}, vol.~20 (1998), pp.~366--384.

\bibitem{Johnson2002}
W.~P. Johnson,
``The curious history of Fa\`a di Bruno's formula,''
\emph{American Mathematical Monthly}, vol.~109 (2002), pp.~217--234.

\bibitem{MoserWyman1955}
L. Moser and M. Wyman,
``An asymptotic formula for the Bell numbers,''
\emph{Transactions of the Royal Society of Canada}, section III, vol.~49
(1955), pp.~49--54.

\bibitem{NISTDLMF}
F.~W.~J. Olver, A.~B. Olde Daalhuis, D.~W. Lozier, B.~I. Schneider,
R.~F. Boisvert, C.~W. Clark, B.~R. Miller, B.~V. Saunders,
H.~S. Cohl, and M.~A. McClain, editors,
\emph{NIST Digital Library of Mathematical Functions},
National Institute of Standards and Technology.

\bibitem{Riordan1958}
J. Riordan,
\emph{An Introduction to Combinatorial Analysis},
Wiley, 1958.

\bibitem{Roman1984}
S. Roman,
\emph{The Umbral Calculus},
Academic Press, 1984.

\bibitem{Spivey2008}
M.~Z. Spivey,
``A generalized recurrence for Bell numbers,''
\emph{Journal of Integer Sequences}, vol.~11 (2008), Article 08.2.5.

\bibitem{Stanley2012}
R.~P. Stanley,
\emph{Enumerative Combinatorics}, vol.~1, second edition,
Cambridge University Press, 2012.

\bibitem{Temme1993}
N.~M. Temme,
``Asymptotic estimates of Stirling numbers,''
\emph{Studies in Applied Mathematics}, vol.~89 (1993), pp.~233--243.

\bibitem{Wilf2005}
H.~S. Wilf,
\emph{generatingfunctionology}, third edition,
A K Peters, 2005.

\bibitem{WikipediaDossier}
The supplied \nolinkurl{Wikipedia/} dossier: reference copies on Stirling,
Bell, Eulerian, Catalan, associahedral, Fa\`a di Bruno, and Lagrange inversion
and reversion topics.

\bibitem{MathWorldDossier}
The supplied \nolinkurl{MathWorld/} dossier: PDF and notebook reference copies
on Bell, Bernoulli--Euler--Genocchi, Catalan, Darboux, Pochhammer, polygamma,
Barnes $G$, Dobi\'nski, and Stirling topics.

\bibitem{SourceArchive}
The five LaTeX drafts in the supplied \nolinkurl{articles/} directory, merged and
normalized for this project.  The draft-by-draft contribution record appears in
Chapter~\ref{ch:merged-concordance}.

\end{thebibliography}

\printindex
\end{document}
